% PH_Foundations_v86.tex
% v86 changes from v85 (October 2026; full record in Appendix L.6 and
% p14d_psion_number_audit/ of the companion repository):
%   1. Sec.13 rewritten after the P14 audit: psion number and CP exact; antimatter
%      = antipsions; B and L required by observation with N_psi/3 = B + L; the
%      eta_B derivation (eqs. 14-15, 33-38) withdrawn; closed mechanisms listed;
%      level bifurcation at reheating as a candidate mechanism ([T3]); P14-B opened.
%   2. F12 suspended; Sec.14.3 direct route marked unsupported; App. A.4, J.1, J.4
%      annotated; Sec.8.2 status note and L.4.6 on the dual sector; open-problem
%      tables gain P14-B; abstract and summaries updated.
%   3. App. L.4.7: hadron-size geometry recorded as phenomenological; P14-C opened
%      (baryon binding scale and internal structure consistent with J.Sigma mixing).
%
% (v85 header follows)
% PH_Foundations_v85.tex
% v85 changes from v84 (October 2026; full record in Appendix L.5 and
% docs/CHANGES_v85.md of the companion repository):
%   1. Dark-matter mass: Landau-pole value (Sec.14.2) adopted as canonical;
%      Sec.8.1 states cascade level != body number; Table 3 relabelled by
%      cascade level; three k = 3 entries of Table 3 corrected.
%   2. Theorem KS: physical application and Corollary 8.2 suspended pending
%      a level rule for fragments (status note, Sec.8.2); verification text
%      replaced by what proofs/P_kinematic_stability.py actually checks.
%   3. F5 split into F5a (exclusion, [T2], E2) and F5b (rate, [T4], E1);
%      the underived 1e-50 cm^2 bound retired; ex-ante kill criterion (9.1).
%   4. F3 downgraded to [T4], conditional on P-PMNS; Sec.17.3 rewritten.
%   5. Table 4 caption: 18 rows (F5a/F5b).
%
% v84 changes from v83 (retained from the v84 source):
%   1. Sec.13.1(2)/13.2: added footnote distinguishing the chirality-based
%      C/CP mechanism used here (epsilon_psi, JL vs JR) from a separate,
%      independently-analysed charge-conjugation question on the six Q_k
%      values -- see Appendix L.4. The two are not the same object; this
%      revision does not claim the L.4 finding invalidates 13.1(2)/13.2.
%   2. Appendix L.4 (new): records (a) a proton/neutron tuple mislabelling
%      found in the companion repository's CHARGE_CONVENTIONS.md and its
%      correction; (b) a closed T1 proof that Theorem V.1's six Q_k values
%      do not satisfy state-wise charge conjugation (Q_k = -Q_{6-k}),
%      independent of particle labelling; (c) the adopted resolution
%      (aggregate-only conjugation, no published value changed); (d) an
%      open, unresolved tension between Sec.6's definite mu/tau
%      charged-lepton identification and Appendix H's "dual-sector"
%      (neutrino/baryon-like, generation ordering open) labelling of the
%      same index triples. See CHARGE_CONVENTIONS.md / SECTION_12_NOTE.md
%      in the companion repository for full derivations.
\documentclass[11pt,letterpaper]{article}
\usepackage[T1]{fontenc}
\usepackage[utf8]{inputenc}
\usepackage{mathptmx}
\usepackage{amsmath,amssymb}
\DeclareUnicodeCharacter{210F}{\ensuremath{\hslash}}
\usepackage[margin=1in]{geometry}
\usepackage{longtable,array}
\usepackage{pifont}
\usepackage{microtype}
\setlength{\parindent}{0pt}
\setlength{\tabcolsep}{3pt}
\setlength{\parskip}{5pt plus 1pt}
\setlength{\LTpre}{6pt}
\setlength{\LTpost}{6pt}
\usepackage[pdftex,hidelinks,pdftitle={The Psion Hypothesis},pdfauthor={Karlos Marden Maia}]{hyperref}
\pagestyle{myheadings}
\markboth{The Psion Hypothesis}{The Psion Hypothesis}
\begin{document}
\begingroup\centering \Huge \textbf{The Psion Hypothesis}\par\endgroup

\begingroup\centering \Large A Framework for the Origin of Matter, Antimatter, Dark Matter, and Gravity\par\endgroup

\begingroup\centering \large \textbf{October 2026}\par\endgroup

\begingroup\centering \normalsize PH\_Foundations v86\par\endgroup

\begingroup\centering \large Karlos Marden Maia\par\endgroup

\begingroup\centering Independent Researcher, Marietta, Georgia\par\endgroup

\begingroup\centering karlosmardenmaia@gmail.com\par\endgroup

{\small \noindent \textbf{[T1] }Exact algebraic theorem or exact Gel'fand--Yaglom functional-determinant calculation. No numerical verification required.\par}

{\small \noindent \textbf{[T2] }Structurally derived and numerically verified. Mechanism identified; derivation complete at leading order.\par}

{\small \noindent \textbf{[T2+] }Structurally closed: mechanism identified, scheme-equivalence complete. One explicit algebraic step remains for T1 closure.\par}

{\small \noindent \textbf{[T3] }Pipeline defined, calculation open. Strategy identified; numerical result not yet obtained.\par}

{\small \noindent \textbf{[T4] }Conjectural: mechanism identified but not derived. Dimensional estimates only.\par}

{\small \noindent \textbf{[SC] }Self-consistency input: the empirical anchor. Two exist: GN (dimensional; mass--gravity--cosmology) and $\Delta$m\textsuperscript{2}\textsubscript{21}/me\textsuperscript{2} (adimensional; neutrino-sector scale, §11).\par}

\begingroup\centering \small \textit{T1 and T2 results are the quantitative backbone of the PH.}\par\endgroup

\begingroup\centering \small \textit{All numerical verifications and derivation scripts: psion-proofs [2].}\par\endgroup

\begingroup\centering \footnotesize MSC: 81V05, 81V17, 83C65  |  Keywords: psion, Z6 symmetry, composite spectrum, cosmological constant, baryogenesis, dark matter, Newton's constant\par\endgroup

\clearpage
\section*{Abstract}

\noindent The Psion Hypothesis (PH) derives the Standard Model's structural constants and the mass--gravity--cosmology observables (CKM/PMNS mixing and individual quark masses are not derived) from a single fermionic field $\psi$ --- the psion --- with Z6 = Z3 $\times$ Z2 internal symmetry and spin j = 5/2. The framework has one dimensional empirical anchor in the mass--gravity--cosmology sector: Newton's constant GN (CODATA 2018: G = 6.67430 $\times$ 10\textsuperscript{-11} m\textsuperscript{3} kg\textsuperscript{-1} s\textsuperscript{-2}, i.e. GN = 6.70883 $\times$ 10\textsuperscript{-39} GeV\textsuperscript{-2} in natural units). This is the uniquely appropriate anchor: the causal hierarchy

\begingroup\centering L(Z6, g\textsuperscript{2}, j)  $\to$  condensate  ---Sakharov$\to$  GN [SC]  $\to$  $\Lambda$meta  $\to$  me\par\endgroup

\noindent places GN two causal levels above me; anchoring to me would be causally circular. From the single dimensional SC input GN, the condensation scale follows: $\Lambda$meta = $\surd$(K($\xi$)/GN) $\cdot$ exp($-$S/2g\textsuperscript{2}) = 1.2609 $\times$ 10\textsuperscript{12} GeV ([T2]), from which the electron mass is a genuine cascade output:

\begingroup\centering me = $\surd$(K($\xi$)/GN) $\cdot$ exp($-$S/2g\textsuperscript{2} $-$ Ne/be) = 0.5089 MeV  (0.42\% error).\par\endgroup

\noindent Here Ne = 10 and be = 325/1152 are exact theorems of Z6 [T1], and K($\xi$) = 1.16335 is the Gel'fand--Yaglom prefactor [T1]. This result is falsifiable: $\delta$me/me = $-$$\tfrac{1}{2}$ $\delta$GN/GN (prediction F2, testable via variation-of-constants cosmology).

\noindent A single condensate action S = 6Ne/be = 212.68 simultaneously yields: $\rho$$\Lambda$ = 2.492 $\times$ 10\textsuperscript{-47} GeV\textsuperscript{4} (0.93\%, [T2+]); GN recovered at 0.83\% self-consistency ([T2]); and the master invariant $\rho$$\Lambda$GN\textsuperscript{2} = 1.121 $\times$ 10\textsuperscript{-123} predicted vs 1.132 $\times$ 10\textsuperscript{-123} observed (0.93\%, [T2+] --- $\Lambda$meta cancels algebraically and the prediction is exactly GN-independent). The invariant admits the closed structural form [Nc/(27$\pi$\textsuperscript{4})]$\cdot$K($\xi$)\textsuperscript{2}$\cdot$e\textasciicircum{}($-$S(1+3/$\pi$\textsuperscript{2})) with S = Nc$\cdot$Ne/be, and is linear in ℏ --- the unique dimensionless completion is $\rho$$\Lambda$GN\textsuperscript{2}ℏ/c\textsuperscript{7} --- so that at fixed classical scales ℏ $\to$ 0 sends $\rho$$\Lambda$ $\to$ 0: the accelerated expansion is a first-order quantum effect, with the falsifiable co-drift law $\delta$$\rho$$\Lambda$/$\rho$$\Lambda$ = $-$2 $\delta$GN/GN (F17, §5.1).

\noindent Exact theorems establish: N$\nu$ = 3 active neutrino species; Ng = 8 gluons; Ngen = 3 charged-lepton generations (the fourth is kinematically forbidden within Z6); and the absolute stability of 39 dark-matter states at body numbers N $\in$ \{9, 12, 15, 18, 21, 27\}, with termination proved to N = 48. Further structural theorems include the unique-electron theorem (exactly one observable Q = $-$1 three-body state), the charge--chirality identity Q = nR $-$ 1, and Theorem DM-S (dark composites are fermions at odd body number; the dark sector is Pauli-structured --- forced-singlet s-wave states at N = 9, 12 and psion nuclei from N = 15 --- while the observable sector is exactly the Pauli-trivial sector). The emergent-gravity sector passes its consistency gauntlet: the condensation channel is a Lorentz scalar (local Lorentz invariance at mean-field grade; c\_gw = c\_$\gamma$ by structure), and the Z6 domain-wall problem resolves by infrared wall tension ($\sigma$\textsuperscript{1/3} = 0.42 MeV, 46 orders below the naive $\Lambda$meta\textsuperscript{3}) with complete chiral-bias annihilation in exact class ratios 5 : 8 : 9. Psion number is exactly conserved, so B + L is fixed and only B $-$ L can change, the reverse of the Standard Model; the earlier derivation of $\eta$B = 6.02 $\times$ 10\textsuperscript{-10} is withdrawn (§13.3), and level bifurcation during the formation of composites after reheating is identified as a candidate mechanism, conditional on a conserved baryon label (P14-B, §13.7). Neff = 3.054 is retained with a status note (§13.5). Two new forces are predicted: the proto-colour force (X17 singlet, 3.59\%, [T3]) and the Chiral-Flip force (CF-CC: structure [T1], rate bounded by |O|\textsuperscript{2} $\le$ 0.370). Seventeen falsifiable predictions span four timescales; hard T1 falsifiers include any detection of a fourth charged lepton, neutrino, or gluon. The framework has two empirical (SC) inputs, GN and $\Delta$m\textsuperscript{2}\textsubscript{21} (the neutrino-sector scale, §11); open problem P2/P12 ([T2]) seeks to close GN from 1 to 0 by solving the fixed-point condition $\Phi$($\Lambda$) = $\Lambda$ of the two-loop NJL kernel, while $\Delta$m\textsuperscript{2}\textsubscript{21} has a separate, unpromoted T4 closure candidate (P26-C, §11). This paper is the third of a four-act programme: from the Keplerian susceptibility G$'$r = Gp, through the phenomenological Relational Gravitational Dynamics (RGD) kernel, to the present Lagrangian, and finally to the derivation of the RGD kernel from this Lagrangian and its resolution of the dark-matter/MOND controversy [12, 13].

\clearpage
\section*{Contents}

{\small \noindent 1  The Primordial Field: The Psion\par}

{\small \noindent 2  The Lagrangian and Cascade Action\par}

{\small \noindent 3  The Cosmological Constant\par}

{\small \noindent 4  Three Predictions from One Action\par}

{\small \noindent 5  Master Invariant\par}

{\small \noindent 6  Composite Spectrum from Z6\par}

{\small \noindent 7  Physical Interpretation\par}

{\small \noindent 8  Kinematic Stability: Theorem KS\par}

{\small \noindent 9  Falsifiable Predictions --- Complete Registry\par}

{\small \noindent 10  Lepton Mass Spectrum\par}

{\small \noindent 11  Neutrino Masses\par}

{\small \noindent 12  Cosmic Accountability\par}

{\small \noindent 13  Baryogenesis: Status after the P14 Audit\par}

{\small \noindent 14  Dark Matter: Identity, Mass, and Relic Density\par}

{\small \noindent 15  Open Problems\par}

{\small \noindent 16  The Psion in Hawking Radiation\par}

{\small \noindent 17  Exotic Signatures and Experimental Context\par}

{\small \noindent 18  Discussion\par}

{\small \noindent 19  Conclusions\par}

{\small \noindent A  Complete Inventory of PH Results\par}

{\small \noindent B  Derivation of the Kink Fluctuation Operator\par}

{\small \noindent C  Exact Identities Used in Theorem 3.1\par}

{\small \noindent D  Theorem P16: Renormalisation of the Z6 Wall Determinant\par}

{\small \noindent E  Unit Conversion: $\rho$$\Lambda$ in GeV\textsuperscript{4}\par}

{\small \noindent F  Principal Conclusions and Spectrum Analysis\par}

{\small \noindent G  PH Proofs Repository\par}

{\small \noindent H  The Psion Diagram: Complete Calculus for Composite Classification\par}

{\small \noindent I  The Proto-Quark: PH Hadronic Constituents and Their Relation to SM Quarks\par}

{\small \noindent J  The Chirality Asymmetry $\varepsilon$$\psi$ = 25/169: Results\par}

{\small \noindent K  Gravitational Sector: Condensate Response\par}

{\small \noindent L  Method Audit Record: Eliminated Routes and Negative Findings\par}

{\small \noindent References\par}

\clearpage
\section*{1  The Primordial Field: The Psion}

\noindent One fermionic field $\psi$ --- a field whose quanta carry half-integer spin --- with Z6 = Z3 $\times$ Z2 internal symmetry generates all observed particles as composites. The cyclic group Z6 of order 6 decomposes as the product of Z3 (proto-colour sector) and Z2 (chirality). The quantum of this field is the psion.

\noindent Each psion state k $\in$ \{0, \ldots{}, 5\} carries three labels: (a) proto-colour --- a Z3 label \{red, green, blue\}, distinct from QCD colour; (b) chirality --- L (left-handed, k $\le$ 2) or R (right-handed, k $\ge$ 3); (c) magnetic quantum number mJ(k) = k $-$ 5/2 $\in$ \{$-$5/2, $-$3/2, $-$1/2, +1/2, +3/2, +5/2\}.

\noindent The electric charge is Qk = (4/9) Re($\omega$\textsuperscript{k mod 3}) $-$ (1/2)($-$1)\textasciicircum{}$\lfloor$k/3$\rfloor$ + 1/6,  $\omega$ = e\textsuperscript{2$\pi$i/3}.

\begin{footnotesize}\begin{longtable}{|>{\raggedright\arraybackslash}p{0.1246\textwidth}|>{\raggedright\arraybackslash}p{0.2354\textwidth}|>{\raggedright\arraybackslash}p{0.1938\textwidth}|>{\raggedright\arraybackslash}p{0.1800\textwidth}|>{\raggedright\arraybackslash}p{0.1662\textwidth}|}
\hline
\textbf{k} & \textbf{Proto-colour} & \textbf{Chirality} & \textbf{mJ(k)} & \textbf{Qk} \\ \hline
0 & red & L & $-$5/2 & +1/9 \\ \hline
1 & green & L & $-$3/2 & $-$5/9 \\ \hline
2 & blue & L & $-$1/2 & $-$5/9 \\ \hline
3 & red & R & +1/2 & +10/9 \\ \hline
4 & green & R & +3/2 & +4/9 \\ \hline
5 & blue & R & +5/2 & +4/9 \\ \hline
\end{longtable}\end{footnotesize}

\begingroup\centering\footnotesize\itshape Table 1: The six psion states.\par\endgroup

\noindent Here SU(2)meta is a global symmetry group of the psion Lagrangian, unrelated to the electroweak SU(2) of the Standard Model. In the charge formula, the three terms encode the Z3 proto-colour contribution, the Z2 chirality contribution, and the universal offset fixing $\Sigma$Qk = 1 (Theorem V.1, App. J). The six states span the spin-5/2 representation of SU(2)meta: J\textsuperscript{2} = j(j + 1) = (5/2)(7/2) = 35/4.

\section*{2  The Lagrangian and Cascade Action}

\subsection*{2.1  Lagrangian}

\begingroup\centering L = $\bar{\psi}$ i$\gamma$\textsuperscript{u}$\partial$\textsubscript{u}$\psi$ $-$ GUV($\bar{\psi}$$\psi$)\textsuperscript{2} $-$ ($\bar{\psi}$J\textsuperscript{a}$\psi$)($\bar{\psi}$$\Sigma$\textsuperscript{a}$\psi$)\hspace{2em}(1)\par\endgroup

\noindent The first term is the standard Dirac kinetic term; the second is a contact (four-fermion) interaction with UV coupling GUV controlling the self-interaction strength at high energies; Ja and $\Sigma$a are the SU(2)meta generators acting on the left and right sectors. The relative coefficient $-$1 is fixed by Z6 invariance --- no free parameter.

\subsection*{2.2  UV Fixed Point}

\noindent |$\beta$0| = NDirac + 2/Nf = 13/3; At the ultraviolet fixed point the beta function vanishes; the resulting UV fixed-point coupling g\textsuperscript{2} = 2$\pi$\textsuperscript{2}/3, $\alpha$meta = $\pi$/6.

\subsection*{2.3  Condensate Action}

\begingroup\centering S = 6Ne/be = 212.677,   be = b0/(ZNP Nc) = 325/1152\hspace{2em}(2)\par\endgroup

\noindent where b0 = 13/6, ZNP = 32/25, Nc = 6. Three projections of S:

\begingroup\centering e\textsuperscript{$-$S} $\to$ $\rho$$\Lambda$\hspace{2em}(3)\par\endgroup

\begingroup\centering e\textsuperscript{$-$S/g\textsuperscript{2}} $\to$ GN\hspace{2em}(4)\par\endgroup

\begingroup\centering e\textasciicircum{}($-$S(1+2/g\textsuperscript{2})) $\to$ $\rho$$\Lambda$GN\textsuperscript{2}\hspace{2em}(5)\par\endgroup

\noindent Their ratio e\textasciicircum{}(S(1$-$3/(2$\pi$\textsuperscript{2}))) $\approx$ 10\textsuperscript{78} resolves the hierarchy problem.

\section*{3  The Cosmological Constant}

\subsection*{3.1  Observed Value (Planck 2018)}

\begingroup\centering $\rho$$\Lambda$\_obs = $\Omega$$\Lambda$ $\cdot$ (3H0\textsuperscript{2}c\textsuperscript{2})/(8$\pi$GN) $\cdot$ ($\hslash$c)\textsuperscript{3} GeV = 2.515 $\times$ 10\textsuperscript{-47} GeV\textsuperscript{4}.\par\endgroup

\subsection*{3.2  Theorem: C = Nc/(27$\pi$\textsuperscript{4})}

\noindent \textbf{\textit{Theorem 3.1 ([T2+]). }}\textit{$\rho$$\Lambda$ = (Nc/27$\pi$\textsuperscript{4}) $\cdot$ $\Lambda$meta\textsuperscript{4} $\cdot$ e\textsuperscript{$-$S}.}

\noindent Proof sketch: (1) Z6 kink; (2) n = 2 Pöschl--Teller operator; (3) det$'$/det = 1/48 [T1]; (4) $\sigma$dw/N0 = 1 [T1]; (5) Z3 Faddeev--Popov factor = 27. Full derivation: Appendix C. $\square$

\subsection*{3.3  Numerical Verification}

\noindent $\rho$$\Lambda$\_pred = 2.492 $\times$ 10\textsuperscript{-47} GeV\textsuperscript{4} (error: 0.93\% against $\rho$$\Lambda$,obs = 2.515 $\times$ 10\textsuperscript{-47}).

\section*{4  Three Predictions from One Action}

\subsection*{4.1  From GN to $\Lambda$meta to Observables}

\noindent \textbf{Causal Hierarchy and the SC Input Rationale. }The PH has one dimensional empirical anchor: GN (Newton's constant, CODATA 2018). The neutrino sector carries one additional adimensional datum, $\Delta$m\textsuperscript{2}\textsubscript{21}/me\textsuperscript{2} (scale role; §11), not yet derived.

\begingroup\centering L(Z6, g\textsuperscript{2}, j) [pure theory] $\to$ condensate ---Sakharov$\to$ GN [SC] [anchor] $\to$ $\Lambda$meta = $\surd$(K($\xi$)/GN) e\textsuperscript{$-$S/2g\textsuperscript{2}} [derived scale] $\to$ me = 0.5089 MeV [genuine prediction]\par\endgroup

\noindent \textbf{Why GN, not me? }GN sits two causal levels above me in the hierarchy. Using me as anchor is causally circular: it defines $\Lambda$meta via a quantity the theory is supposed to derive from $\Lambda$meta. With GN as anchor, me becomes a first-principles output with a falsifiable coupling: $\delta$me/me = $-$$\tfrac{1}{2}$ $\delta$GN/GN.

\noindent From the single dimensional SC input GN = 6.70883 $\times$ 10\textsuperscript{-39} GeV\textsuperscript{-2}, the condensation scale is:

\begingroup\centering $\Lambda$meta = $\surd$(K($\xi$)/GN) $\cdot$ exp($-$S/2g\textsuperscript{2}) = 1.2609 $\times$ 10\textsuperscript{12} GeV.   [SC] $\to$ [T2]\hspace{2em}(6)\par\endgroup

\noindent \textbf{\textit{Remark 4.2 ($\Lambda$meta is not an input). }}$\Lambda$meta is a fixed function of the single dimensional anchor, $\Lambda$meta = f(GN), with f built entirely from derived constants (K($\xi$) [T1], S = 6Ne/be [T1/T2], g\textsuperscript{2} = 2$\pi$\textsuperscript{2}/3 [T1]). It carries no independent empirical content and no continuous freedom: a second dimensional degree of freedom would break the parameter-free closure of the adimensional invariants $\rho$$\Lambda$GN\textsuperscript{2} and GN me\textsuperscript{2} (eq. 9; Table 4, F1--F2), which close at the 0.9\% level with nothing to adjust. The Sakharov relation GN = K($\xi$) $\Lambda$meta\textsuperscript{-2} e\textsuperscript{$-$S/g\textsuperscript{2}} is this same definition inverted --- one equation, not two constraints. The residual theoretical softness is localized and declared: be enters the exponent at [T2] (open problem P1), and its status is tracked by the Class A/B partition of Table 7.

\noindent The logical output cascade:

\begingroup\centering GN [SC] (dimensional empirical anchor)\par\endgroup

\begingroup\centering ↓\par\endgroup

\begingroup\centering $\Lambda$meta = $\surd$(K($\xi$)/GN) e\textsuperscript{$-$S/2g\textsuperscript{2}}  (derived scale, [T2])\par\endgroup

\begingroup\centering ↓\par\endgroup

\begingroup\centering me (0.42\%) $\cdot$ $\rho$$\Lambda$ (0.93\%) $\cdot$ GN [self-consistency] (0.83\%) $\cdot$ mDM (T2) $\cdot$ Neff (<1$\sigma$) $\cdot$ lepton spectrum (5--9\%, P8b)\par\endgroup

\noindent \textbf{\textit{Open Problem 1 (P2/P12 [T2] --- eliminating the last input). }}\textit{Solving $\Phi$($\Lambda$) = $\Lambda$ of the two-loop NJL kernel without observational input would deliver L $\to$ $\Phi$($\Lambda$) = $\Lambda$ $\to$ $\Lambda$meta $\to$ GN $\to$ me, reducing SC inputs from 1 to 0 and making both GN and me first-principles predictions.}

\subsection*{4.2  Dark Energy}

\noindent $\rho$$\Lambda$ = C $\Lambda$meta\textsuperscript{4} e\textsuperscript{$-$S} = 2.492 $\times$ 10\textsuperscript{-47} GeV\textsuperscript{4} (0.93\%, [T2+]).

\subsection*{4.3  Newton's Constant (Self-Consistency Check)}

\begingroup\centering GN = K($\xi$) $\Lambda$meta\textsuperscript{-2} exp($-$S/g\textsuperscript{2}),   K($\xi$) = 1.16335.\hspace{2em}(7)\par\endgroup

\noindent Predicted: 6.653 $\times$ 10\textsuperscript{-39} GeV\textsuperscript{-2}; observed: 6.70883 $\times$ 10\textsuperscript{-39} GeV\textsuperscript{-2} (CODATA 2018, natural units); error: 0.83\%, [T2]. GN emerges via the Sakharov induced-gravity mechanism [3] --- a ground-state tensorial susceptibility of the condensate, not a mediator particle.

\noindent \textbf{\textit{Theorem 4.1 (Falsifiable coupling, [T2]). }}\textit{Differentiating the cascade law me = $\surd$(K($\xi$)/GN) exp($-$S/2g\textsuperscript{2} $-$ Ne/be):}

\begingroup\centering $\delta$me/me = $-$$\tfrac{1}{2}$ $\cdot$ $\delta$GN/GN.\hspace{2em}(8)\par\endgroup

\noindent Parameter-free; testable via variation-of-constants cosmology (prediction F2).

\section*{5  Master Invariant}

\begingroup\centering $\rho$$\Lambda$GN\textsuperscript{2} = (Nc/27$\pi$\textsuperscript{4}) K($\xi$)\textsuperscript{2} exp[ $-$6($\pi$\textsuperscript{2} + 3)Ne / ($\pi$\textsuperscript{2} be) ].\hspace{2em}(9)\par\endgroup

\noindent Predicted: 1.1214 $\times$ 10\textsuperscript{-123} (exactly GN-independent); observed: 2.515 $\times$ 10\textsuperscript{-47} $\times$ GN\textsuperscript{2} = 1.1320 $\times$ 10\textsuperscript{-123}; error: 0.93\% [T2+]. $\Lambda$meta cancels algebraically.

\noindent \textbf{\textit{Remark 5.1 (Irreducibility of exactly one dimensional input). }}The invariant web ($\rho$$\Lambda$GN\textsuperscript{2}, GN$\cdot$me\textsuperscript{2}, eq. 9) shows that $\Lambda$meta and GN cancel in every adimensional combination: all ratios close internally. Conversely, any dimensionful output requires exactly one dimensional input --- a scale-free kernel cannot fix a scale. One dimensional anchor is therefore a structural floor of any theory producing dimensionful observables, not a deficiency of this one.

\subsection*{5.1  Closed Form, ℏ-Linearity, and Falsification Handles}

\noindent Composing the registered chain (GN$\Lambda$meta\textsuperscript{2} = K($\xi$)e\textsuperscript{$-$S/g\textsuperscript{2}}; me = $\Lambda$meta e\textsuperscript{$-$Ne/be}) with the wall coefficient C = Nc/(27$\pi$\textsuperscript{4}) yields the master invariant in closed structural form: $\rho$$\Lambda$GN\textsuperscript{2} = C$\cdot$K($\xi$)\textsuperscript{2}$\cdot$e\textasciicircum{}($-$S(1 + 3/$\pi$\textsuperscript{2})), with S = 6Ne/be = Nc$\cdot$Ne/be = 69120/325 = 212.677 --- the condensate action is exactly Nc cascade depths (the registered S/g\textsuperscript{2} = 32.323 reproduces this to 10\textsuperscript{-6}). Every factor is structural --- \{Nc, $\pi$, K($\xi$), Ne, be\} --- with no dimensional input, which is why this combination is the invariant. Ledger of the 123 decades: ln C = $-$6.08; 2 ln K = +0.30; $-$2S/g\textsuperscript{2} = $-$64.65; $-$S = $-$212.68. Intermediate identity, confronted and resolved: $\rho$$\Lambda$ = C$\cdot$me\textsuperscript{6}/$\Lambda$meta\textsuperscript{2} is the exact rearrangement of Theorem 3.1 ($\rho$$\Lambda$ = C$\cdot$$\Lambda$meta\textsuperscript{4}$\cdot$e\textsuperscript{$-$S}) under the definitional S = 6Ne/be, since e\textsuperscript{$-$S} = (me/$\Lambda$meta)\textsuperscript{6} --- an identity, not a coincidence; the closed form above therefore inherits [T2+] from its components, and the $\le$ 0.1\% criterion reduces to an input-rounding check.

\noindent \textbf{ℏ-linearity. }Restoring SI units, the unique dimensionless combination of \{$\rho$$\Lambda$, GN, ℏ, c\} is $\rho$$\Lambda$GN\textsuperscript{2}ℏ/c\textsuperscript{7}: the exponents (2, 1, $-$7) are forced by the three dimension equations, so the invariant is linear in ℏ by necessity. Within the framework the linearity is structural: $\rho$$\Lambda$ is a one-loop functional determinant (Gel'fand--Yaglom) on the condensate background --- an O(ℏ) quantity. Stated without metaphor: at fixed GN, c, and the classical structure of the condensate, ℏ $\to$ 0 sends $\rho$$\Lambda$ $\to$ 0 --- the de Sitter component of the expansion is a first-order quantum effect. The standard cosmological-constant embarrassment presumes an additive sum of zero-point energies against a UV cutoff; here the invariant is a multiplicative semiclassical factor, e\textasciicircum{}($-$S(1+3/$\pi$\textsuperscript{2})) with S $\approx$ 213 --- instanton-grade suppression, not a 123-decimal cancellation. Loop-order attributions to the dimensionful factors separately are convention-dependent and carried as interpretation; the ℏ-linearity of the invariant is exact.

\noindent \textbf{Falsification handles (three, independent). }(i) The numerical identity: with paper-precision inputs the closed form must land within $\le$ 0.1\% (leading order exact); a persistent residual at O(1/S) = 0.47\% is registered as the two-loop term of the ℏ-linear structure --- with sign and size --- not as failure; a residual $\gg$ 1/S demotes the identities to coincidence. (ii) The co-drift law: because $\rho$$\Lambda$GN\textsuperscript{2} equals a pure structural number, any temporal or environmental drift of GN forces $\delta$$\rho$$\Lambda$/$\rho$$\Lambda$ = $-$2 $\delta$GN/GN --- jointly with F2 ($\delta$me/me = $-$$\tfrac{1}{2}$ $\delta$GN/GN): one drift parameter, three observables (GN, me, $\rho$$\Lambda$), two exact ratios ($-$$\tfrac{1}{2}$, $-$2); a single inconsistent pair falsifies the web (F17). (iii) Any claim that the expansion is tied to ℏ carries meaning only through these dimensionless statements; a drift of ℏ alone is a unit convention, not an observable.

\section*{6  Composite Spectrum from Z6}

\noindent Key definitions. A composite of N psions is a bound state of exactly N psion quanta; N is the body number. Constituents are indexed by k and written as a tuple, e.g. (0, 1, 2); chirality is abbreviated by an L/R prefix (L: k $\le$ 2, R: k $\ge$ 3). A composite is stable if it is the lightest state in its (I, Q, J) sector. Observable composites satisfy I $\equiv$ ($\Sigma$ki) mod 6 $\in$ \{0, 3\}.

\noindent \textbf{Two-body. }|Z3|\textsuperscript{2} $-$ 1 = 8 colour-charged bilinears $\to$ eight gluons (Ng = 8, [T1]). The Z6-singlet spin-1 massless bilinear $\to$ photon.

\noindent \textbf{Three-body (leptons and neutrinos).}

\begin{footnotesize}\begin{longtable}{|>{\raggedright\arraybackslash}p{0.2571\textwidth}|>{\raggedright\arraybackslash}p{0.0857\textwidth}|>{\raggedright\arraybackslash}p{0.0964\textwidth}|>{\raggedright\arraybackslash}p{0.1179\textwidth}|>{\raggedright\arraybackslash}p{0.3429\textwidth}|}
\hline
\textbf{Configuration} & \textbf{I} & \textbf{Q} & \textbf{$\Sigma$mJ} & \textbf{Identity} \\ \hline
LLL (0, 1, 2) & 3 & -1 & -9/2 & e\textsuperscript{-} (unique Q = $-$1, Theorem 6.2) \\ \hline
LLR (0, 1, 5) & 0 & 0 & -3/2 & n / $\nu$$\tau$ (dual-sector) \\ \hline
LLR (0, 2, 4) & 0 & 0 & -3/2 & $\nu$$\mu$ / $\Sigma$\textsuperscript{0}-like (dual-sector) \\ \hline
LRR (0, 4, 5) & 3 & +1 & 3/2 & Proton p (charge check: 1/9 + 4/9 + 4/9 = +1) \\ \hline
LLR (1, 2, 3) & 0 & 0 & -3/2 & $\nu$e / $\Lambda$\textsuperscript{0}-like (dual-sector; generation ordering open) \\ \hline
LRR (1, 3, 5) & 3 & +1 & 3/2 & $\Sigma$\textsuperscript{+}-like \\ \hline
LRR (2, 3, 4) & 3 & +1 & 3/2 & $\Sigma$\textsuperscript{+}-like (p-alternative configuration) \\ \hline
RRR (3, 4, 5) & 0 & +2 & 9/2 & Q = +2 exotic (BSM) \\ \hline
\end{longtable}\end{footnotesize}

\begingroup\centering\footnotesize\itshape Table 2: The eight observable three-body composites (generated census; \path{gen_composite_tables.py}).\par\endgroup

\noindent \textbf{\textit{Theorem 6.2 (Unique electron, [T1]). }}\textit{Among all C(6,3) = 20 three-body combinations, exactly one observable configuration carries Q = $-$1: the all-left state (0, 1, 2).}

\noindent Charged-lepton generations are therefore not distinct k-triples but the eigenvalue ladder of the leptonic transfer matrix on this single charge sector --- mass ratios are $\lambda$i/$\lambda$j (A.2.5) --- with |mJ\_R| $\in$ \{0, 3/2, 5/2\} as excitation labels (Table 5). Neutrinos are the three neutral configurations; $\nu$e $\leftrightarrow$ (1, 2, 3) (generation ordering open, $\nu$-mass sector).

\noindent \textbf{\textit{Theorem 6.3 (Charge--chirality identity, [T1]). }}\textit{For every observable three-body composite, Q = nR $-$ 1, where nR counts right-handed constituents.}

\noindent The per-psion form is Theorem V.1 (Qk+3 = Qk + 1). Corollary: the electron is the unique all-L state --- the chirality vacuum of the census; the same L-preference measured by $\varepsilon$$\psi$ (App.~J) marks the electron as the extremal composite.

\noindent \textbf{Standard Model comparison. }In the SM, electric charge is a U(1) input, the Gell-Mann--Nishijima relation expresses Q through two further inputs (isospin, hypercharge), and the existence of exactly one stable charged lepton is an empirical fact. In the PH all three are counting theorems: charge measures chirality content, and electron uniqueness follows from Pauli exclusion (which restricts the ground-state census to distinct-k triples) combined with Z6 charge arithmetic (which admits exactly one Q = $-$1 triple).

\noindent \textbf{Z\_N structure. }The deep form of Theorem 6.2 is universal: for every anchor group Z2M (M odd) the all-L extremal state is unique, with Qmin = $-$(M$-$1)/2. The literal form is a Z6 discriminant: the number of Q = $-$1 observables is C(M, (M$-$3)/2) = 1, 5, 21, 84 for M = 3, 5, 7, 9 (exact Z10 census in [2]). Two further facts select M = 3: the charge table is rational only for M = 3 among odd M $\ge$ 3 (the real cyclotomic field Q($\zeta$M)+ is rational iff M $\in$ \{1, 2, 3, 4, 6\}), and integer charge quantization of singlets already fails at M = 5 (fractional fifths in the exact Z10 census). Their elevation to a selection principle for Z6 is registered as conjecture C-Q-rat ([T3], [2]).

\subsection*{6.1  Pauli Structure: the SM as the Pauli-Trivial Sector}

\noindent The Dirac kinetic term of eq. (1) fixes the bookkeeping: the psion carries spacetime spin-1/2 and a six-valued internal SU(2)meta label (the internal spin-5/2 multiplet), giving exactly twelve single-particle states --- a counting resolved through the U-count adherence programme, whose registered decider (the first-principles derivation of the relic prefactor) fired on it, with the review of its two residuals recorded in [2]. Two Pauli thresholds follow. From N $\ge$ 7, label repetition is forced and the auto-singlet lemma ([T1]) applies: two fermions on the same label in the same orbital are a forced spin singlet. From N $\ge$ 13, full state repetition is forced and spatial antisymmetrization begins: psion nuclei proper start at N = 15, while N = 9 and 12 are s-wave states. For N = 9 the ground state is decided by a uniqueness theorem ([T1]): among the 32 neutral singlet candidates, exactly one is a pure s-wave configuration (nk $\le$ 2) --- (2, 2, 2, 1, 1, 1), three forced singlets plus a free-spin triplet on the distinct labels \{3, 4, 5\}, with mean-field J$\Sigma$ energy maximally anti-aligned ($\langle$Jz$\rangle$$\langle$$\Sigma$z$\rangle$ = $-$81/2). The free-spin triplet is the same configuration class as the three-body leptons (distinct labels, s-wave, three free spins); the observed lepton spin J = 1/2 anchors the spin selector, giving Jgs = 1/2 and g = 2 --- the derived pair (2, 2$\pi$/3) of the relic prefactor (§14.1). Conversely, every observable SM composite is Pauli-trivial: bilinears are psion--antipsion pairs, and all N = 3 identifications are distinct-k. The SM shadow of §F.3 acquires a structural reading: the Standard Model is exactly the sector of the spectrum that requires no Pauli structure. The historical resonance is inverted with respect to the quark model: there, the Pauli crisis of the $\Delta$++ forced the invention of colour; here, the observable sector never meets the crisis, and the Pauli pressure discharges entirely into the dark sector --- as forced singlets at N = 9, 12 and nuclear structure from N = 15. Stability is unaffected: shell corrections below 13\% of the N = 9 cascade exponent cannot alter any Theorem KS fragmentation bound.

\noindent \textbf{\textit{Theorem 6.1 (N$\nu$ = 3, [T1]). }}\textit{All C(6,3) = 20 three-body Z6-colour-singlet combinations yield exactly three neutral stable states --- the three active neutrino generations.}

\noindent \textbf{Dark matter. }39 stable Z6-singlet states at N $\in$ \{9, 12, 15, 18, 21, 27\}. mDM = 2$\Lambda$meta/NLandau = 7.205 $\times$ 10\textsuperscript{10} GeV, NLandau = 35.

\section*{7  Physical Interpretation}

\noindent S $\approx$ 212.7 resolves the hierarchy problem: e\textasciicircum{}(S(1$-$3/(2$\pi$\textsuperscript{2}))) = e\textasciicircum{}180.4 $\approx$ 10\textsuperscript{78}.

\subsection*{7.1  The Spin-2 State hAB and the Graviton (PC1, [T2])}

\noindent GN arises via Sakharov induced-gravity (ground-state susceptibility of the condensate). hAB is a massive spin-2 composite at $\sim$0.613 GeV, distinct from the GR graviton. It satisfies the Weinberg--Witten theorem and decays to 2$\gamma$ above threshold. Full derivation: Appendix J.

\section*{8  Kinematic Stability: Theorem KS}

\noindent The referee objection that ``Z6-algebraic stability does not imply kinematic stability'' is addressed here by a direct calculation across the complete Z6 composite spectrum.

\subsection*{8.1  The Cascade Mass Function}

\noindent The mass of a state at cascade level $\ell$ is given by the cascade formula:

\begingroup\centering m($\ell$) = $\Lambda$meta exp($-\ell$/be),   $\ell$ $\in$ $\mathbb{N}$\textsuperscript{+}.\hspace{2em}(10)\par\endgroup

\noindent This is a strictly decreasing function of $\ell$: each rung down the cascade suppresses the mass by a factor exp($-$1/be) = exp($-$1152/325) $\approx$ 1/34.6.

\noindent \textbf{Cascade level is not body number.} The level $\ell$ is the ladder rung, an independent quantum number; it is not a function of the body number N. Three anchors fix this: h\textsubscript{AB} (N = 2, $\ell$ = 8), X17 (N = 2, $\ell$ = 9) and the electron (N = 3, $\ell$ = 10) --- two states of equal N at different levels, and a heavier N at a lighter level (repository [2], \path{scripts/cascade_level_anchors.py}). Identifications at leading order: m(8) $\approx$ 610 MeV (nucleon and h\textsubscript{AB} scale), m(9) = 17.62 MeV (X17), m(10) = 0.5089 MeV (electron). The stable dark singlets of §14 are labelled by body number N $\in$ \{9, 12, 15, 18, 21, 27\}; their mass is \emph{not} m(N) but the Landau-pole value m\textsubscript{DM} = 2$\Lambda$meta/35 (§14.2), at the top of the cascade.

\subsection*{8.2  Theorem KS: Kinematic Stability}

\noindent \textbf{\textit{Theorem 8.1 (KS --- Kinematic Stability, [T1]). }}\textit{For any Z6 composite of body number N $\ge$ 2 and any partition N = N1 + N2 + $\cdots$ + Nk with k $\ge$ 2, Ni $\ge$ 1:}

\begingroup\centering $\Sigma$ m(Ni)/m(N) = $\Sigma$ exp((N $-$ Ni)/be) $\ge$ k exp((k $-$ 1)/be) $\ge$ 2 exp(1/be) $\approx$ 69.3.\hspace{2em}(11)\par\endgroup

\noindent \textit{In particular, $\Sigma$ m(Ni) > m(N) for ALL fragmentation decays, universally.}

\noindent Proof. Since N $-$ Ni = $\Sigma$\_\{j$\ne$i\} Nj $\ge$ k $-$ 1 (each Nj $\ge$ 1), each term exp((N $-$ Ni)/be) $\ge$ exp((k $-$ 1)/be). Summing k terms gives $\Sigma$ exp((N $-$ Ni)/be) $\ge$ k exp((k $-$ 1)/be), which for k $\ge$ 2 and Ni $\ge$ 1 satisfies $\ge$ 2 exp(1/be) $\approx$ 69.3 $\gg$ 1. $\square$

\begin{footnotesize}\begin{longtable}{|>{\raggedright\arraybackslash}p{0.2659\textwidth}|>{\raggedright\arraybackslash}p{0.2045\textwidth}|>{\raggedright\arraybackslash}p{0.2148\textwidth}|>{\raggedright\arraybackslash}p{0.2148\textwidth}|}
\hline
\textbf{Cascade level} & \textbf{m($\ell$)} & \textbf{Min k = 2 ratio} & \textbf{Min k = 3 ratio} \\ \hline
$\ell$ = 2 & 1.05 $\times$ 10\textsuperscript{9} GeV & 69.3 & --- \\ \hline
$\ell$ = 8 (nucleon / h\textsubscript{AB} scale) & 610 MeV & 2.88 $\times$ 10\textsuperscript{6} & 1.82 $\times$ 10\textsuperscript{9} \\ \hline
$\ell$ = 10 (electron) & 0.509 MeV & 9.96 $\times$ 10\textsuperscript{7} & 1.21 $\times$ 10\textsuperscript{11} \\ \hline
$\ell$ = 9 (X17) & 17.6 MeV & 5.12 $\times$ 10\textsuperscript{7} & 5.17 $\times$ 10\textsuperscript{9} \\ \hline
$\ell$ = 12 & 4.25 $\times$ 10\textsuperscript{-7} GeV & 3.45 $\times$ 10\textsuperscript{9} & 6.20 $\times$ 10\textsuperscript{12} \\ \hline
$\ell$ = 27 & 3.44 $\times$ 10\textsuperscript{-30} GeV & 3.67 $\times$ 10\textsuperscript{21} & 1.54 $\times$ 10\textsuperscript{28} \\ \hline
\end{longtable}\end{footnotesize}

\begingroup\centering\footnotesize\itshape Table 3: Minimum ratio $\Sigma$m($\ell$\textsubscript{i})/m($\ell$) of the cascade mass function over partitions $\ell$ = $\Sigma\ell$\textsubscript{i} into k parts (the minimum is the balanced partition; --- : no k-part partition exists). Rows are cascade levels, not dark-matter states (dark-matter mass: §14.2). Physical scope: see the status note on Theorem KS. Erratum: earlier prints gave 3597, 5.60 $\times$ 10\textsuperscript{11} and 8.05 $\times$ 10\textsuperscript{13} for the k = 3 entries at $\ell$ = 2, 8, 10; the first was the bound 3e\textsuperscript{2/be}, the other two were overstated. Inequality (11) holds for all.\par\endgroup

\noindent \textbf{Status note (under review).} Theorem KS, as stated, treats the mass as a function of body number and partitions body number, N = $\Sigma$N\textsubscript{i}. Fragmentation conserves psion (body) number, but the cascade level is independent of body number (§8.1), and no rule is yet derived that fixes the levels of the fragments from the level of the parent. The inequality (11) remains exact as arithmetic on the mass function; its application to physical fragmentation decays, and hence Corollary 8.2, is suspended pending that rule. Dark-matter stability (§14.1) does not depend on it: it rests on the exhaustive Z6 fragmentation census (Theorem DM.1, \path{proofs/P_dark_matter.py}) and on psion-number conservation. The module \path{proofs/P_kinematic_stability.py} verifies inequality (11) on the level partitions only and registers this gap as a scope guard. \textbf{v86:} the missing rule is identified as a conserved baryon (rung) label required by nucleon stability (§13.4); it is open problem P14-B.

\subsection*{8.3  Corollary: Z6 is Necessary and Sufficient}

\noindent \textbf{\textit{Corollary 8.2 ([T1]). }}\textit{If Z6 algebra forbids all decay channels of a composite of body number N (topological stability), then the composite is fully dynamically stable: all fragmentation decays --- whether Z6-allowed or not --- are kinematically forbidden by Theorem 8.1.}

\noindent Therefore: Z6 algebraic stability is both necessary and sufficient for full dynamical stability, including 4-momentum conservation.

\noindent \textbf{Scope and caveats. }Theorem KS covers fragmentation decays (N $\to$ $\Sigma$Ni, k $\ge$ 2). It does not cover:

\noindent 1. Radiative transitions (N $\to$ N + $\gamma$, same body number): excluded because m(N) is the ground-state mass by definition of the cascade. No lower-energy state of the same N exists within the PH spectrum.

\noindent 2. Weak rearrangements ($\Delta$N = 0, e.g. $\beta$ decay): covered separately by Z6 charge conservation and the selection rules of App. H. Rearrangement decays ($\Delta$N = 0, e.g. neutron $\beta$-decay: (0, 1, 5) $\to$ (0, 1, 2) + (0, 4, 5) + $\bar{\nu}$) are covered by Z6 charge selection rules (App. H) and are outside the scope of Theorem KS, which addresses fragmentation only.

\noindent \textbf{Anti-matter. }Anti-composites have the same body number N and the same cascade mass m(N) by CPT symmetry. Theorem KS applies to anti-composites by charge conjugation.

\noindent \textbf{T-level. }Theorem KS is [T1]: the proof is an algebraic identity on the cascade mass function. The only T2 input is m(N) = $\Lambda$meta exp($-$N/be) itself (which inherits T2 from the leading-order cascade). Verification in companion repository [2], module \path{proofs/P_kinematic_stability.py} (11 claims; 33 tests in \path{tests/test_kinematic_stability.py}): inequality (11) on all 1,295,920 partitions of every level 2 $\le$ $\ell$ $\le$ 50, the balanced-partition minimum, Table 3, the level/body-number scope guard, and the delegated DM.1 census. Its physical application is subject to the status note above.

\section*{9  Falsifiable Predictions --- Complete Registry}

\begin{footnotesize}\begin{longtable}{|>{\raggedright\arraybackslash}p{0.0708\textwidth}|>{\raggedright\arraybackslash}p{0.3135\textwidth}|>{\raggedright\arraybackslash}p{0.1820\textwidth}|>{\raggedright\arraybackslash}p{0.0910\textwidth}|>{\raggedright\arraybackslash}p{0.0708\textwidth}|>{\raggedright\arraybackslash}p{0.1719\textwidth}|}
\hline
\textbf{ID} & \textbf{Prediction} & \textbf{Value / Error} & \textbf{T-level} & \textbf{E} & \textbf{Experiment} \\ \hline
F1 & $\rho$$\Lambda$GN\textsuperscript{2} ($\Lambda$meta cancels) & 1.121 $\times$ 10\textsuperscript{-123} | 0.93\% & [T2+] & E3 & Precision cosmology \\ \hline
F2 & GN $\times$ me\textsuperscript{2} ($\Rightarrow$ $\delta$me/me = $-$$\tfrac{1}{2}$ $\delta$GN/GN) & 1.737 $\times$ 10\textsuperscript{-45} | 0.83\% & [T2] & E3/E1 & CODATA; clock drift \\ \hline
F6 & N$\nu$ = 3 (counting theorem) & 3 exact & [T1] & E2 & LEP; FCC-ee \\ \hline
F9 & Ng = 8 (|Z3|\textsuperscript{2} $-$ 1) & 8 exact & [T1] & E2 & Lattice QCD \\ \hline
F7 & Neff = 3.054 & 3.054 | <1$\sigma$ & [T2] & E3 & CMB-S4 (§9.1: non-discriminant vs SM) \\ \hline
F12 & $\eta$B $\leftrightarrow$ Neff via $\varepsilon$$\psi$ = 0.1479 & suspended (v86, §13.3) & --- & --- & CMB-S4 + BBN \\ \hline
F5a & No dark-matter state outside the 39 singlets at m\textsubscript{DM}: any direct-detection DM signal at m $\neq$ m\textsubscript{DM} falsifies & exclusion & [T2] & E2 & XENONnT; LZ; PandaX \\ \hline
F5b & $\sigma$\textsubscript{DM--n}(m\textsubscript{DM}) $\sim$ 10\textsuperscript{-54} cm\textsuperscript{2} (magnetic dipole), conditional on zero proto-colour coupling of dark singlets & null at all foreseeable exposures & [T4] & E1 & XENONnT; LZ; XLZD \\ \hline
F8 & UHECR excess above 10\textsuperscript{22} eV & above GZK & [T2] & E1 & Auger; GRAND \\ \hline
F13 & $\gamma$-line at 72 EeV & 72 EeV & [T2] & E1 & AugerPrime \\ \hline
F3 & $\nu_\tau$/$\nu_e$ asymmetry A\textsubscript{te} of the astrophysical flux, conditional on P-PMNS (§17.3); +7\% would require sin\textsuperscript{2}$\theta_{23}$ $\approx$ 0.574 with $\delta_{CP}$ $\approx$ $\pi$ & not derived & [T4] & E1 & IceCube-Gen2; DUNE, Hyper-K ($\theta_{23}$, $\delta_{CP}$) \\ \hline
F10 & GN phase transition at T = $\Lambda$meta & GN changes 10\textsuperscript{15}$\times$ & [T2] & E1 & CMB-S4 tensors \\ \hline
F4 & Proto-quark charges at Q\textsuperscript{2} $\sim$ $\Lambda$meta\textsuperscript{2} & DIS deviation & [T3] & E1 & Future DIS \\ \hline
F14 & Cstable(N) $\in$ \{29, 30, 31\} & $\sim$30/level & [T3] & E2 & Computational (internal) \\ \hline
F11 & DIS $\Sigma$e\textsuperscript{2}q = 47/27 at Q\textsuperscript{2} $\sim$ $\Lambda$meta\textsuperscript{2} & 47/27 & [T4] & E1 & Future DIS ($\gtrsim$ $\Lambda$meta) \\ \hline
F15 & X17 singlet at 17.62 MeV $\dagger$ & 17.62 MeV | 3.59\% & [T3] & E4 & ATOMKI vs MEG II/PADME $\dagger$ \\ \hline
F16 & hAB: $\gamma$ branching = 5 : 3 & 5/3 ratio & [T2] & E1 & Belle II \\ \hline
F17 & Co-drift $\delta$$\rho$$\Lambda$/$\rho$$\Lambda$ = $-$2 $\delta$GN/GN (joint with F2: one drift parameter, ratios $-$$\tfrac{1}{2}$ and $-$2) & exact ratios & [T2] & E1 & GN-drift (LLR/BBN/CMB) $\times$ w(z) constancy \\ \hline
\end{longtable}\end{footnotesize}

\begingroup\centering\footnotesize\itshape Table 4: All 17 falsifiable predictions of the PH (18 rows: F5 is split into F5a and F5b), with T-level (derivation status) and E-class (evidential status).\par\endgroup

\noindent \textbf{E-classification. }T-levels grade how a prediction is derived; E-classes grade what its agreement is worth. E1 --- novel, forward-looking, discriminating: made before the measurement, distinguishable from rival frameworks, data pending on a schedule; the only class with conversion power. E2 --- hard falsifier: forbidden or fixed by theorem; passing is shared with existing knowledge (zero conversion power), but violation kills the framework with no parameter to adjust --- the PH forbids by T1 theorem what the Standard Model merely does not contain. E3 --- postdiction of established data: a consistency scoreboard; necessary, never sufficient. E4 --- postdiction of a disputed anomaly: the weakest class; carried for completeness under quarantine (K.8 discipline). Current census: E1 $\times$ 10, E2 $\times$ 4, E3 $\times$ 3, E4 $\times$ 1 (F5 is split into F5a, E2, and F5b, E1). The framework's scientific weight rests on the E2 rigidity today and on the E1 set tomorrow; the E3 set, however precise, cannot convert.

{\footnotesize \noindent \textit{$\dagger$ X17 (status mid-2026): MEG II excludes the Atomki \textsuperscript{8}Be observation at 94\% CL (2025); PADME reports an independent 2.5$\sigma$ excess at 16.90 MeV --- mutually conflicting. If the PADME value consolidates, the F15 prediction (17.62) moves further from data (3.6\% $\to$ 4.3\%). MEG II (PSI), JINR (Dubna), and NA64 (CERN) found no signal at the Atomki claimed mass and coupling. Experimental status disputed; registered as T3 candidate pending independent confirmation.}\par}

\noindent F6 and F9 are T1 hard falsifiers --- any detection of a fourth neutrino or ninth gluon immediately refutes the framework. F7 and F12 are structurally linked through $\varepsilon$$\psi$ = 0.1479; this two-observable, one-parameter test is uniquely distinct from leptogenesis. The CF-CC structural prediction is discussed in Appendix K §K.8.

\noindent \textbf{P8b dependency.} None of F1--F17 above rests on the lepton-mass ratios $m_\mu/m_e$, $m_\tau/m_e$ or their P8b refinement (§12 Table 5, §A.3.1): those ratios are reported as structural results but are not themselves registered in this falsifiability census. Consequently, closure or non-closure of P8b Passo A.3 ([T2], App.~A) does not add or remove any entry from Table 4; it bears only on the precision claimed for the lepton-mass sector (Appendix L.4.4--L.4.5), not on the falsifiability registry.

\subsection*{9.1  Notes on Selected Predictions}

\noindent \textbf{F7 and F12 (Neff--$\eta$B structural link). }\emph{Status (v86): F12 is suspended, because the $\eta_B$ chain is withdrawn (§13.3); F7 is retained with the status note of §13.5.} Both Neff = 3.054 and $\eta$B = 6.02 $\times$ 10\textsuperscript{-10} arise from the single parameter $\varepsilon$$\psi$ = 0.1479 fixed by SU(2)meta geometry (§12). This structural link --- one parameter, two independent observables --- is uniquely testable by CMB-S4 combined with BBN likelihoods and cannot be reproduced by leptogenesis, which requires right-handed neutrinos absent from the PH spectrum. Honesty note: the PH$-$SM discriminant in Neff is 3.054 $-$ 3.044 = 0.010, below even the CMB-S4 target $\sigma$(Neff) $\approx$ 0.03; F12 can falsify PH and the SM jointly against a third value, but cannot separate them --- the $\varepsilon$$\psi$ two-observable link, not Neff alone, is the discriminating structure.

\noindent \textbf{F2 (variation-of-constants co-drift, honesty note). }The exact ratio $\delta m_e/m_e = -\tfrac{1}{2}\,\delta G_N/G_N$ is a genuine structural prediction (zero free parameters once $\Lambda$meta is fixed), but it is not by itself PH-specific: any scalar-tensor gravity in which $G_N$ and $m_e$ share a common dilaton-like scalar predicts \emph{some} fixed linear relation between the two drifts, generically. What is PH-specific is the particular coefficient $-\tfrac{1}{2}$ (forced by the cascade's $\Lambda$meta-dependence, eq. 6) and its joint closure with F17 ($\delta\rho_\Lambda/\rho_\Lambda=-2\,\delta G_N/G_N$, same single drift parameter, three observables, two fixed ratios). A measured pair of ratios inconsistent with $(-\tfrac{1}{2},-2)$ falsifies PH specifically; a measured pair consistent with some linear relation only confirms the general scalar-tensor class, not PH against its rivals. F2 in isolation is therefore weaker evidence than Table 4's E1 label suggests; the discriminating power is in the F2--F17 joint constraint, not F2 alone.

\noindent \textbf{F5 (direct detection): rationale for the F5a/F5b split and kill criterion, registered ex ante (2 Oct 2026). }\emph{Why split.} Earlier releases carried F5 as $\sigma$\textsubscript{DM--N} < 10\textsuperscript{-50} cm\textsuperscript{2} at [T2], E2. No derivation of that number exists in this paper or in the companion repository [2], so neither label is earned: [T2] requires a derived and verified mechanism, and E2 requires a prohibition fixed by theorem. The number is retired. What carries falsification power is structural and is kept as \textbf{F5a}: Theorem DM.1 ([T1]) and the canonical mass (§14.2, [T2]) leave no dark-matter state outside the 39 singlets at m\textsubscript{DM}, so a signal at any other mass is a component the PH spectrum lacks and contradicts the $\Omega$\textsubscript{DM} accounting of §14.3. It is a hard prohibition with nothing to adjust: [T2], E2. The rate is separated as \textbf{F5b} at [T4] (mechanism identified, dimensional estimates only). For the Landau-pole state (Q = 0, charged constituents, size $\sim$ 1/m\textsubscript{DM}): magnetic-dipole scattering $\sim$ 10\textsuperscript{-54} cm\textsuperscript{2}, charge-radius scattering $\sim$ 10\textsuperscript{-75} cm\textsuperscript{2}. Xenon sensitivity scales linearly with mass at high mass and reaches only $\sim$ 10\textsuperscript{-39} cm\textsuperscript{2} at m\textsubscript{DM} (rough extrapolation), so F5b predicts a null result with $\sim$ 15 orders of margin, with one exception: a coupling of dark singlets to the proto-colour mediator (the X17 family, §A.5) of nucleon-like strength gives $\gtrsim$ 10\textsuperscript{-38} cm\textsuperscript{2}. \emph{Promotion path.} [T3] once the proto-colour form factor of the N = 9 singlet is set up as a calculation (tied to P25, the Z6 composite gauge vertex); [T2] once that coupling is derived or shown to vanish (candidate argument to test: proto-colour singlets decouple at tree level, as neutral hadrons do from the photon). \emph{Kinematic fact.} Under the canonical mass assignment (§14.2) all 39 stable dark singlets sit at m\textsubscript{DM}; no stable neutral singlet exists at N $\le$ 6 (Theorem DM.1), and no other PH state is a halo candidate. The maximum xenon recoil at m\textsubscript{DM} is $\sim$ 1.7 MeV, so kinematics never excludes a signal: the rate does. At m\textsubscript{DM} = 7.2 $\times$ 10\textsuperscript{10} GeV the local number density is $\sim$ 4 $\times$ 10\textsuperscript{-12} cm\textsuperscript{-3}, so any detectable rate in a tonne-scale detector requires $\sigma$ $\gtrsim$ 10\textsuperscript{-39} cm\textsuperscript{2}, some fifteen orders above the F5b estimate. \emph{Kill criterion.} The PH dark sector is falsified by a dark-matter nuclear-recoil signal that (i) reaches $\ge$ 5$\sigma$ global significance in one experiment, (ii) is confirmed at $\ge$ 3$\sigma$ by an independent detector or independent exposure, and (iii) is not attributable to known backgrounds (neutrons, CE$\nu$NS, detector artefacts) by the collaborations' own analyses. Mass and operator do not matter: a signal at the Landau-pole mass violates F5b, and a signal at any other mass violates F5a. Not falsifying: single events, sub-5$\sigma$ excesses, electron-recoil-only excesses. \emph{Current status.} The LZ event of September 2026 (one nuclear recoil at 248 $\pm$ 23 (stat) $\pm$ 23 (sys) keV, 2.6$\sigma$ global, 3.4$\sigma$ local; arXiv:2609.02823) is below threshold and is carried as a watch item under E4 quarantine.

\noindent \textbf{F10 (gravitational phase transition). }Before the condensate forms (T > $\Lambda$meta), GN(T) $\sim$ 12$\pi$\textsuperscript{2}/(Nc$\Lambda$meta\textsuperscript{2}) $\sim$ 10\textsuperscript{15} GN\_obs. The Z6 condensation at T = $\Lambda$meta is a first-order gravitational phase transition reducing GN by 15 orders of magnitude. This imprints on primordial tensor spectra and is in principle detectable via the CMB-S4 tensor-to-scalar ratio.

\noindent \textbf{F11 (DIS structure function at $\Lambda$meta, [T4]). }$\Sigma$e\textsuperscript{2}q = 47/27. At Q\textsuperscript{2} $\ll$ $\Lambda$meta\textsuperscript{2}, PH and QCD give identical DIS predictions. At Q\textsuperscript{2} $\sim$ $\Lambda$meta\textsuperscript{2} $\sim$ 10\textsuperscript{24} GeV\textsuperscript{2}, the proto-quark charges +1/9, +4/9, +5/9 become individually resolvable, shifting the charge-weighted structure function sum from the SM value 1.00 to $\Sigma$e\textsuperscript{2}q = (1 + 16 + 25)/81 $\times$ 3 = 47/27 $\approx$ 1.74. Beyond current technology but precisely computable.

\noindent \textbf{G1--G3 (gravitational sector, App. K). }The driving frequency of [G3] is fdrive $\in$ [0.38, 0.43] GHz with Edrive $\approx$ 1.6 $\times$ 10\textsuperscript{-6} eV, fixed by self-consistency of the susceptibility relation (CN-1). The factor 2 in $\nu$J follows from order-parameter bilinearity (Lemma J.1, [T2+]). A quality factor Q $\sim$ 10\textsuperscript{10} (typical of Josephson junctions) would reduce the threshold driving energy density to $\sim$ 10\textsuperscript{24} GeV\textsuperscript{4}/m\textsuperscript{3} --- beyond current technology but computable. Classifications: [G1][T2+], [G2][T2] conditional on CEC, [G3][T3]. Full treatment: App. K.

\section*{10  Lepton Mass Spectrum}

\noindent \textbf{\textit{Theorem 10.1 (Ne = 10, [T1]). }}\textit{Ne = N$\gamma$ + Ng + N$\varphi$ = 1 + 8 + 1 = 10.}

\noindent Leading-order lepton mass formula (A = 13/12, p = 5/6 from SU(2)meta Clebsch--Gordan structure, [T2]):

\begingroup\centering mn/me = exp( A |mJ\_R|\textasciicircum{}p / be ).\hspace{2em}(12)\par\endgroup

\begin{footnotesize}\begin{longtable}{|>{\raggedright\arraybackslash}p{0.1595\textwidth}|>{\raggedright\arraybackslash}p{0.1367\textwidth}|>{\raggedright\arraybackslash}p{0.1937\textwidth}|>{\raggedright\arraybackslash}p{0.2620\textwidth}|>{\raggedright\arraybackslash}p{0.1481\textwidth}|}
\hline
\textbf{Lepton} & \textbf{|mJ\_R|} & \textbf{m\_pred/me} & \textbf{m\_pred} & \textbf{Error} \\ \hline
e\textsuperscript{-} & 0 & 1 & 0.511 MeV (anchor) & 0.0\% \\ \hline
$\mu$\textsuperscript{-} & 3/2 & 218 & 111.4 MeV & 5.4\% \\ \hline
$\tau$\textsuperscript{-} & 5/2 & 3795 & 1939 MeV & 9.1\% \\ \hline
\end{longtable}\end{footnotesize}

\begingroup\centering\footnotesize\itshape Table 5: Lepton masses (leading order, [T2]). Residual errors target P8b.\par\endgroup

\section*{11  Neutrino Masses}

\noindent SC input: $\Delta$m\textsuperscript{2}\textsubscript{21} = 7.530 $\times$ 10\textsuperscript{-5} eV\textsuperscript{2} --- equivalently the adimensional ratio $\Delta$m\textsuperscript{2}\textsubscript{21}/me\textsuperscript{2} = 2.908 $\times$ 10\textsuperscript{-16}, the framework's second (and only adimensional) empirical datum, serving purely as the neutrino-sector scale. Normal hierarchy and the mass ratios are the sector's predictive content. A closed-form derivation candidate, $\Delta$m\textsuperscript{2}\textsubscript{21} = (18/25) m(15)\textsuperscript{2} ($-$0.07\%), is registered at T4 in the companion repository (C-$\Delta$m21; promotion path P26-C); it is not promoted here.

\begin{footnotesize}\begin{longtable}{|>{\raggedright\arraybackslash}p{0.0931\textwidth}|>{\raggedright\arraybackslash}p{0.1966\textwidth}|>{\raggedright\arraybackslash}p{0.2276\textwidth}|>{\raggedright\arraybackslash}p{0.2276\textwidth}|>{\raggedright\arraybackslash}p{0.1552\textwidth}|}
\hline
\textbf{$\nu$} & \textbf{m$\nu$\_pred} & \textbf{$\Delta$m\textsuperscript{2}\_pred} & \textbf{$\Delta$m\textsuperscript{2}\_obs} & \textbf{Error} \\ \hline
$\nu$e & 7.50 $\times$ 10\textsuperscript{-9} eV & --- & < KATRIN & --- \\ \hline
$\nu$$\mu$ & 8.68 $\times$ 10\textsuperscript{-3} eV & 7.530 $\times$ 10\textsuperscript{-5} eV\textsuperscript{2} & 7.530 $\times$ 10\textsuperscript{-5} & 0.0\% (SC) \\ \hline
$\nu$$\tau$ & 4.73 $\times$ 10\textsuperscript{-2} eV & 2.239 $\times$ 10\textsuperscript{-3} eV\textsuperscript{2} & 2.453 $\times$ 10\textsuperscript{-3} & 8.7\% \\ \hline
\end{longtable}\end{footnotesize}

\begingroup\centering\footnotesize\itshape Table 6: Neutrino mass predictions.\par\endgroup

\section*{12  Cosmic Accountability}

\begin{footnotesize}\begin{longtable}{|>{\raggedright\arraybackslash}p{0.2011\textwidth}|>{\raggedright\arraybackslash}p{0.1915\textwidth}|>{\raggedright\arraybackslash}p{0.1819\textwidth}|>{\raggedright\arraybackslash}p{0.1436\textwidth}|>{\raggedright\arraybackslash}p{0.1819\textwidth}|}
\hline
\textbf{Observable} & \textbf{Predicted} & \textbf{Observed} & \textbf{Error} & \textbf{Status} \\ \hline
$\rho$$\Lambda$ (GeV\textsuperscript{4}) & 2.492 $\times$ 10\textsuperscript{-47} & 2.515 $\times$ 10\textsuperscript{-47} & 0.93\% & [T2+] \\ \hline
$\Delta$m\textsuperscript{2}\textsubscript{21} (eV\textsuperscript{2}) & 7.530 $\times$ 10\textsuperscript{-5} & 7.530 $\times$ 10\textsuperscript{-5} & 0.0\% & SC input \\ \hline
GN (GeV\textsuperscript{-2}) & 6.70883 $\times$ 10\textsuperscript{-39} & 6.70883 $\times$ 10\textsuperscript{-39} & 0.0\% & [SC]\footnote{This row is the single dimensional input (SC) itself, entered at its measured value by construction --- 0.0\% is not a prediction. The second GN row below is the independent, genuinely predicted quantity: GN recovered from the $\Lambda$meta$\to$me cascade run in reverse, at [T2], with a non-zero 0.83\% residual because it is not fitted to the input.} \\ \hline
me & 0.5089 MeV & 0.511 MeV & 0.42\% & cascade output \\ \hline
GN (GeV\textsuperscript{-2}) & 6.653 $\times$ 10\textsuperscript{-39} & 6.709 $\times$ 10\textsuperscript{-39} & 0.83\% & [T2]\footnotemark[\value{footnote}] \\ \hline
$\rho$$\Lambda$GN\textsuperscript{2} & 1.121 $\times$ 10\textsuperscript{-123} & 1.132 $\times$ 10\textsuperscript{-123} & 0.93\% & [T2+] \\ \hline
$\eta$B & withdrawn (§13.3) & 6.13 $\times$ 10\textsuperscript{-10} & --- & open; candidate §13.7 \\ \hline
$\Omega$DM/$\Omega$b & 5.43 / 5.64 & 5.364 & 1.2\% / 5.1\% & direct/inference; mutual 3.9\% (§14.3) \\ \hline
m$\mu$/me & 218 & 206.8 & 5.4\% & b-dep. \\ \hline
m$\tau$/me & 3795 & 3477 & 9.1\% & b-dep. \\ \hline
$\Delta$m\textsuperscript{2}\textsubscript{31} (eV\textsuperscript{2}) & 2.239 $\times$ 10\textsuperscript{-3} & 2.453 $\times$ 10\textsuperscript{-3} & 8.7\% & indep. \\ \hline
N$\nu$ & 3 & 3 & exact & theorem \\ \hline
Ng & 8 & 8 & exact & theorem \\ \hline
Neff & 3.054 & 2.99 $\pm$ 0.17 & < 1$\sigma$ & indep.; SM: 3.044 \\ \hline
\end{longtable}\end{footnotesize}

\begingroup\centering\footnotesize\itshape Table 7: Cosmic accountability table. SC = self-consistency input. b-dep. = sensitive to cascade parameter b at $\lesssim$ 35\% per 1\% change in b.\par\endgroup

\noindent \textbf{\textit{Remark 12.1 (Cascade output vs. empirical anchor). }}The cascade output is unique: me = $\Lambda$meta exp($-$Ne/be) = $\surd$(K($\xi$)/GN) exp($-$S/2g\textsuperscript{2} $-$ Ne/be) = 0.5089 MeV (0.42\% against the observed 0.511 MeV), since K($\xi$) sits inside $\Lambda$meta (eq. 6). A K($\xi$)-free cascade would give 0.4714 MeV ($-$7.8\%); the prefactor is part of $\Lambda$meta itself (eq. 6). In P8b, the observed 0.511 MeV serves only as a reference for computing lepton mass ratios m$\mu$/me, m$\tau$/me; it does not enter the GN $\to$ $\Lambda$meta $\to$ me derivation chain.

\section*{13  Baryogenesis: Status after the P14 Audit}

\noindent \textbf{Summary. }Earlier releases carried $\eta_B = 6.02\times10^{-10}$ as a derived output at [T2]. The P14 audit withdraws that claim. Equation (1) conserves psion number exactly, field-theory CP is exact, and the Standard-Model sphaleron conversion used in the former eq.~(14) does not apply in the PH. Neither eq.~(1) nor its minimal extension produces the asymmetry (13.3). One candidate mechanism remains, level bifurcation in the pair production of unconfined composites after inflationary reheating (13.7). Its structure follows from symmetry, but its magnitude is not yet explained. This section replaces the former 13.1--13.5.

\subsection*{13.1  What Equation (1) Conserves}

\noindent \textbf{Psion number. }$N_\psi$ is the Noether charge of the overall phase and is exactly conserved by every term of eq.~(1). In the broken phase it is the only exact additive charge (exact charge audit over 20 triples and four condensate variants, zero further candidates).

\noindent \textbf{Continuous internal symmetry. }The full Lie algebra of internal transformations leaving $(\bar\psi J^a\psi)(\bar\psi\Sigma^a\psi)$ invariant is two-dimensional and abelian, $U(1)_N\times U(1)_M$, with $M=\sum m_k$; no continuous chiral extension survives the $G_{\rm UV}$ term. The winding condensate of App.~J.6.2 breaks $U(1)_M$, leaving a residual discrete charge $D=(2M-6Q+3)\bmod 4$.

\noindent \textbf{CP. }The coupling tensor is real and $\varepsilon_\psi$ is CP-even: field-theory CP is exact. The Z6 vacua break CP spontaneously, but conjugate vacua are degenerate ($\Delta C(s)$ symmetric).

\noindent \textbf{Antimatter. }Only one reading is consistent with observed reactions: antiparticles are composites of antipsions ($N_\psi=-3$, charge $-Q$). The index inversion $k\to(6-k)\bmod 6$ (App.~A.4) describes the Z6 label of the antipsion, not a psion composite. The alternative reading (antiparticle $=$ psion triple with inverted indices) fails psion-number conservation already in $\beta^\pm$ decay, fails charge conservation in $n\to pe^-\bar\nu$ and $\mu^-\to e^-\bar\nu_e\nu_\mu$, and gives $\bar\Sigma^+$ the charge $+1$.

\subsection*{13.2  The Sakharov Conditions [16], Revisited}

\noindent \textbf{(1) Baryon-number violation. }Nothing built from psions violates $N_\psi$, including an emergent sphaleron. With $N_\psi/3=B+L$ (13.4), $B+L$ is conserved and only $B-L$ can change: the reverse of the Standard Model. The sphaleron factor 8/23 and the coefficient $C_{\rm core}$ of the former eq.~(14) therefore have no basis in the PH. In addition, $C_{\rm core}=8.4434\times10^{-4}$ could not be reproduced from its stated inputs.

\noindent \textbf{(2) C and CP violation. }Eq.~(1) supplies none (13.1). The former statement that the chirality-flip amplitude $\propto(m_e/\Lambda_{\rm meta})^2$ is a CP-odd phase is withdrawn: a real amplitude carries no phase.

\noindent \textbf{(3) Departure from equilibrium. }Available: the formation of composites is out of equilibrium when the contact rate is comparable to the expansion rate (13.7). Temperature histories alone redistribute abundances among cascade levels but never create an asymmetry.

\subsection*{13.3  Withdrawn Results and Closed Mechanisms}

\noindent \textbf{Withdrawn. }(i) The former eqs.~(14)--(15), $\eta_B = C_{\rm core}(g_{1/2}/g_S)\varepsilon_\psi = 6.02\times10^{-10}$: the sphaleron input has no PH basis, and $C_{\rm core}$ is not reproducible. (ii) Eqs.~(37)--(38) of App.~J: $\xi^*=434.09$ was computed from the target $\eta_B$, so eq.~(38) is an identity, not a check. Independent routes to $\xi^*$ give $\xi^*\le1.08$ (Potts Z3; second-order NJL transition with $T_c=9.6$ MeV) and $\xi^*=3.6$--$39$ (Coleman--Weinberg with thermal inflation). (iii) The claim that the PH satisfies all three Sakharov conditions without additional ingredients. (iv) Prediction F12 is suspended (§9).

\noindent \textbf{Reclassified. }$\varepsilon_\psi=25/169$ remains a result of SU(2)$_{\rm meta}$ geometry (App.~J, conditional on P23). Its link to $\eta_B$ is not established.

\noindent \textbf{Closed mechanisms. }Each was tested under a pre-registered definition, and each fails for a stated reason:
\begin{itemize}
\item Emergent sphalerons: they conserve $N_\psi$.
\item Z6 domains with spontaneous CP: they form at $T\approx3$ MeV, when assembly is frozen ($\Gamma/H=3\times10^{-37}$). Their maximal comoving size today is $\sim30$ pc, against the $\gtrsim1$ Gpc required of antimatter regions [19]. Degenerate conjugate vacua give a zero global average.
\item Temperature histories: no asymmetry from symmetric initial data.
\item CP terms built from PH operators: 16 commutator and rearrangement patterns, alone or in pairs, fit the observed composition no better than random directions.
\item A phase inserted in $J\!\cdot\!\Sigma$: it generates an asymmetry, but exotic triples receive 2--2500 times the proton asymmetry. Without a conserved label, no CP source selects the observed composites.
\item Species mixing in the minimal extension of 13.4: the proton lifetime bounds its coefficient to $c\le1.3\times10^{-8}$, which places the active window above $2.7\times10^{15}$ GeV, outside the contact description.
\end{itemize}

\subsection*{13.4  Baryon Number and the Level Rule}

\noindent \textbf{What observation requires. }With the cascade level treated as a quantum number, the observed reactions $\beta^\pm$ decay, $e^-$ capture, $\pi^+\to\mu^+\nu_\mu$, $\mu$ decay, $\nu_\mu n\to\mu^-p$, $\bar\nu_ep\to e^+n$, $\Lambda\to p\pi^-,n\pi^0$, $\Sigma^0\to\Lambda\gamma$ and neutrino oscillations require exactly two additive charges: $B$ (states of the level-8 class) and $L$ (level-10 and neutrino classes), with $N_\psi/3=B+L$.

\noindent \textbf{What eq.~(1) allows. }With only $N_\psi$ and $Q$ conserved, seven unobserved processes are allowed: $p\to\nu\pi^+$, $n\to\nu\pi^0$, $n\to3\nu$, $p\,e^-\to\nu_\mu\nu_\tau$, and three level jumps within one triple ($n\to\nu_\tau\gamma$, $\Lambda^0\to\nu_e\gamma$, $\Sigma^0\to\nu_\mu\gamma$). All have $\Delta B=-\Delta L$. Against the measured limits ($p\to\bar\nu\pi^+>3.9\times10^{32}$ yr and $n\to\bar\nu\pi^0>1.1\times10^{33}$ yr [17]; $n\to$ invisible $>5.8\times10^{29}$ yr [18]), the required amplitude suppression is $\lesssim10^{-8}$ relative to a contact operator at $\Lambda_{\rm meta}$. For same-triple jumps it is $\lesssim5\times10^{-30}$ relative to an electromagnetic transition (reference lifetime $10^{29}$ yr).

\noindent \textbf{What eq.~(1) cannot supply. }$B$ is not a Noether charge: it lies outside the span of $(M,Q,N_\psi)$. It is not topological: the vacuum manifold has $\pi_2=\pi_3=0$, so no skyrmion-like baryon number exists [21]. And the level is not an internal degree of freedom, since states are labelled by the triple alone. Hence the dual-sector identification of \S6 and L.4.4, in which the same triple is either a baryon or a neutrino ((0,1,5) $=n/\nu_\tau$, (0,2,4) $=\Sigma^0/\nu_\mu$, (1,2,3) $=\Lambda^0/\nu_e$), is incompatible with neutron stability as eq.~(1) stands. Correction to \S6: there are five neutral three-body triples, not three. They are (0,1,4), (0,1,5), (0,2,4), (0,2,5) and (1,2,3); the two additional ones carry $D=2$.

\noindent \textbf{Minimal extension (registered, not adopted). }A two-valued species index, $\psi\to\psi_s$ with $s\in\{B,L\}$ and the currents of eq.~(1) summed over species, preserves $U(1)_B\times U(1)_L$ if $SU(2)_{\rm species}$ is broken. All observed reactions are allowed, and all seven unobserved processes are forbidden exactly. Its costs:
\begin{itemize}
\item The $L$ sector contains six positively charged matter leptons that must be heavy.
\item The numerical predictions survive only if the species index does not enter $|\beta_0|=N_{\rm Dirac}+2/N_f$. Doubling $N_f$ gives $m_e=0.065$ MeV.
\item A label carried by the constituents cannot be chosen at formation.
\end{itemize}
This is recorded as open problem \textbf{P14-B}: a conserved baryon (rung) label.

\subsection*{13.5  Link to $N_{\rm eff}$ (status note)}

\begingroup\centering $\Delta$Neff\_PH = 3 $\cdot$ (15/7$\pi$\textsuperscript{2}) $\cdot$ ($\mu$$\nu$/T)\textsuperscript{2} = 0.010,   Neff\_PH = 3.054.\hspace{2em}(16)\par\endgroup

\noindent Equation (16) is retained, but its shared-parameter status with $\eta_B$ lapses with 13.3. If $\mu_\nu/T=0.123$ is read as a lepton-number asymmetry, it conflicts with psion-number bookkeeping. With total $N_\psi=0$, $B=-L$ and charge neutrality, the antineutrino excess is two per hydrogen atom, a degeneracy $\xi_\nu\approx2.5\times10^{-9}$, some $5\times10^7$ times smaller than 0.123. Equation (16) must therefore rest on a different reading of $\mu_\nu$ or be reclassified.

\subsection*{13.6  Structural Predictions}

\noindent \textbf{(a) }No nucleon decay with $\Delta B=\Delta L$: $p\to e^+\pi^0$ and similar modes change $N_\psi$ and are forbidden exactly [T1]. \textbf{(b) }Any nucleon decay, if observed, has $\Delta B=-\Delta L$ ($p\to\nu\pi^+$, $n\to\nu\pi^0$, $n\to3\nu$). \textbf{(c) }No $0\nu\beta\beta$ and no $n$--$\bar n$ oscillation, since both change $N_\psi$. \textbf{(d) }If the asymmetry formed with total $N_\psi=0$: a cosmic antineutrino excess of two per hydrogen atom. \textbf{(e) }Dark matter as a thermal relic of the cascade top would require level $\approx4.8$ ($\approx56$ TeV); the canonical $m_{\rm DM}=7.2\times10^{10}$ GeV requires a non-thermal origin (§14.3).

\subsection*{13.7  A Candidate Mechanism: Level Bifurcation at Reheating}

\noindent \textbf{Unconfined spectrum. }By the confinement rule of App.~H/I (summed index $I \bmod 6\in\{0,3\}$), exactly eight triples are unconfined: (0,1,2), (0,1,5), (0,2,4), (0,4,5), (1,2,3), (1,3,5), (2,3,4) and (3,4,5). These are the triples with one constituent of each proto-colour, and they are the states identified with observed particles. All fractional-charge triples and both $D=2$ neutrals are confined. Exotic triples therefore do not constrain any mechanism below $\Lambda_{\rm meta}$.

\noindent \textbf{Postulates. }Below $\Lambda_{\rm meta}$ only these composites and radiation exist, so composites are pair-produced from radiation. Each triple has four channels, and $t_B\bar t_L$ and $t_L\bar t_B$ are CP conjugates from a self-conjugate initial state, the minimal setting for CP violation compatible with CPT and unitarity [20]. (P1) The rung (baryonic $t_B$ or leptonic $t_L$) is fixed at production and conserved by every process except annihilation back into radiation. (P2) The conjugate channels are biased by $\varepsilon_t$, with source $\varepsilon_t$ times the net production rate, so the source vanishes in equilibrium. (P3) Production occurs after reheating at $T_{\rm RH}<\Lambda_{\rm meta}$ with the natural coupling, so a single condensate vacuum covers the observable Universe.

\noindent \textbf{Structural consequences ([T1] given P1--P2). }Each production event changes $B$ and $L$ oppositely and creates a triple with its antitriple. Hence $B=-L$, total charge zero and $\Delta N_\psi=0$, and the late-time state is hydrogen plus two net antineutrinos per baryon.

\noindent \textbf{CP source. }The vacuum phase of the winding condensate changes $M$ by two units, and $U(1)_M$ is exact in eq.~(1); since the condensate itself generates the bond of J.6.3, a $U(1)_M$ rotation removes the phase exactly in every vacuum. The phase is physical only relative to an explicit breaking of strength $\delta$, with CP-odd effects $\propto\delta^2$. The only explicit breaking in the PH, $\kappa_{\rm CP}=(m_e/\Lambda_{\rm meta})^2$ (K.10), gives a bias $\sim10^{-62}$. The PH as written has no CP source of adequate size.

\noindent \textbf{Efficiency and open problem. }With the natural coupling, $\eta_B/\varepsilon$ reaches $\sim0.1$ for $T_{\rm RH}\sim4$--$7\times10^9$ GeV, and the asymmetry is washed out above $\sim1.5\times10^{10}$ GeV. The result is insensitive to gradual reheating (peak temperature up to 60 $T_{\rm RH}$). The observed $\eta_B$ needs a bias of $10^{-7}$--$10^{-6}$, far above what the PH supplies, and the sign depends on the binding law and on $T_{\rm RH}$. The reheating temperature is not a prediction.

\noindent \textbf{Status ([T3]) and conditions. }(C1) Derive the rung label, for instance as a winding charge changed only through the deconfined core, where the action per mode $S/N_e$ gives a suppression $e^{-21\Delta\ell}$, enough for 13.4. (C2) Identify a derived explicit breaking of $U(1)_M$ (or another CP-odd ingredient) of strength $\delta\sim10^{-3}$. (C3) Derive the pair-production rate of composites.

\section*{14  Dark Matter: Identity, Mass, and Relic Density}

\subsection*{14.1  Dark Matter as Stable Z6 Singlets}

\begin{footnotesize}\begin{longtable}{|>{\raggedright\arraybackslash}p{0.5108\textwidth}|>{\raggedright\arraybackslash}p{0.0973\textwidth}|>{\raggedright\arraybackslash}p{0.2919\textwidth}|}
\hline
\textbf{State} & \textbf{I} & \textbf{Stability} \\ \hline
(0, 0, 1, 1, 2, 2, 3, 4, 5) & 0 & Ground configuration (unique s-wave state: 3 forced singlets + free-spin triplet) \\ \hline
(0, 0, 0, 0, 0, 0, 1, 1, 4) & 0 & Neutral singlet (census; high shell cost) \\ \hline
(0, 0, 0, 0, 0, 1, 1, 1, 3) & 0 & Neutral singlet (census; high shell cost) \\ \hline
\end{longtable}\end{footnotesize}

\begingroup\centering\footnotesize\itshape Table 8: Representative neutral N = 9 singlets (generated census, \path{gen_composite_tables.py} [2]); stability filter per [2], pending the parity-correct re-enumeration (P-Pauli).\par\endgroup

\noindent \textbf{Theorem DM-S (statistics by parity, [T1]). }All stable dark singlets carry Q = 0 (neutral by the Z6 charge formula), with statistics fixed by body-number parity: fermionic at N $\in$ \{9, 15, 21, 27\}, bosonic at N $\in$ \{12, 18\}; scalar states are possible only at even N. DM is dark by group-theoretic necessity (Q = 0), and its stability is protected independently of spin by psion-number conservation (gauge bosons are $\psi$$\bar{\psi}$ bilinears with net psion number zero), Theorem KS, and --- for odd N --- fermion-number superselection.

\noindent \textbf{Absolute stability (Theorem KS + App. G, [T1]): }By App. G (exhaustive enumeration), no Z6-allowed decay channel exists. By Theorem KS (§8), ALL fragmentation decays N $\to$ \{Ni\} are kinematically forbidden by the cascade mass function, independently of Z6 quantum numbers ($\Sigma$m(Ni)/m(N) $\ge$ 69 for any 2-body decay). Together these establish full dynamical stability: Z6 algebraic stability is both necessary and sufficient.

\noindent \textbf{DM stability: }39 stable neutral singlets at N $\in$ \{9, 12, 15, 18, 21, 27\} by exhaustive enumeration (dynamic programming + brute force, two independent algorithms; termination proved to N = 48). Mass scale: §14.2. Relic abundance: §14.3.

\subsection*{14.2  Dark Matter Mass from the Landau Pole}

\begingroup\centering mDM = 2$\Lambda$meta/NLandau = 2 $\times$ 1.2609 $\times$ 10\textsuperscript{12}/35 = 7.205 $\times$ 10\textsuperscript{10} GeV.\hspace{2em}(17)\par\endgroup

\noindent NLandau = 35 counts cascade modes to the Landau pole (distinct from Ne = 10 binding levels, Theorem G.1).

\subsection*{14.3  Relic Density Derivation}

\noindent Two partially independent routes to $\Omega$DM/$\Omega$b; both are quoted against the observed ratio 5.364 = $\Omega$c h\textsuperscript{2}/$\Omega$b h\textsuperscript{2} (Planck 2018 base chain: $\Omega$c h\textsuperscript{2} = 0.1200, $\Omega$b h\textsuperscript{2} = 0.02237, h = 0.6736).

\noindent \textbf{Inference route. }$\Omega$DM = 1 $-$ $\Omega$$\Lambda$(PH) $-$ $\Omega$b(PH), with both terms computed from PH predictions and no observed cosmological density imported: $\Omega$$\Lambda$(PH) = $\rho$$\Lambda$,pred/$\rho$\_crit(h) = 2.492 $\times$ 10\textsuperscript{-47}/3.673 $\times$ 10\textsuperscript{-47} = 0.6784, and $\Omega$b(PH) = $\eta$10,pred/(273.9 h\textsuperscript{2}) = 0.0484. Result: $\Omega$DM/$\Omega$b = 0.2732/0.0484 = 5.64 (5.1\% error) [T2]. Declared inputs: h = 0.6736, TCMB = 2.7255 K, flatness. This route makes no assumption about the DM production mechanism.

\noindent \textbf{Status (v86). }The direct route below uses $\xi^*$ = 434.09, obtained from the predicted $\eta$B through eq.~(33); both are withdrawn (§13.3). The direct route is therefore carried as unsupported until $\xi^*$ is derived independently, and the mutual-consistency and sensitivity statements below are conditional on it. The inference route is unaffected if the measured $\Omega$b is used.

\noindent \textbf{Direct route. }$\Omega$DM/$\Omega$b = (4$\pi$/3) $\times$ (mDM/mp) $\times$ e\textsuperscript{$-$$\xi$DM} = 5.43 (1.2\% error) [T2+], with $\xi$DM = (mDM/$\Lambda$meta) $\times$ $\xi$* = (2/35) $\times$ 434.09 = 24.805 and $\xi$* = 434.09 obtained self-consistently from the predicted $\eta$B = 6.02 $\times$ 10\textsuperscript{-10} via §13.3. The exponential factor e\textsuperscript{$-$$\xi$DM} is derived from the baryogenesis sector and is proved independent of TCMB: comoving yields Y = n/s eliminate s(T) in the ratio, and Tnuc cancels in the nucleation-epoch density ratio. The prefactor 4$\pi$/3 has geometric motivation (SO(3) isotropy) and is under active derivation (P-4$\pi$/3, [T2+]).

\noindent \textbf{Mutual consistency. }The two routes share the baryogenesis output $\eta$B (which sets $\Omega$b in the inference route and $\xi$* in the direct route); their discriminating second ingredients are $\rho$$\Lambda$ (inference) and the Landau-pole mass mDM (direct). Mutual agreement: 3.9\% --- fully accounted by propagation of the upstream Class-A deficits through the route elasticities (see Sensitivity below): the inference route amplifies the $\rho$$\Lambda$ deficit ($-$0.92\% in $\Omega$$\Lambda$, elasticity $-$2.5) and the $\eta$B deficit ($-$1.75\% in $\Omega$b, elasticity $-$1.2), while the direct route carries only the mild GN\textsuperscript{$-$1/2} dependence of mDM. The spread is error propagation, not an independent tension; no new problem is registered. Both routes are independent of be (Class A).

\noindent \textbf{Sensitivity. }The direct route is exponentially sensitive to the $\eta$B convention: d ln($\Omega$DM/$\Omega$b)/d ln $\eta$B = +$\xi$DM/3 = +8.27, so switching $\xi$* from the predicted $\eta$B (434.09) to a measured one (434.6) moves the direct value from 5.43 to 5.27 --- a 2.9\% swing, larger than the quoted agreement itself. The predicted-$\eta$B convention is canonical (zero-parameter purity; §J.4). The inference route elasticities are $-$2.5 (wrt $\Omega$$\Lambda$) and $-$1.2 (wrt $\Omega$b). Headline percentages in this section are convention-conditional at these levels.

\subsection*{14.4  Experimental Signatures}

\noindent Direct detection: no signal at any mass other than m\textsubscript{DM} (F5a), and a rate $\sim$ 10\textsuperscript{-54} cm\textsuperscript{2} at m\textsubscript{DM} conditional on zero proto-colour coupling (F5b, [T4]); see §9.1 for the kill criterion. Monochromatic $\gamma$-line DM2 + DM2 $\to$ $\gamma$$\gamma$ at 72 EeV (F13; AugerPrime). UHECR excess above 10\textsuperscript{22} eV (F8; Auger, GRAND). XRISM (2025) null keV result consistent with PH; falsifies keV sterile-$\nu$ DM.

\section*{15  Open Problems}

\begin{footnotesize}\begin{longtable}{|>{\raggedright\arraybackslash}p{0.1345\textwidth}|>{\raggedright\arraybackslash}p{0.3931\textwidth}|>{\raggedright\arraybackslash}p{0.1241\textwidth}|>{\raggedright\arraybackslash}p{0.2483\textwidth}|}
\hline
\textbf{ID} & \textbf{Problem} & \textbf{Status} & \textbf{Blocks} \\ \hline
P23 & $\kappa$ = 1 normalisation identity (Lemma I.1') & [T2+] & $\varepsilon$$\psi$ $\to$ T1 ($\eta$B link withdrawn, §13.3) \\ \hline
P14-B & Conserved baryon (rung) label; sign and size of the bifurcation bias (§13.4, §13.7) & open & $\eta$B, nucleon stability, Theorem KS \\ \hline
P14-C & Baryon binding scale ($(\hslash c)/a\approx163$ MeV) and internal structure consistent with $J\!\cdot\!\Sigma$ mixing; hexagonal radii are phenomenological (App.~L.4.7) & open & hadron radii, $m_n-m_p$ \\ \hline
P-PMNS & Sector-bridging PMNS texture (e\textsuperscript{-}: I = 3; $\nu$: I = 0); the CC bridge is the proto-W (I = 3) --- texture open; flavour basis fixed, eq. (1) alone gives sin\textsuperscript{2}$\theta_{23}$ $\approx$ 0.21 (§17.3) & [T3] & mixing sector; F3 \\ \hline
P-|O| & Overlap sector for absolute rates; first datum: |O|\textsuperscript{2}(CF) $\le$ 0.370 (2$\sigma$) & [T2] & absolute rates $\to$ T1 \\ \hline
P-$\varepsilon$conf & Formation-overlap democracy behind $\varepsilon$conf = 4/9; bounded empirically at 1--2\% via the +8.27 lever & [T2] & withdrawn with eq.~(33) (§13.3) \\ \hline
P-Pauli & Pauli structure restated: auto-singlets from N $\ge$ 7; spatial shells from N $\ge$ 13 (nuclei at N = 15+); per-level structure open; 39-count re-audit (floor) & [T2] & DM spin sector \\ \hline
C-Q-rat & Charge rationality / integer quantization as Z6 selection principle & [T3] & why Nc = 6 \\ \hline
F-fork & RESOLVED $\to$ twelve-state bookkeeping: the U3 decider fired ((2, 2$\pi$/3)); residuals R-pair (auto-singlet + uniqueness + lepton anchor) and R-norm (bias/formation split, 0.9\% seal) reviewed and held & [T2] & closed \\ \hline
C-web & Closed-form web RESOLVED: S definitional; $\rho$$\Lambda$ = C$\cdot$me\textsuperscript{6}/$\Lambda$meta\textsuperscript{2} exact rearrangement of Theorem 3.1; O(1/S) = two-loop registration & [T2+] & master invariant; F17 \\ \hline
P26 & $\nu$-quantum number locking (ce = 11/6 vs 2) & [T2] & P23 (locking mechanism) \\ \hline
P-4$\pi$/3 & Prefactor 4$\pi$/3 in $\Omega$DM/$\Omega$b direct & [T2+] & $\Omega$DM $\to$ T1 \\ \hline
P1 & be = 325/1152 from 2PI kernel & [T2] & mass ratios Class B \\ \hline
P2/P12 & $\Phi$($\Lambda$) = $\Lambda$ without me; closes all adimensional ratios; 1 dimensional datum remains (irreducibility) & [T2] & me $\to$ prediction \\ \hline
\end{longtable}\end{footnotesize}

\begingroup\centering\footnotesize\itshape Table 9: Twelve critical open problems. Complete registry: Appendix K.7 and psion-proofs.\par\endgroup

\section*{16  The Psion in Hawking Radiation}

\noindent The Z6 condensate modifies black hole thermodynamics when the Hawking temperature approaches $\Lambda$meta. The Planck mass MPl = ($\hslash$c/GN)\textasciicircum{}\{1/2\} = 1.22 $\times$ 10\textsuperscript{19} GeV is the scale at which quantum gravitational effects become order-unity.

\noindent \textbf{• Critical mass: }M* = MPl\textsuperscript{2}/(8$\pi$$\Lambda$meta) = 8.38 g --- the black hole mass below which the horizon temperature reaches $\Lambda$meta.

\noindent \textbf{• Deconfinement radius: }rdc = (GN M/$\Lambda$meta\textsuperscript{4})\textasciicircum{}\{1/3\} for M > M*.

\noindent \textbf{• Modified evaporation rate: }(dM/dt)PH/(dM/dt)SM = 127.75/106.75 = 1.197 (19.7\% faster above M*).

\noindent Full derivation of M* = 8.38 g, deconfinement radius, and the 19.7\% modified Hawking rate: companion PH P19 Black Holes (psion-proofs [2]).

\section*{17  Exotic Signatures and Experimental Context}

\subsection*{17.1  ANITA: Z6 Exotic Carrier}

\noindent The ANITA balloon experiment observed anomalous upward-going radio pulses inconsistent with standard neutrino interactions. The PH identifies the carrier as $\nu$S = (1, 4, 4): a colour-singlet (I = 3) composite with Q = +1/3, M $\approx$ 7.5 GeV, via:

\begingroup\centering $\nu$S(1, 4, 4) + n(0, 1, 5) $\to$ $\tau$\textsuperscript{-}(0, 1, 5)DI + X(0, 0, 3).\hspace{2em}(18)\par\endgroup

\noindent Composite identification: Appendix H (PD protocol). Mass M $\approx$ 7.5 GeV: §6 four-body mass formula. The carrier is charged (Q = +1/3) and hence not sterile; the identification is carried under E4 quarantine (outside the Table 4 census) pending confrontation with fractional-charge limits (ArgoNeuT, milliQan).

\subsection*{17.2  CF-CC Channel: Structural Prediction (see Appendix K)}

\noindent The Chiral-Flip Charged Current (CF-CC) is a channel predicted by the Z6 composite structure:

\begingroup\centering $\nu$$\mu$(0, 4, 2) + n(0, 1, 5) $\to$ e\textsuperscript{-}(0, 1, 2) + p(0, 4, 5).\hspace{2em}(19)\par\endgroup

\noindent Structure ([T1]): selection rules and existence. The absolute rate factorizes as RCF = NDirac g\textsuperscript{2}meta (me/m$\mu$)\textsuperscript{2} $\times$ |O|\textsuperscript{2} = 6.16 $\times$ 10\textsuperscript{-4} |O|\textsuperscript{2}, with |O|\textsuperscript{2} the BSE formation overlap (P-|O|); the rate is energy- and baseline-independent at every accessible order (Theorem K.8.1). MicroBooNE bounds |O|\textsuperscript{2} $\le$ 0.370 (2$\sigma$). Full treatment and the e\textsuperscript{+}e\textsuperscript{-} uniqueness theorem: Appendix K §K.8.

\subsection*{17.3  IceCube Tau Flavor Ratio (F3, conditional)}

\noindent F3 concerns the $\tau$--$e$ asymmetry of the astrophysical flux at Earth, A\textsubscript{te} = (f\textsubscript{$\tau$} $-$ f\textsubscript{e})/(f\textsubscript{$\tau$} + f\textsubscript{e}) for a pion-decay source. This quantity is fixed by the PMNS matrix alone. With standard oscillations it spans roughly $-$10\% to +9\% over the NuFIT 6.0 ranges; +7\% corresponds to sin\textsuperscript{2}$\theta_{23}$ $\approx$ 0.574 with $\delta_{CP}$ $\approx$ $\pi$, and with real mixing ($\delta$ $\in$ \{0, $\pi$\}) $\delta$ = 0 gives A\textsubscript{te} < 0 for every allowed $\theta_{23}$. The PH does not derive the PMNS matrix (open problem P-PMNS), and no derivation of the 7\% figure is contained in this release; earlier versions of this section cited one in [2]. F3 is therefore carried at [T4], conditional on P-PMNS.

\noindent \textbf{Status of the PMNS programme} ([2], \path{proofs/docs/PMNS_PROGRAMME.md}). The flavour basis is fixed exactly: the generation of a neutral triple is the proto-colour of its single right-handed constituent, and the charged-current flip is colour-diagonal. With the charged leptons taken diagonal in that basis, both terms of eq. (1) together with the locking gap give sin\textsuperscript{2}$\theta_{12}$ close to the measured value, but sin\textsuperscript{2}$\theta_{23}$ $\approx$ 0.21 and A\textsubscript{te} $\approx$ $-$34\%, the opposite sign to +7\%. The charged-lepton rotation, which could change $\theta_{23}$, requires the full P8b transfer matrix and has not been tested.

\noindent \textbf{Experimental context.} IceCube's flavour measurement (f\textsubscript{e} : f\textsubscript{$\mu$} : f\textsubscript{$\tau$} = 0.30 : 0.37 : 0.33, 5 TeV--10 PeV) does not constrain A\textsubscript{te} at the required level. IceCube-Gen2 and KM3NeT target $\sim$5\% precision on the flavour ratio, while DUNE and Hyper-K will measure $\theta_{23}$ and $\delta_{CP}$ directly, which determine A\textsubscript{te} in any case.

\subsection*{17.4  Sterile Neutrino DM vs. PH}

\noindent N$\nu$ = 3 (T1 theorem) forbids a fourth neutral fermionic stable composite. XRISM (2025) found no keV X-ray lines. PH and sterile-neutrino DM are mutually exclusive. N$\nu$ = 3 T1 proof: Theorem 6.1 (§6); Appendix G §G.5.2; prediction F6 (§9).

\subsection*{17.5  Neff and Baryogenesis}

\noindent $\varepsilon$$\psi$ = 0.1479 generates $\delta$L = 0.370, giving $\mu$$\nu$/T = 0.123 and Neff = 3.054 (F12). Structurally distinct from leptogenesis: no right-handed neutrinos. The $\eta$B link of this chain is withdrawn in v86 (§13.3); Neff is retained with the status note of §13.5. Predictions F7 and F12 (§9). Structural derivation of $\varepsilon$$\psi$: Appendix J §J.6.

\section*{18  Discussion}

\subsection*{18.1  Class A vs. Class B Results}

\noindent Class A (independent of b): $\Omega$DM/$\Omega$b, N$\nu$ = 3, Ng = 8, normal hierarchy, DM darkness. Class B (b-dependent, $\lesssim$ 35\% per 1\% change): m$\mu$/me, m$\tau$/me, $\Delta$m\textsuperscript{2}\textsubscript{31}. Registered observation: Koide's empirical relation Q $\equiv$ (me + m$\mu$ + m$\tau$)/($\surd$me + $\surd$m$\mu$ + $\surd$m$\tau$)\textsuperscript{2} = 2/3, satisfied at the \textasciitilde{}10\textsuperscript{-5} level with no SM explanation, is listed among the b-dependent targets of the lepton sector; no PH derivation is registered.

\subsection*{18.2  SM Gauge Group from Z6}

\noindent Z6 = Z3 $\times$ Z2 = Center(SU(3) $\times$ SU(2)). The SM gauge symmetry SU(3) $\times$ SU(2) $\times$ U(1) is derived, not postulated.

\subsection*{18.3  Framework Comparison: PH vs. SM and GUTs}

\begin{footnotesize}\begin{longtable}{|>{\raggedright\arraybackslash}p{0.5260\textwidth}|>{\raggedright\arraybackslash}p{0.3740\textwidth}|}
\hline
\textbf{Framework} & \textbf{Free inputs} \\ \hline
Standard Model & $\sim$19 \\ \hline
Minimal GUT (SU(5)) & 5--8 \\ \hline
PH (this work) & 2 (GN, Newton's constant; $\Delta$m\textsuperscript{2}\textsubscript{21}, neutrino-sector scale, §11) \\ \hline
PH (P2/P12 closed) & 1 ($\Delta$m\textsuperscript{2}\textsubscript{21} remains; a separate candidate closure, P26-C, is registered at T4 and not yet promoted, §11) \\ \hline
\end{longtable}\end{footnotesize}

\begingroup\centering\footnotesize\itshape Table 10: Input count comparison.\par\endgroup

\subsection*{18.4  What the PH Derives That the SM Assumes}

\noindent Five structural facts that enter the SM as inputs or empirical accidents are theorems here: (i) electron uniqueness (Theorem 6.2) --- the SM offers no reason why exactly one stable charged lepton exists; (ii) the charge--chirality identity (Theorem 6.3) --- where Gell-Mann--Nishijima expresses Q through two further inputs, the PH reads Q as chirality content; (iii) Pauli-triviality of the observable sector (§6.1) --- the inverse of the quark-model $\Delta$++ story; (iv) the statistics of dark matter (Theorem DM-S) --- fixed by parity, not chosen; (v) the Dirac large number --- the electric-to-gravitational force ratio between two electrons, RFG = $\alpha$em $\times$ $\xi$cond\textsuperscript{-1} $\times$ ($\Lambda$meta/me)\textsuperscript{2} = 4.17 $\times$ 10\textsuperscript{42}, and the closed relation ke/GN = $\Lambda$meta\textsuperscript{2}/K($\xi$) (A.2.4): where the SM holds two independent couplings and a 42-order-of-magnitude mystery, both inverse-square laws are bookkeepings of one field --- charge as chirality content (Theorem 6.3), mass as cascade binding --- sharing one condensate scale. Each carries an internal falsification handle through the Z\_N anchor suite ([2]: $\varepsilon$$\psi$(Z4) = 9/25 in the chirality sector; five Q = $-$1 states in the exact Z10 charge census).

\section*{19  Conclusions}

\noindent The Psion Hypothesis derives the observable universe from a single fermionic field $\psi$ with Z6 = Z3 $\times$ Z2 and spin j = 5/2, with one empirical anchor: Newton's constant GN.

\noindent \textbf{• Exact theorems ([T1]): }N$\nu$ = 3, Ng = 8, Ngen = 3; 39 stable DM states; det$'$/det = 1/48; $\sigma$dw/N0 = 1.

\noindent \textbf{• Class A results ([T2]): }Neff = 3.054 (status note, §13.5); $\Omega$DM/$\Omega$b = 5.64 (inference, 5.1\%); the direct route (5.43) is unsupported in v86 (§14.3). $\eta$B: derivation withdrawn; candidate mechanism in §13.7.

\noindent \textbf{• Master invariant ([T2+]): }$\rho$$\Lambda$GN\textsuperscript{2} = 1.121 $\times$ 10\textsuperscript{-123} predicted vs 1.132 $\times$ 10\textsuperscript{-123} observed (0.93\%; $\Lambda$meta cancels and the prediction is GN-independent). Closed form: C$\cdot$K\textsuperscript{2}$\cdot$e\textasciicircum{}($-$S(1+3/$\pi$\textsuperscript{2})) with S = Nc$\cdot$Ne/be; ℏ-linear --- the accelerated expansion is a first-order quantum effect, with co-drift falsifier $\delta$$\rho$$\Lambda$/$\rho$$\Lambda$ = $-$2 $\delta$GN/GN (F17).

\noindent \textbf{• SC input: }1 dimensional anchor (GN) + 1 adimensional datum ($\Delta$m\textsuperscript{2}\textsubscript{21}/me\textsuperscript{2}) $\to$ 15+ predictions. P2/P12 ([T2]): closure $\to$ all adimensional ratios internal; 1 dimensional datum remains (Remark 5.1).

\noindent \textbf{• Hard T1 falsifiers: }any fourth charged lepton, fourth active neutrino, or ninth gluon refutes the framework.

\noindent \textbf{• Structural theorems ([T1]): }unique electron; charge--chirality identity Q = nR $-$ 1; DM statistics by parity (DM-S); Pauli-trivial observable sector; CF-CC flatness with |O|\textsuperscript{2} $\le$ 0.370 (2$\sigma$).

\noindent \textbf{• Gravity-sector consistency: }the condensation channel is a Lorentz scalar (K.9; c\_gw = c\_$\gamma$ by structure) and the Z6 relic-wall problem resolves by infrared tension and exact chiral bias (K.10).

\noindent \textbf{• Critical open step (P23): }the normalisation identity $\delta$$\chi$ = Ind(D$\chi$)/(2j)\textsuperscript{2} --- one equation from T1 closure of $\varepsilon$$\psi$ (its link to $\eta$B is withdrawn, §13.3). \textbf{New open step (P14-B): }a conserved baryon (rung) label, required by nucleon stability and by the candidate baryogenesis mechanism (§13.4, §13.7).

\clearpage
\section*{A  Complete Inventory of PH Results}

\noindent This appendix catalogues all results by category with T-level classification. The framework derives everything from $\psi$ + Z6 + NJL fixed point + GN ([SC]); nothing beyond GN is a free parameter.

\subsection*{A.1  Derived Fundamental Parameters}

\noindent be = 325/1152 $\approx$ 0.2821 (cascade parameter, T1 given the 2PI kernel). Ne = 10 (Z6 theorem, T1). S = 6Ne/be = 212.677 (condensate action, T1). K($\xi$) = 1.16335 (Gel'fand--Yaglom coefficient, T1). f$\pi$ = $\surd$(3/$\pi$) me = 0.4994 MeV (condensate decay constant; confirmed by three independent checks: Yukawa coupling, vortex-action consistency, GY $\surd$$\xi$). $\Lambda$meta = $\surd$(K($\xi$)/GN) exp($-$S/2g\textsuperscript{2}) ([SC] $\to$ [T2], eq. 6).

\subsection*{A.2  Exact Theorems (T1)}

\subsubsection*{\itshape A.2.1  Counting Theorems}

\noindent N$\nu$ = 3 active neutrino species (Z6 singlet counting, [T1]). Ng = 8 gluons (adjoint of the emergent colour group, [T1]). Ne = 10 single-particle modes (Z6 combinatorial theorem, [T1]). Ngen = 3 charged-lepton generations: the fourth would require JR = 7/2 > jmax = 5/2 and is kinematically forbidden within Z6 ([T1]). This is a kinematic upper bound, not a constructive identification of the second and third generations; the latter is addressed separately via the Leptonic Transfer Matrix (P8b, App.\ A.2.5; Appendix L.4.4). SM gauge group SU(3) $\times$ SU(2) $\times$ U(1) is the unique continuous gauge group consistent with the Z6 centre ([T1]).

\subsubsection*{\itshape A.2.2  Algebraic Identities}

\noindent $\sigma$dw/N0 = 1 (domain-wall tension over zero-mode normalisation; one-line algebra, [T1]). GY ratio = 1/48 (Gel'fand--Yaglom domain-wall ratio; three independent verifications, [T1]). $\zeta$$'$4D(0) = $-$3/(2$\pi$) (Kramers--Minakshisundaram spectral zeta, [T1]). Z3 overcounting = 27 (fixed-point Z3 degeneracy, [T1]). $\lambda$ = 9/2 in V($\sigma$) = $\lambda$($\sigma$\textsuperscript{2} $-$ f$\pi$\textsuperscript{2})\textsuperscript{2}: forced by mZ6 = 6f$\pi$ (Theorem W.1: m\textsuperscript{2} = 8$\lambda$f$\pi$\textsuperscript{2}, imposing mZ6 = 6f$\pi$ forces $\lambda$ = 9/2 exactly, [T1]). k $-$ 2j = (2j + 1) $-$ j for j = 5/2, k = 3: unique self-consistency condition of the leptonic sector ([T1]). 1/r\textsuperscript{2} for both Coulomb and Newton (Gauss's theorem in 3D from masslessness, [T1]).

\subsubsection*{\itshape A.2.3  Z6 Gauge Structure Theorems}

\noindent V.1: Qk+3 = Qk + 1; unit Z3 triality; $\Sigma$Qk = 1 structural ([T1]). V.2: the flip 2 $\to$ 3 conserves triality; intermediaries are colour singlets ([T1]). V.3: Eflip must be rational, killing one-meta-boson exchange ([T1]). V.4: the charge interlock is a flux-quantum count, 2j quanta of size 1/N ([T1]). V.5: Z6 = Z3 $\times$ Z2 factorisation removes the BKT caveat, restoring sharp self-duality ([T1]). V.7: the equivalence web --- five reformulations of Lemma I.1' are exactly equivalent; no further reformulation can yield T1 ([T1]).

\subsubsection*{\itshape A.2.4  The $\nu$ Quantum Number}

\noindent $\nu$ $\equiv$ $\lfloor$k/3$\rfloor$ $\oplus$ $\chi$ is a hidden good quantum number of the Z6 structure. Theorem $\nu$.1: $\nu$ = 0 $\Leftrightarrow$ locked subspace, verified exactly on all 12 single-particle states ([T1]). Theorem $\nu$.2: G $\cdot$ Punlocked = G $\cdot$ N$\nu$ with N$\nu$ = $-$(1 $-$ $\hat{Z}$2 $\cdot$ $\chi$) --- the deep-gap projection operator of the I.6 algebraic model is the $\nu$ charge, exact ([T1]).

\subsubsection*{\itshape A.2.5  Leptonic Transfer Matrix (P8b)}

\noindent t12 = $-$8 (J\textsuperscript{+} element, exact, [T1]). $\langle$3|Hdd|1$\rangle$ = t12\textsuperscript{2} = 8 (three-body symmetry cancellation, exact, [T1]). 15 active pairs (direct count). k $-$ 2j = (2j + 1) $-$ j identity. Algebraic cancellation of be in t13 ([T1]). C2 diagonals = \{10.75, 8.75, 6.75\} (exhaustive Q = $-$1 enumeration, [T1]). Mass ratios m$\mu$/me and m$\tau$/me: $\lambda$i/$\lambda$j vary < 0.01\% for $\pm$28\% change in be (insensitivity theorem, repository [2] §I.9.3, [T1]). CF-CC structure (selection rules, existence): [T1]; absolute rate 6.16 $\times$ 10\textsuperscript{-4} $\times$ |O|\textsuperscript{2}, ceiling |O|\textsuperscript{2} $\le$ 0.370 (2$\sigma$) ([T2], P-|O|). Absolute DM stability: four-body Z6 singlets have no kinematically accessible decay channel (exhaustive enumeration, termination proved to N = 48, [T1]).

\subsubsection*{\itshape A.2.6  Gravitational Structure}

\noindent GN = K($\xi$) $\Lambda$meta\textsuperscript{-2} exp($-$S/g\textsuperscript{2}) (T1 at leading order). PC1 ([T2], resolved): hAB is a massive spin-2 composite ($\sim$0.6 GeV), satisfying the Weinberg--Witten theorem by being massive; GN emerges via the Sakharov mechanism [3] as a collective condensate property; there is no massless composite graviton in the PH. The 1/r\textsuperscript{2} law for gravity follows from masslessness of the effective long-range coupling and 3D isotropy ([T1]).

\subsection*{A.3  Structural Results (T2)}

\subsubsection*{\itshape A.3.1  Mass Spectrum (Class B --- b-dependent)}

\noindent me (cascade output) = $\Lambda$meta exp($-$Ne/be) = 0.5089 MeV (0.42\% error; the observed 0.511 MeV is the matrix anchor in P8b). m$\mu$/me: +0.46\%. m$\tau$/me: +0.56\%. Zero free parameters. mp/m$\tau$ = 0.52805 at 0.000\% error (structural, b-independent). GN = 6.653 $\times$ 10\textsuperscript{-39} GeV\textsuperscript{-2} at 0.83\% error against CODATA 2018 in natural units (6.70883 $\times$ 10\textsuperscript{-39} [11]).

\subsubsection*{\itshape A.3.2  Cosmological Results (Class A --- b-independent)}

\noindent $\eta$B: derivation withdrawn in v86 (§13.3); level bifurcation is a candidate mechanism conditional on P14-B (§13.7, [T3]). Neff = 3.054 (< 1$\sigma$ from CMB; CMB-S4 target). $\Omega$DM/$\Omega$b = 5.43 (1.2\% error, direct, [T2+]) / 5.64 (5.1\% error, inference, [T2]); mutual 3.9\% (error propagation, §14.3). $\rho$$\Lambda$ = 2.492 $\times$ 10\textsuperscript{-47} GeV\textsuperscript{4} (0.93\% error). $\rho$$\Lambda$GN\textsuperscript{2} = 1.121 $\times$ 10\textsuperscript{-123} predicted vs 1.132 $\times$ 10\textsuperscript{-123} observed (0.93\%; the prediction is GN-independent).

\subsubsection*{\itshape A.3.3  Force and Coupling Results}

\noindent ke/GN = $\Lambda$meta\textsuperscript{2}/K($\xi$) at 0.000\% error (P22, [T2]). CZ6 = (2j + 1)\textsuperscript{2} $\times$ 2 = 72 (Z6-exact integer; observed ratio $\alpha$meta/$\alpha$em = 71.75 differs by 0.35\% = predicted running of $\alpha$em). RFG = Fem/Fgrav = $\alpha$em $\times$ $\xi$cond\textsuperscript{-1} $\times$ ($\Lambda$meta/me)\textsuperscript{2} = 4.17 $\times$ 10\textsuperscript{42} (three PH factors, no free parameters). The Coulomb and Newton forces share common Z6 origin: both electric charge and gravitational mass are quantum numbers of psion composites built from the same field $\psi$.

\subsubsection*{\itshape A.3.4  Condensate Structure}

\noindent C = Nc/(27$\pi$\textsuperscript{4}) (dark-energy coefficient, [T2+]). f$\pi$ = $\surd$(3/$\pi$) me = 0.4994 MeV (Branch B; confirmed by three independent internal checks: Yukawa coupling, vortex-action consistency, Gel'fand--Yaglom $\xi$). ucond = f$\pi$\textsuperscript{4} = (3/$\pi$)\textsuperscript{2} me\textsuperscript{4} $\approx$ 0.912 me\textsuperscript{4} (Branch B; Conjecture CEC). Self-consistency triangle: GN, $\rho$$\Lambda$, and me are three projections of the single action S = 6Ne/be through the coupling g\textsuperscript{2} = 2$\pi$\textsuperscript{2}/3.

\subsubsection*{\itshape A.3.5  Locking Mechanism ($\nu$-channel, P26)}

\noindent The gap operator G $\cdot$ Punlocked acquires a derived dynamical origin: f1 = 0 (the direct J$\Sigma$ $\to$ $\nu$ route is killed by kinematic identity, exact). f2 = 1 (universal scalar piece). f3 = 3/4 (tensor-induced $\nu$ self-energy, exact rational). f4 = 1/12 (bubble). Mechanism: tensor-induced $\nu$ condensation in two stages --- $\Delta$T (tensor condensate) $\to$ $\Delta$$\nu$ (induced $\nu$ gap) through the 3/4 coefficient. Quantitative open step: ceff = f2 + f3 $\cdot$ w3 + f4 = 11/6 with independently registered multiplicities, vs target 2 required by the operational definition; fork R1/R2 registered (P26, [T2]).

\subsection*{A.4  Ontological Derivations}

\noindent Particles are Z6 tuples (k1, \ldots{}, kN); antiparticles are the same tuples made of antipsions ($N_\psi=-3$, charge $-Q$), whose Z6 labels are the inversions k $\to$ (6 $-$ k) mod 6 (v86 correction; §13.1). Dark matter: absolutely stable Z6-singlet composites (Q = 0, I $\in$ \{0, 3\}; statistics by body-number parity, Theorem DM-S); stability proved by exhaustive enumeration (39 states at N $\in$ \{9, \ldots{}, 27\}, termination theorem, [T1]). Mass scale mDM = 7.205 $\times$ 10\textsuperscript{10} GeV ([T2]). The full composite spectrum contains 76,761 states; the SM is the 8-state shadow under the projections Q $\in$ $\mathbb{Z}$, N $\le$ 3, stable ground states.

\subsection*{A.5  New Forces Predicted}

\noindent \textbf{Force 1 --- Proto-colour force ([T3]): }mediated by the tower of Z6 proto-colour bilinears. The observable member is the proto-colour singlet of the four neutral bilinears ((1,2)$\otimes$(5,4): I = 0, Q = 0, net psion number zero --- the exact analogue of hadrons as colour singlets); individual bilinears carry I $\in$ \{1, 5\} and are confined. The lightest observable member is identified with the X17 singlet (F15): mX17 = $\Lambda$meta exp($-$9/be) = 17.62 MeV (observed 17.01 $\pm$ 0.16 MeV, 3.59\% error, Atomki [8]). Universal coupling across three nuclei follows from Z6 group structure. The tower predicts heavier members at each cascade level N < 9.

\noindent \textbf{Force 2 --- Chiral-Flip force ([T1]): }a new interaction exchanging chirality within the Z6 internal space ($\Delta$$\chi$ = 1, $\Delta$Q = 0, $\Delta$I = 0), absent from SU(3) $\times$ SU(2) $\times$ U(1). Structure [T1]; absolute rate 6.16 $\times$ 10\textsuperscript{-4} $\times$ |O|\textsuperscript{2} with |O|\textsuperscript{2} $\le$ 0.370 (2$\sigma$, MicroBooNE exclusive 1e1p [14]) --- see App. K §K.8.

\noindent \textbf{Force 3 candidate --- resonant gravitational coupling ([T4], P18): }oscillatory $\delta$GN/GN $\sim$ 10\textsuperscript{-10} at fdrive $\in$ [0.38, 0.43] GHz (UHF). Mechanism conjectural ([T4]): requires Z6 composite gauge structure (PC1-sub, open) and resolution of the static-vs-dynamic screening inconsistency (PC2, frozen). Not falsifiable until both close.

\subsection*{A.6  Experimental Anomalies Illuminated}

\noindent CF-CC channel (structural [T1] prediction): RCF = 6.16 $\times$ 10\textsuperscript{-4}; MiniBooNE 4.8$\sigma$ excess was consistent, but MicroBooNE (2022) found no excess --- experimental status disputed (see App. K §K.8). X17 anomalous boson (Atomki, three nuclei): mX17 = 17.62 MeV (F15, [T3], disputed$\dagger$). ANITA anomalous upward-going pulses: candidate carrier $\nu$S = (1, 4, 4), a Z6 composite with I = 3, Q = +1/3, M $\approx$ 7.5 GeV, derived from Z6 selection rules independently of the experimental result.

\subsection*{A.7  Falsifiable Predictions (16 in Table 4)}

\noindent Hard T1 falsifiers (immediate): fourth charged lepton, fourth active neutrino, ninth gluon. Near-term: F2 ($\delta$me/me = $-$$\tfrac{1}{2}$ $\delta$GN/GN); F12 (Neff = 3.054, CMB-S4). Existing data: F15 (X17, Atomki, disputed); F5a/F5b (direct detection: exclusion outside m\textsubscript{DM}, E2; rate at m\textsubscript{DM}, [T4]; XENONnT/LZ). See Table 4 for the complete registry.

\subsection*{A.8  Principal Open Problems}

\noindent P23 ([T2+], critical path): Lemma I.1', $\delta$$\chi$ = Ind(D$\chi$)/(2j)\textsuperscript{2}; closure cascades to $\varepsilon$$\psi$ = 25/169 (T1); the former $\varepsilon$$\psi$ $\to$ $\eta$B link is withdrawn (§13.3). P14-B (open; critical for baryogenesis and nucleon stability): a conserved baryon (rung) label (§13.4, §13.7). Two open components: the unlocked-gap mechanism (P26; mechanism at T2, quantitative closure open) and the winding placement with Ind = 1 in the §2.3 construction. P2/P12 ([T2]): self-consistent fixed-point $\Phi$($\Lambda$) = $\Lambda$, determining $\Lambda$meta from internal consistency; closure makes all adimensional ratios internal; exactly one dimensional datum remains, by the irreducibility established in §5 (a scale-free kernel cannot fix a scale). P1 ([T2]): formal derivation of be = 325/1152 from the 2PI kernel (28\% theoretical gap; proven not to propagate to mass ratios; repository [2] §I.9.3). P8b Passo A.3 ([T2]): formal integration of $\Sigma$2 = be $\times$ 15/8. PC1-sub ([T3]): minimum energy scale to perturb GN via collective modes.

\section*{B  Derivation of the Kink Fluctuation Operator}

\noindent The second derivative of the kink potential V($\sigma$) = $\lambda$($\sigma$\textsuperscript{2} $-$ f$\pi$\textsuperscript{2})\textsuperscript{2} evaluated on the kink solution $\sigma$k(x) = f$\pi$ tanh(mZ6 x/2) gives:

\begingroup\centering V$''$($\sigma$k) = mZ6\textsuperscript{2} $-$ (3/2) mZ6\textsuperscript{2} sech\textsuperscript{2}(mZ6 x/2).\hspace{2em}(20)\par\endgroup

\noindent Substituting y = mZ6 x/2 yields the fluctuation operator:

\begingroup\centering M = (mZ6\textsuperscript{2}/4) [ $-$d\textsuperscript{2}/dy\textsuperscript{2} + 4 $-$ 6 sech\textsuperscript{2}(y) ],   n(n + 1) = 6, n = 2.\hspace{2em}(21)\par\endgroup

\noindent This is the n = 2 Pöschl--Teller Hamiltonian --- a quantum-mechanical potential well exactly solvable in closed form, with two bound states and a reflectionless continuum.

\noindent The two bound states have eigenvalues $\varepsilon$0 = 0 (zero mode / Goldstone) and $\varepsilon$1 = 3mZ6\textsuperscript{2}/4 (shape mode). The zero mode corresponds to translation of the kink along x and must be excluded from the functional determinant; the shape mode contributes the lowest eigenvalue of the fluctuation operator, entering the Gel'fand--Yaglom calculation of Appendix D.

\section*{C  Exact Identities Used in Theorem 3.1}

\noindent The coefficient C = Nc/(27$\pi$\textsuperscript{4}) rests on three exact identities.

\noindent \textbf{Domain-wall tension. }$\sigma$dw = 4f$\pi$\textsuperscript{3}. This follows from integrating the kink solution $\sigma$k(x) = f$\pi$ tanh(mZ6 x/2) through the potential: $\sigma$dw = $\int$[V($\sigma$k)]\textasciicircum{}\{1/2\} d$\sigma$k = 4f$\pi$\textsuperscript{3}.

\noindent \textbf{Zero-mode normalisation. }N0 = (2/3) mZ6 f$\pi$\textsuperscript{2}. The zero-mode wave function $\psi$0(x) $\propto$ sech\textsuperscript{2}(mZ6 x/2) integrates to N0 = (2/3) mZ6 f$\pi$\textsuperscript{2}.

\noindent \textbf{Ratio (exact, T1).}

\begingroup\centering $\sigma$dw/N0 = 4f$\pi$\textsuperscript{3} / [(2/3) $\times$ 6f$\pi$ $\times$ f$\pi$\textsuperscript{2}] = 4f$\pi$\textsuperscript{3}/4f$\pi$\textsuperscript{3} = 1. $\checkmark$ [T1]\hspace{2em}(22)\par\endgroup

\noindent \textbf{The Z3 Faddeev--Popov factor. }The Z3 fixed-point degeneracy contributes an overcounting factor of 27, giving C = Nc/(27$\pi$\textsuperscript{4}) with Nc = 6.

\section*{D  Theorem P16: Renormalisation of the Z6 Wall Determinant}

\noindent This appendix proves Theorem P16 (T1): the explicit renormalisation equation $\sigma$wall\_ren = $\sigma$wall\_bare + $\delta$$\sigma$bulk for the Z6 domain wall. The proof uses the exact heat-kernel trace $\Delta$K(t) = erf(m$\surd$t) + e\textsuperscript{$-$3m\textsuperscript{2}t/4} erf(m$\surd$t/2) (Lemma I), Seeley--DeWitt locality, analytic pole cancellation in Mellin space, and the exact finite part $\Gamma$ren/A = $-$$\surd$3 mZ6\textsuperscript{3}/(96$\pi$).

\subsection*{D.1  The Kirsten--McKane Spectral Zeta Result}

\noindent The Kirsten--McKane formula [6] for the operator L = $-$$\Delta$3 + A(x) in d = 3 transverse dimensions:

\begingroup\centering $\zeta$$'$l(0) = (4$\pi$)\textsuperscript{$-$3/2} $\int$\textsubscript{0}\textasciicircum{}$\infty$ dt t\textsuperscript{$-$5/2} Kdiff(t),\hspace{2em}(23)\par\endgroup

\noindent where Kdiff(t) = e\textsuperscript{$-$$\varepsilon$1 t} and $\varepsilon$1 = 3mZ6\textsuperscript{2}/4. Evaluating the Laplace transform:

\begingroup\centering $\zeta$$'$4D,diff(0) = $-$3/(2$\pi$).   [exact, [T1]]\hspace{2em}(24)\par\endgroup

\subsection*{D.2  The Gel'fand--Yaglom Result}

\noindent The GY method [6] gives the fully renormalised determinant ratio (including bulk subtraction):

\begingroup\centering det$'$(Mk)/det(Mvac) = 1/48.   [T1], exact.\hspace{2em}(25)\par\endgroup

\noindent Three independent verifications: (i) analytic zeta function of exact $\Delta$K(t) --- $\zeta$$'$(0) = ln 12; (ii) lattice spectral pairing (L = 120, N = 32000): 0.08333461 vs 1/12 = 0.08333333; (iii) Jost threshold limit: [$\Pi$(n $-$ j)/(n + j)]/(2n)\textsuperscript{2} = 1/48.

\begingroup\centering $\Gamma$GY = (1/2) log(1/48) = $-$(3/2) log 2 $-$ (1/2) log 4 = $-$1.93574\ldots{}   [T1]\hspace{2em}(26)\par\endgroup

\subsection*{D.3  Theorem P16: Renormalisation Equation ([T2+])}

\noindent \textbf{\textit{Theorem D.1 (Scheme-equivalence of domain-wall determinant, [T2+]). }}\textit{The fluctuation determinant of the kink background admits two renormalisation-scheme representatives:}

\noindent \textbf{• GY scheme: }det$'$(Mk)/det(Mvac) = 1/48 (y-units; 1/12 unit-gap).

\noindent \textbf{• KM scheme: }$\zeta$$'$4D,diff(0) = $-$3/(2$\pi$).

\noindent The mismatch defines the bulk-subtraction cocycle: $\Gamma$bulk$-$ct = $\Gamma$GY $-$ $\Gamma$KM = $-$(3/2) log 2 + 3/(4$\pi$) = $-$0.901889\ldots{} This cocycle is a scheme-transformation artefact absorbed by the domain-wall renormalisation condition. By the standard renormalisation principle, physical observables are scheme-independent: det\_ren\textasciicircum{}$\zeta$ = det\_ren\textasciicircum{}GY. Therefore C = Nc/(27$\pi$\textsuperscript{4}) is scheme-independent.

\noindent \textbf{Consequence ([T2+]): }Theorem P16 closes the renormalisation equation at T1. The prediction $\rho$$\Lambda$ = C $\Lambda$meta\textsuperscript{4} e\textsuperscript{$-$S} = 2.492 $\times$ 10\textsuperscript{-47} GeV\textsuperscript{4} (0.93\% error) is [T2+]: the renormalisation equation is proved (T1), and the C = Nc/(27$\pi$\textsuperscript{4}) assembly retains [T2+] status pending the identification $\mu$4D = 1/$\pi$\textsuperscript{4}.

\noindent The physical coefficient C is constructed from the GY-renormalised wall partition function (Steps 1--4 of Theorem 3.1). The KM result $\zeta$$'$4D,diff(0) = $-$3/(2$\pi$) is not an additional multiplicative factor in C --- it is the heat-kernel representative of the same determinant. The cocycle $\Gamma$bulk$-$ct is absorbed by the domain-wall renormalisation condition and does not appear in physical observables.

\noindent Factor decomposition: C = Nc $\times$ |Z3|\textsuperscript{-3} $\times$ $\mu$4D with Nc = 6, |Z3|\textsuperscript{-3} = 1/27, $\mu$4D = 1/$\pi$\textsuperscript{4} holds independently of regularisation scheme. The 1/$\pi$\textsuperscript{4} origin: the Faddeev--Popov path-integral measure over S\textsuperscript{1}-valued transverse fluctuations contributes a factor $\pi$\textsuperscript{-4} in the 4D functional integral (see P16, [T2]).

\section*{E  Unit Conversion: $\rho$$\Lambda$ in GeV\textsuperscript{4}}

\begingroup\centering $\rho$$\Lambda$\_obs = $\Omega$$\Lambda$ $\cdot$ (3H0\textsuperscript{2}c\textsuperscript{2})/(8$\pi$GN) $\cdot$ ($\hslash$c)\textsuperscript{3}/(1 GeV) = 2.515 $\times$ 10\textsuperscript{-47} GeV\textsuperscript{4}.\hspace{2em}(27)\par\endgroup

\noindent Using $\hslash$c = 0.197327 GeV$\cdot$fm, H0 = 67.36 km/s/Mpc, $\Omega$$\Lambda$ = 0.6847 (Planck 2018 TT,TE,EE+lowE+lensing [1]); equivalently $\rho$\_crit = 3.673 $\times$ 10\textsuperscript{-47} GeV\textsuperscript{4}.

\noindent \textbf{Convention note: }every observational comparator in this paper ($\rho$$\Lambda$,obs; $\rho$$\Lambda$GN\textsuperscript{2}; $\Omega$DM/$\Omega$b = 5.364) is drawn from this single chain --- do not mix Planck chains.

\section*{F  Principal Conclusions and Spectrum Analysis}

\subsection*{F.1  $\alpha$meta = $\pi$/6: Formal Proof}

\noindent \textbf{\textit{Theorem F.1 ($\alpha$meta = $\pi$/Nf, [T1] conditional). }}\textit{The Molien series for Z3 acting on $\mathbb{C}$\textsuperscript{6}: F(x) = (1/3)[1/(1 $-$ x)\textsuperscript{6} + 2/(1 $-$ x\textsuperscript{3})\textsuperscript{2}]. The identity gmeta = 2$\pi$/$\surd$Nf gives $\alpha$meta = g\textsuperscript{2}/(4$\pi$) = $\pi$/Nf = $\pi$/6. $\square$}

\subsection*{F.2  PH Composite Spectrum: Molien Series}

\noindent The generating function F(x) counts Z6-invariant composites at each body number N. The coefficient of x\textasciicircum{}N gives the total number of colour-singlet states.

\begin{footnotesize}\begin{longtable}{|>{\raggedright\arraybackslash}p{0.1000\textwidth}|>{\raggedright\arraybackslash}p{0.1778\textwidth}|>{\raggedright\arraybackslash}p{0.1444\textwidth}|>{\raggedright\arraybackslash}p{0.1556\textwidth}|>{\raggedright\arraybackslash}p{0.1444\textwidth}|>{\raggedright\arraybackslash}p{0.1778\textwidth}|}
\hline
\textbf{N} & \textbf{C(N) total} & \textbf{Cstable} & \textbf{Cunstable} & \textbf{\% stable} & \textbf{Cum. stable} \\ \hline
1 & 2 & 2 & 0 & 100\% & 2 \\ \hline
2 & 7 & 6 & 1 & 86\% & 8 \\ \hline
3 & 20 & 14 & 6 & 70\% & 22 \\ \hline
4 & 42 & 20 & 22 & 48\% & 42 \\ \hline
5 & 84 & 28 & 56 & 33\% & 70 \\ \hline
6 & 156 & 29 & 127 & 19\% & 99 \\ \hline
7 & 252 & 30 & 222 & 12\% & 129 \\ \hline
8 & 462 & 30 & 432 & 6\% & 159 \\ \hline
9 & 670 & 35 & 635 & 5\% & 194 \\ \hline
10 & 1001 & 30 & 971 & 3\% & 224 \\ \hline
12 & 2066 & 30 & 2036 & 1\% & 284 \\ \hline
$\Sigma$ (1--20) & 76761 & 534 & 76227 & 0.7\% & 534 \\ \hline
\end{longtable}\end{footnotesize}

\begingroup\centering\footnotesize\itshape Table 11: PH composite spectrum. Cstable(N) $\in$ \{29, 30, 31\} for all N $\ge$ 10 (Molien saturation, falsifiable prediction F14, [T3]).\par\endgroup

\noindent N = 9 (Cstable = 35) is the final pre-saturation transition level; Cstable(N) $\in$ \{29, 30, 31\} applies strictly for N $\ge$ 10. Any value Cstable(N) $\ge$ 32 falsifies the saturation mechanism (prediction F14, [T3]; testable by running \path{enumerate_z6.py}).

\subsection*{F.3  What the Standard Model Represents}

\noindent The SM is the triple projection of the PH spectrum onto: (1) Q $\in$ $\mathbb{Z}$; (2) N $\le$ 3; (3) stable ground states. This selects 8 states from 76,761.

\subsection*{F.4  PH Boson Spectrum: 36 Bilinears}

\noindent All spin-1 force carriers are bilinears B\textasciicircum{}$\mu$\_kk$'$ = $\bar{\psi}$k $\gamma$\textasciicircum{}$\mu$ $\psi$k$'$ with Z6 gauge charge Ibln = (k$'$ $-$ k) mod 6.

\begin{footnotesize}\begin{longtable}{|>{\raggedright\arraybackslash}p{0.2620\textwidth}|>{\raggedright\arraybackslash}p{0.1253\textwidth}|>{\raggedright\arraybackslash}p{0.1139\textwidth}|>{\raggedright\arraybackslash}p{0.1709\textwidth}|>{\raggedright\arraybackslash}p{0.2278\textwidth}|}
\hline
\textbf{Family} & \textbf{Count} & \textbf{Q} & \textbf{Mass} & \textbf{Force} \\ \hline
Photon (diagonal) & 1 & 0 & massless & EM ($\checkmark$) \\ \hline
Proto-W$\pm$ & 6 & $\pm$1 & MW & Weak ($\checkmark$) \\ \hline
Z\textsuperscript{0} (mixing) & --- & 0 & MZ & Weak ($\checkmark$) \\ \hline
Gluons (Z3) & 8 & 0 & massless & Strong ($\checkmark$) \\ \hline
X17 family & 4 & 0 & 17.6 MeV & Proto-colour \\ \hline
Exotic bilinears & 20 & $\pm$n/9 & cascade & Z6 sector \\ \hline
\end{longtable}\end{footnotesize}

\begingroup\centering\footnotesize\itshape Table 12: PH boson spectrum. Two forces beyond the SM are predicted: the proto-colour force (X17 family, scale $\Lambda$meta) and Z6 Sector Isolation (topological).\par\endgroup

\section*{G  PH Proofs Repository}

\noindent All numerical verifications, enumeration algorithms, and derivation scripts are maintained in psion-proofs (GitHub: \url{https://github.com/karlosmardenmaia/psion-proofs} [2]).

\begin{footnotesize}\begin{longtable}{|>{\raggedright\arraybackslash}p{0.3477\textwidth}|>{\raggedright\arraybackslash}p{0.1432\textwidth}|>{\raggedright\arraybackslash}p{0.4091\textwidth}|}
\hline
\textbf{Module} & \textbf{Section} & \textbf{Content} \\ \hline
\path{proofs/P_dark_matter.py} & §14.1 & 39 stable Z6 dark singlets ([T1]) \\ \hline
\path{proofs/P_dark_energy.py} & §5.2 & C = Nc/(27$\pi$\textsuperscript{4}) functional det ([T1]) \\ \hline
\path{proofs/P_eps_psi.py} & §13.2, App. J & $\varepsilon$$\psi$ = 25/169 chain ([T2]) \\ \hline
\path{proofs/P_counting_theorems.py} & §7--§9 & N$\nu$ = 3, Ng = 8, Ngen = 3 ([T1]) \\ \hline
\path{proofs/P_wall_tensions.py} & §5.2, App. B--C & $\sigma$dw/N0 = 1, det ratio 1/48 ([T1]) \\ \hline
\path{proofs/P_yield_ratio.py} & former §13.3--13.4 & $\eta$B derivation (withdrawn in v86) \\ \hline
\path{p14d_psion_number_audit/} & §13 & P14 audit; level bifurcation ([T3]) \\ \hline
\path{scripts/cascade_anchor.py} & §3 & GN SC input; me cascade output \\ \hline
\path{scripts/dm_prova_real.py} & §14.2 & DM mass verification ([T2]) \\ \hline
\path{scripts/proof_T_CMB_cancels.py} & §14.3 & TCMB-independence ([T1]) \\ \hline
\path{scripts/cascade_level_anchors.py} & §14.2 & Cascade level $\to$ body number ([T1]) \\ \hline
\path{docs/DARK_MATTER_PROOFS.md} & §14 & Full DM proof inventory \\ \hline
\path{docs/PROGRESS_AND_ROADMAP.md} & App. J & P23 derivation log \\ \hline
\end{longtable}\end{footnotesize}

\begingroup\centering\footnotesize\itshape Table 13: Core modules of the psion-proofs repository.\par\endgroup

\noindent All proofs run against a pytest suite ($\ge$ 25 tests, all passing): pip install -r \path{requirements.txt} \&\& pytest tests/. Reproducible on any machine with Python 3.10+.

\noindent \textbf{\textit{Theorem G.1 (Ne = 10, [T1]). }}\textit{Ne = N$\gamma$ + Ng + N$\varphi$ = 1 + 8 + 1 = 10.}

\noindent Proof. N$\gamma$ = 1: unique Z6-singlet spin-1 massless mode (proved in §7). Ng = 8: |Z3|\textsuperscript{2} $-$ 1 = 8 coloured gluonic bilinears (proved T1 in §7). N$\varphi$ = 1: unique Z6-singlet scalar mode. Total: 1 + 8 + 1 = 10. This is a counting theorem of Z6 = Z3 $\times$ Z2; no free parameter. $\square$

\noindent \textbf{Falsifiable prediction F14 ([T3]). }For all N $\ge$ 10, Cstable(N) $\in$ \{29, 30, 31\}. Any value Cstable(N) $\ge$ 32 falsifies the saturation mechanism. Testable by running \path{enumerate_z6.py} for N = 7, 8, 11, 12.

\noindent \textbf{DM stability theorem ([T1]). }39 stable neutral singlets at N $\in$ \{9, 12, 15, 18, 21, 27\} by exhaustive enumeration using two independent algorithms (dynamic programming + brute force); termination proved to N = 48. Code: \path{proofs/P_dark_matter.py}.

\section*{H  The Psion Diagram: Complete Calculus for Composite Classification}

\noindent This appendix makes the PH composite spectrum self-contained. It provides: (H.1) the six psion states and their quantum numbers; (H.2) composite formation rules; (H.3) the five-step Psion Diagram (PD) protocol; (H.4) the complete N = 1--4 spectrum; (H.5) the Molien series and proof of Ne = 10; (H.6) exactly forbidden channels; (H.7) the proto-quark sector; (H.8) the $\alpha$meta = $\pi$/6 formal proof; (H.9) worked examples. No external reference is required to verify these results.

\subsection*{H.1  Composite Formation Rules}

\noindent \textbf{Rule 1 (Z6 colour singlet, [T1]): }a composite is colour-singlet (observable) if and only if I = ($\Sigma$ki) mod 6 $\in$ \{0, 3\}.

\noindent \textbf{Rule 2 (Electric charge): }Q = $\Sigma$Qki. Observable particles have Q $\in$ $\mathbb{Z}$ or Q $\in$ \{$\pm$1/3, $\pm$2/3\} (fractional, confined below $\Lambda$meta).

\noindent \textbf{Rule 3 (Spin): }total spin J follows from coupling individual mJ(ki) values under SU(2)meta. Jmax = N $\times$ 5/2.

\noindent \textbf{Stability criterion: }a composite is stable if it is the lightest state in its (I, Q, J) sector.

\subsection*{H.2  The Five-Step Psion Diagram (PD) Protocol}

\noindent \textbf{Step 1: Assign Z6 triples. }Compute Q = $\Sigma$Qki, I = ($\Sigma$ki) mod 6, chirality. For unknown particles: enumerate all Z6-allowed triples consistent with observed Q and J.

\noindent \textbf{Step 2: Apply conservation laws at every vertex. }(a) Qin = Qout. (b) Iin = Iout mod 6. Sector Isolation Theorem: I = 0 and I = 3 sectors never mix perturbatively.

\noindent \textbf{Step 3: Classify diagram topology. }L = 0 (tree): rate $\sim$ g\textsuperscript{2}meta. L = 1 (one loop): rate $\sim$ g\textsuperscript{4}meta. CF-CC is the unique inter-composite channel exchanging a single constituent between a three-body lepton and a three-body baryon.

\noindent \textbf{Step 4: Determine amplitude type. }O1 (scalar): |$\Delta$mJ| > 1. O2 (spin): |$\Delta$mJ| = 1. Tree-level chiral-flip rate: R $\sim$ NDirac $\times$ g\textsuperscript{2}meta $\times$ (mass ratio)\textsuperscript{2}.

\noindent \textbf{Step 5: Predict observables. }Conservation laws give selection rules. Rate (leading order): BSE overlap integral required for precision. Mass: m $\sim$ $\Lambda$meta $\times$ exp($-$Ncascade/b).

\subsection*{H.3  Complete Composite Spectrum: N = 1 to N = 4}

\noindent \textbf{N = 1. }Only k = 0 (I = 0) and k = 3 (I = 3) are colour-singlets, but both are unstable. No observable N = 1 particles.

\noindent \textbf{N = 2 (bilinears).}

\begin{footnotesize}\begin{longtable}{|>{\raggedright\arraybackslash}p{0.3400\textwidth}|>{\raggedright\arraybackslash}p{0.0800\textwidth}|>{\raggedright\arraybackslash}p{0.0900\textwidth}|>{\raggedright\arraybackslash}p{0.0800\textwidth}|>{\raggedright\arraybackslash}p{0.1700\textwidth}|>{\raggedright\arraybackslash}p{0.1400\textwidth}|}
\hline
\textbf{Representative state} & \textbf{I} & \textbf{Q} & \textbf{J} & \textbf{SM identity} & \textbf{Status} \\ \hline
(k, k$'$), I = 0, Q = 0, J = 1 & 0 & 0 & 1 & Photon $\gamma$ & [T1] \\ \hline
|Z3|\textsuperscript{2} $-$ 1 = 8 coloured states & 0 & $\pm$frac & 0/1 & Gluons (8) & [T1], Ng = 8 \\ \hline
(k, k$'$), J = 2, I = 0, Q = 0 & 0 & 0 & 2 & hAB $\sim$ 0.6 GeV & [T2], not graviton \\ \hline
Others $\notin$ \{0, 3\} & --- & --- & --- & Confined & --- \\ \hline
\end{longtable}\end{footnotesize}

\noindent \textbf{N = 3 (leptons and nucleons).}

\begin{footnotesize}\begin{longtable}{|>{\raggedright\arraybackslash}p{0.2571\textwidth}|>{\raggedright\arraybackslash}p{0.0857\textwidth}|>{\raggedright\arraybackslash}p{0.0964\textwidth}|>{\raggedright\arraybackslash}p{0.1179\textwidth}|>{\raggedright\arraybackslash}p{0.3429\textwidth}|}
\hline
\textbf{Configuration} & \textbf{I} & \textbf{Q} & \textbf{$\Sigma$mJ} & \textbf{Identity} \\ \hline
LLL (0, 1, 2) & 3 & -1 & -9/2 & e\textsuperscript{-} (unique Q = $-$1, Theorem 6.2) \\ \hline
LLR (0, 1, 5) & 0 & 0 & -3/2 & n / $\nu$$\tau$ (dual-sector) \\ \hline
LLR (0, 2, 4) & 0 & 0 & -3/2 & $\nu$$\mu$ / $\Sigma$\textsuperscript{0}-like (dual-sector) \\ \hline
LRR (0, 4, 5) & 3 & +1 & 3/2 & Proton p (charge check: 1/9 + 4/9 + 4/9 = +1) \\ \hline
LLR (1, 2, 3) & 0 & 0 & -3/2 & $\nu$e / $\Lambda$\textsuperscript{0}-like (dual-sector; generation ordering open) \\ \hline
LRR (1, 3, 5) & 3 & +1 & 3/2 & $\Sigma$\textsuperscript{+}-like \\ \hline
LRR (2, 3, 4) & 3 & +1 & 3/2 & $\Sigma$\textsuperscript{+}-like (p-alternative configuration) \\ \hline
RRR (3, 4, 5) & 0 & +2 & 9/2 & Q = +2 exotic (BSM) \\ \hline
\end{longtable}\end{footnotesize}

\noindent \textbf{N $\ge$ 4 (dark matter: neutral singlets).}

\begin{footnotesize}\begin{longtable}{|>{\raggedright\arraybackslash}p{0.4725\textwidth}|>{\raggedright\arraybackslash}p{0.0900\textwidth}|>{\raggedright\arraybackslash}p{0.3375\textwidth}|}
\hline
\textbf{Configuration} & \textbf{I} & \textbf{Status} \\ \hline
(0, 0, 1, 1, 2, 2, 3, 4, 5) & 0 & Ground configuration (closed shell + valence) --- stable, neutral \\ \hline
(0, 0, 0, 0, 0, 0, 1, 1, 4) & 0 & Neutral singlet (census; high shell cost) \\ \hline
(0, 0, 0, 0, 0, 1, 1, 1, 3) & 0 & Neutral singlet (census; high shell cost) \\ \hline
All others & --- & Confined \\ \hline
\end{longtable}\end{footnotesize}

\subsection*{H.4  Exactly Forbidden Channels (T1)}

\begin{footnotesize}\begin{longtable}{|>{\raggedright\arraybackslash}p{0.2352\textwidth}|>{\raggedright\arraybackslash}p{0.3989\textwidth}|>{\raggedright\arraybackslash}p{0.2659\textwidth}|}
\hline
\textbf{Proposed transition} & \textbf{Reason forbidden} & \textbf{Experimental test} \\ \hline
Flavour-changing NC vertex ($\nu$$\mu$ $\to$ $\nu$e at one vertex) & I = 0 and I = 3 sectors never mix at a perturbative vertex; mass-eigenstate oscillation is outside the theorem's scope (P-PMNS) & FCNC suppression \\ \hline
$\nu$$\mu$ $\to$ $\nu$S (exotic) & Single-carrier vertex needs $\Delta$I = 3 with a neutral mediator; the only I = 3 bilinears (proto-W) carry Q = $\pm$1 & No mixing; XRISM null (2025) \\ \hline
4th neutrino & N$\nu$ = 3 is T1 theorem; no 4th Q = 0, I $\in$ \{0, 3\} three-body singlet & Neff < 4; LEP/FCC null \\ \hline
DM $\to$ $\gamma$ direct decay & Net psion number (DM: N $\ge$ 9) vs $\gamma$ (0); no lighter same-N state (ground-state definition); odd N additionally fermion-protected (DM-S) & No keV $\gamma$ line; XRISM 2025 \\ \hline
FCNC (u $\to$ c, d $\to$ s) & Each transition changes I; Sector Isolation forbids single-vertex FCNC & GIM mechanism; FCNC suppressed \\ \hline
Any Qin $\ne$ Qout & Electric charge exactly conserved at every vertex & Universal \\ \hline
\end{longtable}\end{footnotesize}

\begingroup\centering\footnotesize\itshape Table 14: Exactly forbidden transitions and their group-theoretic reasons.\par\endgroup

\subsection*{H.5  Worked Examples: PD Protocol Applied}

\noindent \textbf{Example 1 --- CF-CC reaction (structure, [T1]; experimental status disputed).}

\begingroup\centering $\nu$$\mu$(0, 4, 2) + n(0, 1, 5):  Iin = 12 mod 6 = 0,  Qin = 0.\hspace{2em}(28)\par\endgroup

\begingroup\centering e\textsuperscript{-}(0, 1, 2) + p(0, 4, 5):  Iout = 12 mod 6 = 0 $\checkmark$,  Qout = 0 $\checkmark$.\hspace{2em}(29)\par\endgroup

\noindent RCF = NDirac g\textsuperscript{2}meta (me/m$\mu$)\textsuperscript{2} = 4 $\times$ (2$\pi$\textsuperscript{2}/3) $\times$ (0.511/105.66)\textsuperscript{2} = 6.16 $\times$ 10\textsuperscript{-4}. Structural prediction: RCF = 6.16 $\times$ 10\textsuperscript{-4} $\times$ |O|\textsuperscript{2} (structure T1; overlap open, P-|O|). MicroBooNE bounds |O|\textsuperscript{2} $\le$ 0.370 at 2$\sigma$ (see App. K §K.8).

\noindent \textbf{Example 2 --- Photon identification ([T1]). }Two-body (k, k$'$) with I = 0, Q = 0, J = 1: unique solution is diagonal bilinear B\textasciicircum{}$\mu$\_kk = $\bar{\psi}$k $\gamma$\textasciicircum{}$\mu$ $\psi$k, I = 0, Q = 0, J = 1, massless. $\Rightarrow$ Photon. $\checkmark$

\noindent \textbf{Example 3 --- Dark matter stability ([T2]). }Composite (0, 1, 2, 3): I = 0, Q = 0, J = 0. Decay DM $\to$ $\nu$e + X? X must have J = 1/2, Q = 0, I $\in$ \{0, 3\}. No such lighter composite exists $\Rightarrow$ DM absolutely stable. $\checkmark$

\section*{I  The Proto-Quark: PH Hadronic Constituents and Their Relation to SM Quarks}

\noindent This appendix introduces the term proto-quark for the PH hadronic constituent and develops the formal argument distinguishing it from the SM quark: (I.1) why a new name is warranted; (I.2) three exact structural distinctions; (I.3) what the two frameworks agree on; (I.4) falsifiable predictions; (I.5) responses to expected objections; (I.6) summary table.

\subsection*{I.1  Why a New Name is Warranted}

\noindent The SM quark is defined operationally: charge Q $\in$ \{+2/3, $-$1/3\}, colour charge under SU(3), spin 1/2, and confinement. The PH does not dispute these observations. What it disputes is the ontological status: the SM treats the quark as fundamental. The PH derives the same observable object as a composite of psions.

\noindent \textbf{\textit{Definition I.1 (Proto-quark). }}\textit{A proto-quark is an N = 3 psion composite in a colour-charged Z6 sector (I $\notin$ \{0, 3\}), permanently confined below $\Lambda$meta. Observable SM charges $\pm$2/3, $\pm$1/3 emerge as the low-Q\textsuperscript{2} effective description at Q\textsuperscript{2} $\ll$ $\Lambda$meta\textsuperscript{2}.}

\subsection*{I.2  Three Exact Structural Distinctions}

\noindent \textbf{Confinement: arithmetic vs. dynamical. }In QCD, confinement is dynamical (non-Abelian SU(3) running coupling). In the PH, confinement is arithmetic: I $\notin$ \{0, 3\} ([T1]). No confinement problem exists in the PH.

\noindent \textbf{Fractional charge: axiomatic vs. derived. }In the SM, Q = +2/3 and Q = $-$1/3 are inputs. In the PH, all charges are derived from the Qk formula (§1).

\noindent \textbf{Asymptotic freedom: present vs. absent. }QCD: $\alpha$s(Q\textsuperscript{2}) decreases logarithmically. The PH: no running coupling between proto-quarks below $\Lambda$meta --- proto-quarks are permanently confined for all Q\textsuperscript{2} < $\Lambda$meta\textsuperscript{2}.

\subsection*{I.3  What the Two Frameworks Agree On}

\begin{footnotesize}\begin{longtable}{|>{\raggedright\arraybackslash}p{0.2148\textwidth}|>{\raggedright\arraybackslash}p{0.2352\textwidth}|>{\raggedright\arraybackslash}p{0.2966\textwidth}|>{\raggedright\arraybackslash}p{0.1534\textwidth}|}
\hline
\textbf{Observable} & \textbf{SM QCD} & \textbf{PH Proto-Quark} & \textbf{Agreement} \\ \hline
Proton charge Q(p) = +1 & +2/3 + 2/3 $-$ 1/3 & (0, 4, 5): +1/9 + 4/9 + 4/9 = +1 & Exact \\ \hline
Neutron charge Q(n) = 0 & +2/3 $-$ 1/3 $-$ 1/3 & (0, 1, 5): +1/9 $-$ 5/9 + 4/9 = 0 & Exact \\ \hline
Ng gluons & 8 (SU(3) adjoint) & |Z3|\textsuperscript{2} $-$ 1 = 8 ([T1]) & Exact \\ \hline
Quark confinement & Always confined & I $\notin$ \{0, 3\}: always confined & Exact \\ \hline
mp/m$\tau$ ratio & 0.52805 & exp($-$$\Delta$XD/b) = 0.52805 & 0.000\% ([T2]) \\ \hline
Integer hadron charges & By construction & By Z6 singlet rule & Exact \\ \hline
\end{longtable}\end{footnotesize}

\noindent Every measurement on hadrons to date is consistent with both frameworks. The PH reduces to the SM at Q\textsuperscript{2} $\ll$ $\Lambda$meta\textsuperscript{2}.

\subsection*{I.4  Falsifiable Predictions}

\noindent \textbf{CF-CC inter-composite scattering (structure [T1]; App. K §K.8). }RCF = NDirac g\textsuperscript{2}meta (me/m$\mu$)\textsuperscript{2} $\times$ |O|\textsuperscript{2} = 6.16 $\times$ 10\textsuperscript{-4} |O|\textsuperscript{2}, flat in energy and baseline (Theorem K.8.1); MicroBooNE bounds |O|\textsuperscript{2} $\le$ 0.370 at 2$\sigma$ [14], and SBND's percent-level $\nu$e-like rate measurement tightens the ceiling directly.

\noindent \textbf{Deep inelastic scattering at Q\textsuperscript{2} $\sim$ $\Lambda$meta\textsuperscript{2}. }At Q\textsuperscript{2} $\ll$ $\Lambda$meta\textsuperscript{2}, PH and QCD predictions are identical. At Q\textsuperscript{2} $\sim$ $\Lambda$meta\textsuperscript{2} $\sim$ 10\textsuperscript{24} GeV\textsuperscript{2}, individual proto-quark charges become resolvable and F2(x, Q\textsuperscript{2}) deviates from QCD. Prediction F4 ([T3]).

\subsection*{I.5  Addressing Expected Objections}

\noindent \textbf{Objection 1: ``DIS has directly confirmed quark charges $\pm$2/3, $-$1/3.'' }DIS measures $\Sigma$q e\textsuperscript{2}q q(x) --- not individual charges directly. PH proto-quark charges give the same result at Q\textsuperscript{2} $\ll$ $\Lambda$meta\textsuperscript{2}.

\noindent \textbf{Objection 2: ``Asymptotic freedom is confirmed.'' }Correct --- and this is a PH prediction for Q\textsuperscript{2} $\ll$ $\Lambda$meta\textsuperscript{2}. The PH predicts departure from asymptotic freedom at Q\textsuperscript{2} $\sim$ 10\textsuperscript{24} GeV\textsuperscript{2}.

\noindent \textbf{Objection 3: ``The constituent quark model predicts hadron masses.'' }The PH mass formula gives the proton/tau ratio at 0.000\% error with zero free parameters.

\noindent \textbf{Objection 4: ``Why introduce new terminology?'' }The SM quark is defined as fundamental. Calling the PH constituent a `quark' would contradict the framework. `Proto-quark' signals the correct relationship: the quark is the effective description; the proto-quark is its sub-structure.

\subsection*{I.6  Summary Table}

\begin{footnotesize}\begin{longtable}{|>{\raggedright\arraybackslash}p{0.2148\textwidth}|>{\raggedright\arraybackslash}p{0.2045\textwidth}|>{\raggedright\arraybackslash}p{0.2761\textwidth}|>{\raggedright\arraybackslash}p{0.2045\textwidth}|}
\hline
\textbf{Property} & \textbf{SM Quark} & \textbf{PH Proto-Quark} & \textbf{Observationally distinct?} \\ \hline
Fractional charge & Input axiom & Derived from Qk formula & At Q\textsuperscript{2} $\sim$ $\Lambda$meta\textsuperscript{2} only \\ \hline
Individual charges & +2/3, $-$1/3 & +1/9, +4/9, +5/9 & At Q\textsuperscript{2} $\sim$ $\Lambda$meta\textsuperscript{2} only \\ \hline
Charge sums & Q(p) = +1, Q(n) = 0 & Q(p) = +1, Q(n) = 0 & Identical \\ \hline
Asymptotic freedom & Yes (running $\alpha$s) & No (condensate at $\Lambda$meta) & At Q\textsuperscript{2} $\sim$ 10\textsuperscript{24} GeV\textsuperscript{2} \\ \hline
CF-CC scattering & Not in QCD & RCF = 6.16 $\times$ 10\textsuperscript{-4} & Testable now \\ \hline
Mass formula & Fitted parameters & exp(A|mJ\_R|\textasciicircum{}p/b); p, A derived & Testable (mass ratios) \\ \hline
\end{longtable}\end{footnotesize}

\noindent \textbf{Epistemic summary: }The proto-quark is not a falsification of the quark --- it is its sub-structure. The SM quark is the correct effective description at Q\textsuperscript{2} $\ll$ $\Lambda$meta\textsuperscript{2}. Two predictions are testable at current energies: the CF-CC rate (MicroBooNE) and the zero-parameter mass formula.

\clearpage
\section*{J  The Chirality Asymmetry $\varepsilon$$\psi$ = 25/169: Results}

\noindent This appendix contains the complete result set behind the chirality asymmetry parameter. Full derivations, verification scripts, and the problem registry are maintained in the companion repository [2].

\subsection*{J.1  Physical Mechanism}

\noindent The contact interaction $-$($\bar{\psi}$J\textsuperscript{a}$\psi$)($\bar{\psi}$$\Sigma$\textsuperscript{a}$\psi$) is odd under Z2 $\subset$ Z6: three-body LLL composites bind more efficiently than RRR composites at the Z6 phase transition. \textbf{v86:} LLL and RRR are not matter and antimatter (antimatter is made of antipsions, §13.1), and $\varepsilon$$\psi$ is CP-even; it does not by itself produce a baryon excess (§13.3).

\subsection*{J.2  The Chiral Asymmetry Formula}

\begingroup\centering $\varepsilon$$\psi$ = 4$\delta$$\chi$/(1 + $\delta$$\chi$)\textsuperscript{2}.\hspace{2em}(30)\par\endgroup

\noindent For $\delta$$\chi$ = 1/25: $\varepsilon$$\psi$ = (4/25)/(26/25)\textsuperscript{2} = 100/676 = 25/169 $\approx$ 0.1479.

\subsection*{J.3  Structural Origin: $\delta$$\chi$ = 1/(2j)\textsuperscript{2} ([T2])}

\noindent The binding asymmetry between the L-sector (k = 0, 1, 2) and R-sector (k = 3, 4, 5) is a second-order effect of the Ja$\Sigma$a coupling; second-order corrections in a representation of spectral width 2j scale as 1/(2j)\textsuperscript{2}. Hence:

\begingroup\centering $\delta$$\chi$ = 1/(2j)\textsuperscript{2} = 1/25,\hspace{2em}(31)\par\endgroup

\begingroup\centering $\delta$$\chi$ = 1/25  $\Rightarrow$  $\varepsilon$$\psi$ = 25/169 = 0.14793\ldots{}.\hspace{2em}(32)\par\endgroup

\noindent The value (2j)\textsuperscript{2} = 25 is the squared spectral width that fixes the six-state Z6 assignment (j = 5/2 is the minimal half-integer spin with exactly six states). $\varepsilon$$\psi$ is an exact rational of the spin-5/2 structure, not a fitted number; its first-principles derivation is the open problem P23 ([T2+]).

\subsection*{J.4  Former Baryogenesis Chain (withdrawn in v86)}

\noindent \textbf{Status (v86). }The chain below is withdrawn (§13.3): $\xi^*$ in eq.~(37) was computed from the target $\eta_B$, so eq.~(38) is an identity, not a verification. It is retained for the record. Former chain:

\begingroup\centering $\eta$B = $\varepsilon$$\psi$ $\times$ $\varepsilon$conf $\times$ (n/n$\gamma$) $\times$ ($\xi$*)\textsuperscript{-3},\hspace{2em}(33)\par\endgroup

\begingroup\centering $\varepsilon$$\psi$ = 25/169 = 0.14793  (exact rational fraction),\hspace{2em}(34)\par\endgroup

\begingroup\centering $\varepsilon$conf = N3/(N2 + N3 + N4) = 4/9,\hspace{2em}(35)\par\endgroup

\begingroup\centering n/n$\gamma$ = 0.749,\hspace{2em}(36)\par\endgroup

\begingroup\centering $\xi$* = 434.09  (computed from the target $\eta$B; see status note),\hspace{2em}(37)\par\endgroup

\begingroup\centering $\eta$B = 0.147929 $\times$ 0.444444 $\times$ 0.749/434.09\textsuperscript{3} = 6.02 $\times$ 10\textsuperscript{-10}.\hspace{2em}(38)\par\endgroup

\noindent \textbf{Note on $\varepsilon$conf. }$\varepsilon$conf = 4/9 counts four lepton-number three-body composites (e\textsuperscript{-}, $\mu$\textsuperscript{-}, $\tau$\textsuperscript{-}, $\nu$e) out of nine three-body families. Verification: 4/9 $\Rightarrow$ $\xi$* = 434.09 $\checkmark$; 4/13 $\Rightarrow$ $\xi$* = 384 \ding{55}. Declared assumptions ([T2]): $\varepsilon$conf = 4/9 and n/n$\gamma$ = 0.749 are thermal-democracy counts --- equal formation weight across the three-body families at the T $\sim$ $\Lambda$meta transition --- and the factorization of eq. (33) (asymmetry $\times$ counting $\times$ density, no cross terms) is likewise declared. The direct $\Omega$DM/$\Omega$b route (elasticity +8.27, §14.3) doubles as an empirical bound on democracy violations at the 1--2\% level. Full derivation: repository [2], P14.

\subsection*{J.5  Index-Normalisation Conjecture}

\noindent \textbf{\textit{Conjecture J.1 (P23 reformulated, [T2] candidate).}}

\begingroup\centering $\delta$$\chi$ = Ind(D$\chi$)/(2j)\textsuperscript{2} = 1/25.\hspace{2em}(39)\par\endgroup

\noindent The Dirac operator in the Z6 kink background,

\begingroup\centering D$\chi$ = i$\gamma$$\mu$$\partial$$\mu$ + g $\Phi$(x) Ja$\Sigma$a,   $\Phi$(x) = f$\pi$ tanh(mZ6 x/2),\hspace{2em}(40)\par\endgroup

\noindent carries one chiral zero mode by the Jackiw--Rebbi mechanism [7]:

\begingroup\centering Ind(D$\chi$) = nL $-$ nR = 1,   [T1]\hspace{2em}(41)\par\endgroup

\begingroup\centering W = $\Delta$mJ = 5/2 $-$ ($-$5/2) = 5 = 2j.   [T1]\hspace{2em}(42)\par\endgroup

\begingroup\centering $\delta$$\chi$ = Ind(D$\chi$)/(2j)\textsuperscript{2} = 1/25  [conditional on Conjecture J.1].\hspace{2em}(43)\par\endgroup

\noindent Ind = 1 counts the topological unit of chirality; (2j)\textsuperscript{2} is the spectral volume. Once the identity is proved, $\varepsilon$$\psi$ = 25/169 becomes a T1 theorem (its link to $\eta$B is withdrawn, §13.3).

\subsection*{J.6  The Chiral Mixing Angle: Exact Theorem Suite}

\subsubsection*{\itshape J.6.1  Theorem I.1 ([T1]) --- Bare J$\cdot$$\Sigma$ diagonal in the locked subspace}

\noindent In the chirality-locked subspace (|$-$j\ldots{}$-$1/2, L$\rangle$, |+1/2\ldots{}+j, R$\rangle$, $\sigma$ = sign m), the bare coupling Jz$\Sigma$z + $\lambda$(J\textsuperscript{+}$\Sigma$\textsuperscript{-} + J\textsuperscript{-}$\Sigma$\textsuperscript{+}) is purely diagonal, equal to |m|: J\textsuperscript{+}$\Sigma$\textsuperscript{-} and J\textsuperscript{-}$\Sigma$\textsuperscript{+} cannot connect two locked states. The coupling producing the asymmetry must therefore come from the condensate.

\subsubsection*{\itshape J.6.2  Theorem I.2 ([T1]) --- Winding condensate activates a single bond}

\noindent The Z6 winding condensate (winding fraction $\varepsilon$ = 1/6 per kink) generates J\textsuperscript{+}$\Sigma$\textsuperscript{+} + J\textsuperscript{-}$\Sigma$\textsuperscript{-}, activating exactly one bond |$-$1/2, L$\rangle$ $\leftrightarrow$ |+1/2, R$\rangle$ with the exact SU(2) element:

\begingroup\centering $\langle$+1/2|J\textsuperscript{+}|$-$1/2$\rangle$ = j + 1/2 = 3,   t = $\varepsilon$(j + 1/2) = 1/2.\hspace{2em}(44)\par\endgroup

\subsubsection*{\itshape J.6.3  Theorem I.3 ([T1]) --- Universality of t = 1/2 and the zero mode}

\begingroup\centering t = $\varepsilon$(j + 1/2) = (j + 1/2)/(2j + 1) = 1/2  $\forall$j;   H = diag(|m|) + t(|L$\rangle$$\langle$R| + |R$\rangle$$\langle$L|) has an exact E = 0 mode.\hspace{2em}(45)\par\endgroup

\noindent Verified for j = 1/2 through 9/2: the winding normalisation $\varepsilon$ = 1/Nc and the Jackiw--Rebbi index theorem are equivalent statements in this algebraic model.

\subsubsection*{\itshape J.6.4  The Z6 Charge Interlock ([T1])}

\noindent The zero mode is |0$\rangle$ = (|k = 2, L$\rangle$ $-$ |k = 3, R$\rangle$)/$\surd$2. With the Z6 charges of §1:

\begingroup\centering Q2 = $-$5/9,  Q3 = +10/9,  $\Delta$Q = 5/3,\hspace{2em}(46)\par\endgroup

\begingroup\centering $\langle$0|Q$\Sigma$\textsuperscript{z}|0$\rangle$ = $\Delta$Q/2 = 5/6  [offset-independent],\hspace{2em}(47)\par\endgroup

\begingroup\centering Nc$\langle$0|Q$\Sigma$\textsuperscript{z}|0$\rangle$ = 6 $\times$ 5/6 = 5 = 2j $\cdot$ Ind(D$\chi$).   $\star$ EXACT INTERLOCK $\star$\hspace{2em}(48)\par\endgroup

\noindent \textbf{Universality ([T1]). }For every ZN, the proto-colour charge coefficient obeys aN = [(N $-$ 1)/N]/(1 $-$ cos(2$\pi$/N)), giving $\Delta$Q = (2N $-$ 1)/N and hence Nc$\langle$Q$\Sigma$\textsuperscript{z}$\rangle$0 = 2N $-$ 1 = 2j identically. The Z3 coefficient 4/9 is derived, not input. Corollary ($\Sigma$Qk = 1, [T1]): Qk = q(k mod 3) + $\lfloor$k/3$\rfloor$ with q = (1/9, $-$5/9, $-$5/9), so $\Sigma$Qk = 2$\Sigma$q + 3 = 1 --- the unit sum is the integer Z2 chirality shift, a structural corollary of the Z3 $\times$ Z2 factorisation; the §J.6.4 charge difference is offset-independent.

\subsubsection*{\itshape J.6.5  Chiral Mixing Angle ([T1])}

\noindent The formula $\varepsilon$$\psi$ = 4$\delta$$\chi$/(1 + $\delta$$\chi$)\textsuperscript{2} is the double-angle identity under $\delta$$\chi$ = tan\textsuperscript{2}$\theta$$\chi$:

\begingroup\centering $\varepsilon$$\psi$ = sin\textsuperscript{2}(2$\theta$$\chi$);   tan $\theta$$\chi$ = 1/5  $\Rightarrow$  sin 2$\theta$$\chi$ = 5/13, cos 2$\theta$$\chi$ = 12/13,\hspace{2em}(49)\par\endgroup

\begingroup\centering $\varepsilon$$\psi$ = (5/13)\textsuperscript{2} = 25/169.   [Pythagorean triple 5\textsuperscript{2} + 12\textsuperscript{2} = 13\textsuperscript{2}]\hspace{2em}(50)\par\endgroup

\subsubsection*{\itshape J.6.6  Lemma I.1$'$ --- the single open step ([T2+])}

\begingroup\centering tan $\theta$$\chi$ = Nc$\langle$0|Q$\Sigma$\textsuperscript{z}|0$\rangle$/(2j)\textsuperscript{2} = 5/25 = 1/5.   [[T2+], open]\hspace{2em}(51)\par\endgroup

\noindent By spectral reduction (using only the interlock and t = 1/2, both T1), Lemma I.1$'$ is algebraically equivalent to tan $\theta$$\chi$ = t/j: the charge structure drops out; what remains is a pure spectral statement.

\subsubsection*{\itshape J.6.7  Jackiw--Rebbi Anchor and Structural Protection ([T1])}

\noindent At j = 1/2, Lemma I.1$'$ gives tan $\theta$$\chi$ = 1 $\Rightarrow$ $\varepsilon$$\psi$ = 1 --- exactly the Jackiw--Rebbi result (the kink zero mode is 100\% chirally polarised) [7]. PASS. Edge protection: at j = 1/2 there is no leakage channel, so $\varepsilon$$\psi$(Z2) = 1 for all condensate gaps G and all $\lambda$. Z4 invariant-subspace protection: the sector S = \{($-$1/2, L), (+1/2, R), ($-$3/2, R), (+3/2, L)\} satisfies H($\lambda$, G)$\cdot$S $\subseteq$ S for all $\lambda$, G, forcing $\varepsilon$$\psi$(Z4) = 9/25 exactly for any eigenstate in S. Finite-gap correction law: $\varepsilon$$\psi$(G) = 25/169 $-$ $\kappa$6/G\textsuperscript{2} + O(1/G\textsuperscript{3}) for Z6, with exact leading term $-$480/2197 from first-order perturbation theory.

\subsubsection*{\itshape J.6.8  General Formula}

\begin{footnotesize}\begin{longtable}{|>{\raggedright\arraybackslash}p{0.1011\textwidth}|>{\raggedright\arraybackslash}p{0.1618\textwidth}|>{\raggedright\arraybackslash}p{0.1315\textwidth}|>{\raggedright\arraybackslash}p{0.1315\textwidth}|>{\raggedright\arraybackslash}p{0.2427\textwidth}|>{\raggedright\arraybackslash}p{0.1315\textwidth}|}
\hline
\textbf{j} & \textbf{Group} & \textbf{tan $\theta$$\chi$} & \textbf{$\varepsilon$$\psi$} & \textbf{Pythagorean triple} & \textbf{sec(2$\theta$$\chi$)} \\ \hline
1/2 & Z2 (JR) & 1 & 1 & --- & --- \\ \hline
3/2 & Z4 & 1/3 & 9/25 & (3, 4, 5) & 5/4 \\ \hline
5/2 & Z6 & 1/5 & 25/169 & (5, 12, 13) & 13/12 \\ \hline
7/2 & Z8 & 1/7 & 49/625 & (7, 24, 25) & 25/24 \\ \hline
\end{longtable}\end{footnotesize}

\noindent Pattern: tan $\theta$$\chi$ = 1/(2j); triple (2j, 4j\textsuperscript{2} $-$ 1, 4j\textsuperscript{2} + 1). Internal falsification target: a full Z4 dynamical calculation must give $\varepsilon$$\psi$ = 9/25 (protected exactly by the invariant-subspace theorem).

\subsection*{J.7  Z6 Gauge Structure Theorems ([T1])}

\noindent \textbf{V.1 (Z3 $\times$ Z2 charge decomposition). }Qk = q(k mod 3) + $\lfloor$k/3$\rfloor$; the Z2 factor acts as Q $\to$ Q + 1; every psion carries unit Z3 triality, 9Qk $\equiv$ 1 (mod 3).

\noindent \textbf{V.2 (Selection rules). }All Q = $-$1 lepton composites and their flip intermediates are colour singlets; the flip 2 $\to$ 3 conserves triality. The flip cost comes from the Z6 flux/contact sector, not from confinement.

\noindent \textbf{V.3 (Rationality). }Eflip = 2j$\cdot$t = 5/2 must be rational; one-meta-boson exchange (carrying $\alpha$meta = $\pi$/6) is excluded as the mechanism.

\noindent \textbf{V.4 (Flux-quantum count). }The flip transports $\Delta$Q = 5/3 = 2j quanta of size 1/3; for every Z2N, quanta = Nc$\Delta$Q/2 = 2N $-$ 1 = 2j identically --- the charge interlock is a flux count.

\noindent \textbf{V.5 (Factorisation). }Z6 $\cong$ Z3 $\times$ Z2 (CRT) and the charge table factorises exactly; the relevant statistical system is a product of a Z2 and a Z3 system, each with a unique Kramers--Wannier self-dual transition (no BKT window).

\noindent \textbf{\textit{Theorem J.3 (V.7: equivalence web, [T1]). }}\textit{The following are exactly equivalent, verified on all four anchor groups: (1) Lemma I.1$'$; (2) tan $\theta$$\chi$ = t/j; (3) Eflip = 2j$\cdot$t = j; (4) the flip intermediate sits at the locked spectral edge; (5) flux costs equal (t, N $-$ 1) = (half bond splitting, bond$\to$edge distance). Conditionally on L1 ([T2]), all are further equivalent to (6): the colour factor sits at its Wegner self-dual point. No algebraic reformulation can close the set to T1; a dynamical input is required.}

\subsection*{J.8  Self-Duality Results ([T1])}

\noindent \textbf{SD-1. }The winding coupling $\varepsilon$ = 1/Nc equals the Villain 4D self-dual coupling $\alpha$sd(ZN) = 1/(2N) exactly, on all four anchor groups: the Z3 colour factor of the condensate sits exactly at the Villain self-dual point.

\noindent \textbf{SD-2. }ZN is Pontryagin self-dual; invariance under q $\mapsto$ N $-$ q forces E(flux q) = f(min(q, N $-$ q)) --- the conjugacy min-rule is derived by group theory, not chosen.

\noindent \textbf{SD-3. }With $\Phi$($\Lambda$) = $\Lambda$ (P2/P12), the Z6-uniform condensate $\langle$$\bar{\psi}$k$\psi$k$\rangle$ = $\sigma$/Nc factorises the effective action exactly: Seff = Seff(Z3) $\times$ Seff(Z2).

\noindent \textbf{SD-4. }The single matter coupling t = 1/2 places both factors simultaneously at their self-dual couplings: $\alpha$Z3 = t/3 = 1/6 = $\alpha$sd(Z3), $\alpha$Z2 = t/2 = 1/4 = $\alpha$sd(Z2), with no free parameter. SD-4 fixes the matter coupling; Eflip runs through the gauge stiffness $\kappa$ --- a distinct coupling --- so self-duality alone does not close P23.

\subsection*{J.9  The $\nu$ Quantum Number and the Locking Mechanism}

\noindent \textbf{\textit{Theorem J.4 ($\nu$.1, [T1]). }}\textit{$\nu$ $\equiv$ $\lfloor$k/3$\rfloor$ $\oplus$ $\chi$ is a hidden good quantum number: $\nu$ = 0 is exactly the locked subspace, verified state by state on all twelve.}

\noindent \textbf{\textit{Theorem J.5 ($\nu$.2, [T1]). }}\textit{Punlocked = $-$(1 $-$ $\hat{Z}$2$\cdot$$\chi$): the deep-gap projector of the algebraic model is the $\nu$ charge, G$\cdot$Punlocked = G$\cdot$N$\nu$, exactly.}

\noindent \textbf{Locking mechanism ([T2], exact coefficients). }Projections of the §2 interactions onto the $\nu$-channel: f1 = 0 (direct route killed by kinematic identity), f2 = 1 (scalar), f3 = 3/4 (tensor-induced, exact), f4 = 1/12 (bubble). Locking is tensor-induced $\nu$ condensation: the tensor condensate $\Delta$T forms in the J$\Sigma$ channel, then breeds the $\nu$ gap through the exact 3/4 coefficient --- a self-reinforcing bootstrap in which $\nu$-conservation is emergent in the deep-gap limit. Quantitative closure is open: the registered multiplicities give ceff = f2 + 2f3 + f4 = 11/6 against the 2G convention of the operational definition; resolving the discrepancy requires the self-consistent non-commuting two-stage solve (P26). Falsifiable corollary: if $\nu$ condenses, the $\nu$ = 1 sector is a physical multiplet at mass $\sim$2G, to be confronted with the hAB resonance ($\sim$0.6 GeV).

\subsection*{J.10  The $\kappa$/G Barrier: Exact Form and Interpretation}

\noindent Every analytic route from the §2 action to the k-ladder gauge stiffness yields the same exact ratio ([T2]): electromagnetic duality with $\beta$mag = Nc (which follows from g\textsuperscript{2} = 2$\pi$\textsuperscript{2}/3 via $\beta$elec $\times$ $\beta$mag = Nc\textsuperscript{2}/4$\pi$\textsuperscript{2}, itself a new T2 result) gives $\kappa$/G = Nc\textsuperscript{2}/(4$\pi$\textsuperscript{2}) = 9/$\pi$\textsuperscript{2} $\approx$ 0.912 --- an 8.8\% deficit against the pass criterion $\kappa$/G = 1, fixed by Nc = 6 versus 2$\pi$. Ward identities do not determine $\kappa$ (charge is incompressible in the deep gap); $\kappa$ is a construction datum of the emergent-gauge map, not a symmetry consequence.

\subsubsection*{\itshape J.10.1  The sinc\textsuperscript{2} Identity ([T2])}

\begingroup\centering $\kappa$/G = sinc\textsuperscript{2}($\pi$/Nc),   sinc(x) = sin(x)/x;   sinc\textsuperscript{2}($\pi$/6) = (1/2)\textsuperscript{2}/($\pi$/6)\textsuperscript{2} = 9/$\pi$\textsuperscript{2}.\hspace{2em}(52)\par\endgroup

\noindent sinc is the lattice-to-continuum energy ratio of the fundamental ZNc mode: the deficit is the Nyquist sampling error of the massless psion on a six-point lattice; $\kappa$/G $\to$ 1 as Nc $\to$ $\infty$. Geometric form: $\surd$($\kappa$/G) = chord/arc of a Z6 circular sector (half-angle 30$^{\circ}$): the discrete jump between Z6 sectors versus continuous propagation on the phase circle. The gap G is itself lattice-generated (the massless psion has Gcont = 0), so both $\kappa$ and G are intrinsically lattice objects and the sinc\textsuperscript{2} is a genuine physical ratio, not a removable artefact.

\subsubsection*{\itshape J.10.2  The Psionic Fine-Structure Constant ([T2])}

\begingroup\centering $\alpha$$\psi$ = 1 $-$ sinc\textsuperscript{2}($\pi$/Nc) = 1 $-$ 9/$\pi$\textsuperscript{2} $\approx$ 0.0881,   1/$\alpha$$\psi$ $\approx$ 11.35.\hspace{2em}(53)\par\endgroup

\noindent The massless psion's fine structure separates into three constants: $\alpha$meta = $\pi$/6 (dynamical coupling), $\varepsilon$ = 1/6 (winding unit), $\alpha$$\psi$ $\approx$ 0.0881 (geometric sampling deficit) --- a separation that is itself the signature of masslessness (a massive field collapses all readings to one constant, as the electron does to 1/137). Numerical observation, not a derivation: $\alpha$$\psi$ $\approx$ 12$\alpha$ to 0.6\%.

\subsubsection*{\itshape J.10.3  The Fundamental Identity ([T2])}

\begingroup\centering sinc\textsuperscript{2}($\pi$/Nc) $\times$ g\textsuperscript{2} $\times$ gL $\times$ t/(2j + 1) = 1.\hspace{2em}(54)\par\endgroup

\noindent Exact: [Nc\textsuperscript{2}/4$\pi$\textsuperscript{2}] $\times$ [2$\pi$\textsuperscript{2}/3] $\times$ [2] $\times$ [(1/2)/6] = 1, every factor independently [T1]. The identity states that the compensating factor 1/sinc\textsuperscript{2} = $\pi$\textsuperscript{2}/9 equals g\textsuperscript{2} $\times$ gL $\times$ t/(2j + 1); it establishes consistency of $\kappa$/G = 1 with the derived parameter set, but not that the effective stiffness $\kappa$eff takes this form --- that derivation is the §2.3 $\to$ k-ladder map, the open content of P23.

\subsubsection*{\itshape J.10.4  Scale-Hierarchy Link ([T2])}

\noindent Nc/(2$\pi$) = sinc($\pi$/Nc) exactly. The observed hierarchy ln($\Lambda$meta/mZ6) = 33.68 lies between S/(2$\pi$) = 33.85 (continuum, 0.5\%) and S/Nc = 35.45 (discrete, 5.2\%); their ratio is exactly the P23 gap factor sinc($\pi$/Nc) = 3/$\pi$. The physical scale separation is set by the continuum psion dispersion while the $\kappa$/G deficit reflects the discrete Z6 lattice --- two faces of one tension. Additional exact identity: S $\times$ be = 6Ne = 60 = 10Nc.

\subsection*{J.11  The f$\pi$ Branch Question}

\noindent Two exact, mutually exclusive forms for the condensate decay constant coexist in the framework: Branch A, f$\pi$ = me (from the App. B identity $\sigma$dw/N0 = 1 with mZ6 = 6me, P24); Branch B, f$\pi$ = $\surd$(3/$\pi$) me (explicit in §3.2, where the normalisation $\lambda$ = 9/2 is forced by mZ6 = 6f$\pi$ --- Theorem W.1, [T1] --- and 3/$\pi$ = Nc/(2$\pi$) gives Branch B the shape of a derived relation). The branches differ by 2.28\%; existing numerical routes sit between them. Decisive question: the provenance of the §3.2 matching. Quarantine ([T1] in scope): the baryogenesis pipeline (eq. 33) contains no f$\pi$ or $\sigma$dw dependence, so the fork cannot propagate to $\eta$B, Neff, or the dark-matter chain; it is contained to \{P24, Conjecture CEC, §3.2\}.

\subsection*{J.12  Axiom I.1 --- The $\kappa$ = G Postulate}

\noindent \textbf{Postulate P-$\kappa$ (Axiom I.1). }$\kappa$ = G: in the deep-gap limit of the k-ladder construction, the electric stiffness of the emergent Z6 gauge sector equals the matter gap.

\noindent \textbf{Status: }postulate, not theorem. Every algebraic route is blocked by Theorem V.7; every analytic route from §2 yields $\kappa$/G = sinc\textsuperscript{2}($\pi$/Nc). The axiom is equivalent to asserting $\alpha$$\psi$ = 0 at the physical point --- that the emergent-gauge construction supplies the 1/sinc\textsuperscript{2} enhancement, whose required form is exactly the fundamental identity (eq. 54).

\noindent \textbf{What it delivers: }$\kappa$ = G $\Rightarrow$ Eflip = j $\Rightarrow$ Lemma I.1$'$ $\Rightarrow$ $\varepsilon$$\psi$ = 25/169 (T1 conditional); the further step to $\eta$B is withdrawn (§13.3).

\noindent \textbf{Falsifiability: }any explicit construction of the §2.3 $\to$ k-ladder emergent-gauge map yielding $\kappa$ $\ne$ G refutes it; conversely, a construction delivering the eq. 54 factor set proves it.

\section*{K  Gravitational Sector: Condensate Response}

\noindent If GN is the tensorial susceptibility of the Z6 condensate, the vacuum is a medium with characteristic frequencies, resonances, and penetration depths.

\subsection*{K.1  Convention CN-1 (Planck Convention)}

\noindent The susceptibility relation GN = K($\xi$) $\Lambda$meta\textsuperscript{-2} exp($-$S/g\textsuperscript{2}) is written in the standard (non-reduced) convention, GN = MPl\textsuperscript{-2} with GN = 6.70883 $\times$ 10\textsuperscript{-39} GeV\textsuperscript{-2} (CODATA 2018 in natural units) and MPl = 1.22089 $\times$ 10\textsuperscript{19} GeV. Any appearance of 8$\pi$GN must cite this note. Self-consistency closes exactly:

\begingroup\centering S/g\textsuperscript{2} = 2 ln(MPl/$\Lambda$meta) + ln K($\xi$) = 32.323  (identically 6Ne/(be g\textsuperscript{2})).\hspace{2em}(55)\par\endgroup

\subsection*{K.2  Lemma J.1 --- Kinematic Gravitational Josephson Frequency ([T2+])}

\noindent \textbf{Premise P1 (open): }the order parameter relevant to GN is bilinear in the lightest composite, pair rest energy 2mec\textsuperscript{2}. 

\noindent \textbf{Premise P2 ([T1], established): }in the static weak field, a state of rest energy E accumulates phase at $\dot{\theta}$ = $-$E(1 + $\varphi$/c\textsuperscript{2})/$\hslash$ (gravitational redshift of the Compton clock; Chou et al., Science 329, 1630, 2010).

\begingroup\centering fJ = (2mec\textsuperscript{2}/h) $\cdot$ ($\Delta$$\varphi$/c\textsuperscript{2}) = 2.6946 $\times$ 10\textsuperscript{4} Hz/m $\times$ L  (terrestrial).\hspace{2em}(56)\par\endgroup

\noindent The factor 2 is the exact analogue of the Cooper-pair charge 2e. Prediction [G1] ([T2+]): parameter-free.

\subsection*{K.3  Proposition J.2 --- Topological Clock Identity (conditional on P24)}

\begingroup\centering ftop = mZ6c\textsuperscript{2}/(6h) = mec\textsuperscript{2}/h = fe   [if mZ6 = 6me].\hspace{2em}(57)\par\endgroup

\noindent Arithmetic identity given P24; if P24 closes at T1, the vortex tick is the electron Compton clock. Cross-check: m$\mu$/6 = 17.6097 MeV differs from mX17 = 17.62 MeV by 0.06\% --- under the general-level reading m$\mu$ = Nc$\cdot$m(9), mX17 = m(9), this is an identity of the cascade level N = 9 (numerical observation; repository [2]).

\subsection*{K.4  Conjecture CEC and Proposition J.3}

\noindent \textbf{Conjecture CEC ([T4]; [T3] with pipeline): }the Z6 vacuum energy density equals the Compton-scale density of its lightest composite:

\begingroup\centering ucond = (mec\textsuperscript{2})\textsuperscript{4}/($\hslash$c)\textsuperscript{3} = 1.42 $\times$ 10\textsuperscript{24} J/m\textsuperscript{3}.\hspace{2em}(58)\par\endgroup

\noindent Evidence: inverting fplasma = 18 mHz gives 6.6 $\times$ 10\textsuperscript{-14} GeV\textsuperscript{4} vs me\textsuperscript{4} = 6.82 $\times$ 10\textsuperscript{-14} (3\%); no other framework scale reproduces an 18 mHz plasma frequency. A single ratio at 3\% is weak evidence in isolation; the conjecture is internally falsifiable by any derivation of $\rho$cond.

\noindent \textbf{Proposition J.3 ([T2] conditional on CEC): }the gravitational plasma frequency is parameter-free:

\begingroup\centering fplasma = (1/2$\pi$) $\surd$( 4$\pi$GN (mec\textsuperscript{2})\textsuperscript{4} / (($\hslash$c)\textsuperscript{3} c\textsuperscript{2}) ) = 18.3 mHz,\hspace{2em}(59)\par\endgroup

\noindent placing the [G2] cutoff in the LISA band; the interpretation as collective screening remains [T3] pending PC2.

\subsection*{K.5  Consistency Problem PC2 --- Static vs. Dynamic Screening}

\noindent [Open] A static reading of the penetration depth $\lambda$grav = c/(2$\pi$fplasma) $\approx$ 2.6 $\times$ 10\textsuperscript{9} m is falsified by planetary orbits; only dynamical perturbations below fplasma may experience collective response, and this distinction is not yet established in the PH. Until PC2 resolves, [G2] and [G3] are not rigorously falsifiable.

\subsection*{K.6  Predictions [G1]--[G3]}

\noindent \textbf{[G1] ([T2+]): }fJ = (2.6946 $\times$ 10\textsuperscript{4} Hz/m) $\times$ L (Lemma J.1, eq. 56); parameter-free given Premise P1.

\noindent \textbf{[G2] ([T2] cond. CEC; interpretation [T3]): }suppression of stochastic GW power below 18.3 mHz (LISA band); PTA detections of the nHz background require a quantitative amplitude prediction before confrontation.

\noindent \textbf{[G3] ([T3]): }resonant $\delta$GN/GN $\sim$ 10\textsuperscript{-10} at fdrive $\in$ [0.38, 0.43] GHz, via the pipeline effective Lagrangian $\to$ per-cell coupling (Hypothesis A2, [T4]: the driving field couples photon-per-coherence-cell, depositing Edrive per composite pair; under A2 the identification Edrive = h$\cdot$fdrive is legitimate --- without it, Edrive is an energy density and the frequency conversion does not follow) $\to$ $\chi$($\omega$) $\to$ $\delta$GN($\omega$).

\subsection*{K.7  Open Problem Registry}

\noindent Full registry in repository [2], \path{docs/PROGRESS_AND_ROADMAP.md}.

\begin{footnotesize}\begin{longtable}{|>{\raggedright\arraybackslash}p{0.1300\textwidth}|>{\raggedright\arraybackslash}p{0.4200\textwidth}|>{\raggedright\arraybackslash}p{0.1100\textwidth}|>{\raggedright\arraybackslash}p{0.2400\textwidth}|}
\hline
\textbf{ID} & \textbf{Problem} & \textbf{Status} & \textbf{Blocks} \\ \hline
P23 & $\kappa$ = G (Lemma I.1$'$) --- §2.3 emergent-gauge construction & [T2+] & $\varepsilon$$\psi$ $\to$ T1, Neff $\to$ T1 ($\eta$B link withdrawn) \\ \hline
P14-B & Conserved baryon (rung) label; sign and size of the bifurcation bias (§13.4, §13.7) & open & $\eta$B, nucleon stability, Theorem KS \\ \hline
P14-C & Baryon binding scale ($(\hslash c)/a\approx163$ MeV) and internal structure consistent with $J\!\cdot\!\Sigma$ mixing; hexagonal radii are phenomenological (App.~L.4.7) & open & hadron radii, $m_n-m_p$ \\ \hline
P26 & $\nu$-gap quantitative closure (ceff = 11/6 vs 2G) & [T2] & P23 (locking mechanism) \\ \hline
P25 & Z6 composite gauge vertex & [T2] & P23, [G3], PC1-sub \\ \hline
P-4$\pi$/3 & Prefactor 4$\pi$/3 in $\Omega$DM/$\Omega$b direct formula & [T2+] & $\Omega$DM $\to$ T1 \\ \hline
P1 & be = 325/1152 from 2PI kernel & [T2] & mass ratios Class B \\ \hline
P2/P12 & $\Phi$($\Lambda$) = $\Lambda$ fixed point; closure $\to$ adimensional ratios internal; 1 dimensional datum (Remark 5.1) & [T2] & GN, me $\to$ predictions \\ \hline
P8b & Passo A.3 integration ($\Sigma$2 = be $\times$ 15/8) & [T2] & lepton mass ratios precision \\ \hline
P24 & mZ6 = 6me vs 6$\surd$(3/$\pi$)me (f$\pi$ branch, §J.11) & [T2] & Prop. J.2, CEC \\ \hline
C-$\Delta$m21 & Derive $\Delta$m\textsuperscript{2}\textsubscript{21} = (18/25) m(15)\textsuperscript{2} from the $\nu$-channel (P26-C) & [T4] & second datum $\to$ prediction \\ \hline
PC1 & Gravity: emergent vs. ontological (hAB vs. graviton) & [T2] & [G3] pipeline \\ \hline
PC2 & Static vs. dynamic screening of $\lambda$grav & open & [G2], [G3] falsifiability \\ \hline
P-PMNS & Sector-bridging PMNS texture; CC bridge = proto-W (I = 3); texture open; flavour basis fixed, eq. (1) alone gives sin\textsuperscript{2}$\theta_{23}$ $\approx$ 0.21 (§17.3) & [T3] & mixing sector; F3 \\ \hline
P-|O| & Overlap sector for absolute rates; |O|\textsuperscript{2}(CF) $\le$ 0.370 (2$\sigma$) & [T2] & rates $\to$ T1 \\ \hline
P-$\varepsilon$conf & Formation-overlap democracy ($\varepsilon$conf = 4/9, n/n$\gamma$) & [T2] & withdrawn with eq.~(33) (§13.3) \\ \hline
P-Pauli & Pauli structure (restated): auto-singlets from N $\ge$ 7; spatial shells from N $\ge$ 13, nuclei at N = 15+; per-level structure open; 39-count floor & [T2] & DM spin sector \\ \hline
C-Q-rat & Charge rationality / integer quantization as Z6 selection principle & [T3] & why Nc = 6 \\ \hline
P-Lorentz & Local Lorentz invariance: D1 + R1 resolved (scalar channel; conserved-charge theorem); open: loop-level E/$\Lambda$meta vs SME; induced-metric universality & [T2+] & emergent-metric T-status \\ \hline
P-walls & ZKO relic walls: resolved-conditional (IR tension + exact chiral-bias spectrum 5:8:9); open: $\kappa$CP$\cdot$V\textsubscript{0} normalization & [T2/T3] & cosmological safety \\ \hline
F-fork & RESOLVED $\to$ twelve-state bookkeeping via U-count (U3: (2, 2$\pi$/3)); residuals reviewed and held & [T2] & closed \\ \hline
C-web & Closed-form master web (§5.1) RESOLVED: exact rearrangements of registered structure; O(1/S) two-loop slot & [T2+] & master invariant; F17 \\ \hline
\end{longtable}\end{footnotesize}

\begingroup\centering\footnotesize\itshape Table 15: Open problems and cascading dependencies. P23 is the critical path; P2/P12 is the deepest.\par\endgroup

\subsection*{K.8  The Chiral-Flip Channel and the e\textsuperscript{+}e\textsuperscript{-} Sector}

\noindent The Z6 composite structure predicts the Chiral-Flip Charged Current, a zero-parameter structural channel ([T1]: selection rules, existence) with no analogue in three-neutrino QFT:

\begingroup\centering $\nu$$\mu$(0, 4, 2) + n(0, 1, 5) $\to$ e\textsuperscript{-}(0, 1, 2) + p(0, 4, 5),   RCF = NDirac g\textsuperscript{2}meta (me/m$\mu$)\textsuperscript{2} $\times$ |O|\textsuperscript{2} = 6.16 $\times$ 10\textsuperscript{-4} |O|\textsuperscript{2}.\hspace{2em}(60)\par\endgroup

\noindent Here |O|\textsuperscript{2} $\in$ [0, 1] is the BSE formation overlap (open sector P-|O|).

\noindent \textbf{\textit{Theorem K.8.1 (flatness, [T1]). }}\textit{Z6 vertex conservation forces any t-channel mediator to carry (I, Q) = (3, +1). Exhaustive enumeration of the 36 bilinears shows the only such states are the six proto-W members at MW; their propagator is flat to < 2 $\times$ 10\textsuperscript{-4} over accelerator kinematics, and hAB (I = 0, Q = 0) cannot mediate. RCF is therefore energy- and baseline-independent at every accessible order.}

\noindent \textbf{Experimental status. }A flat ratio cannot reproduce a spectrally localized excess for any |O|\textsuperscript{2}; MicroBooNE's exclusive 1e1p channel [10, 14] bounds |O|\textsuperscript{2} $\le$ 0.244 (90\% CL) and $\le$ 0.370 (2$\sigma$), i.e. RCF $\le$ 2.28 $\times$ 10\textsuperscript{-4} --- below the MiniBooNE anomaly and decoupled from it. Confrontation protocol: a genuine CF-CC signal appears as a 1e1p excess identical (after flux scaling) at SBND (110 m) and ICARUS (600 m), incompatible with any L/E fit, at RCF $\le$ 2.28 $\times$ 10\textsuperscript{-4}; SBND's percent-level $\nu$e-like rate measurement tightens the |O|\textsuperscript{2} ceiling directly.

\noindent \textbf{\textit{Theorem K.8.2 (e\textsuperscript{+}e\textsuperscript{-} uniqueness, MB-U, [T1]). }}\textit{A composite can decay to e\textsuperscript{+}e\textsuperscript{-} only if its net psion number vanishes --- the mixed $\psi$$\bar{\psi}$ sector; every pure-psion composite is excluded by conservation. Enumeration of the neutral $\psi$$\bar{\psi}$ spectrum leaves a single massive candidate below 0.6 GeV: the proto-colour singlet of the X17 family at cascade level 9 (17.62 MeV).}

\noindent The next neutral $\psi$$\bar{\psi}$ states --- the level-8 family member (0.610 GeV) and hAB (0.613 GeV) --- lie above the short-baseline excess window and decay to two-shower $\gamma$$\gamma$ topologies, outside the single-ring e\textsuperscript{+}e\textsuperscript{-} class. Any dark-sector e\textsuperscript{+}e\textsuperscript{-} interpretation of short-baseline data is therefore confined, within the PH, to this one state. The corresponding conjecture (C-MB-X17) is registered in the repository under E4 quarantine with ex-ante kill criteria and is not carried in the prediction registry (Table 4).

\subsection*{K.9  Local Lorentz Invariance of the Emergent Metric}

\noindent Every condensate-based emergent gravity owes an account of why local Lorentz invariance is exact or suppressed below experimental bounds: a condensate generically selects a frame. The PH clears this at every currently computable level. (i) [T1] The contact interaction of eq. (1) carries no Dirac structure --- J, $\Sigma$ act on the internal space and commute with every $\gamma$\textsuperscript{u} --- so the available condensation channels $\langle$$\bar{\psi}$$\psi$$\rangle$, $\langle$$\bar{\psi}$J$\Sigma$$\psi$$\rangle$ are Lorentz scalars; a scalar condensate preserves Lorentz exactly, as the QCD chiral condensate does. (ii) [T2] The Fierz-induced vector and axial channels cannot condense: they couple to conserved (axial-)charge densities, which vanish in the $\mu$ = 0 vacuum, and the mean-field solution is uniquely trivial; the C-invariance of the interaction independently forbids the C-odd vector order parameter, and the tensor projection is subcritical (relative Fierz coefficient $-$1/8). (iii) Consequences: c\_gw = c\_$\gamma$ by structure --- the photon bilinear and the metric response propagate on one Lorentz-invariant substrate (GW170817 satisfied without tuning) --- and the equivalence principle as single-substrate universality (argument grade; the induced-metric computation is the open deliverable). Residual: loop-level operators in powers of E/$\Lambda$meta ($\approx$ 8 $\times$ 10\textsuperscript{-10} at 1 TeV) against the SME tables; if the suppression scale is $\Lambda$meta rather than MPl, dimension-six coefficients sit near future sensitivity --- a testable window.

\subsection*{K.10  The Fate of the Z6 Domain Walls}

\noindent Spontaneous breaking of the discrete Z6 at the first-order transition T = $\Lambda$meta Kibble-produces a relic wall network --- the Zel'dovich--Kobzarev--Okun problem, fatal for a naive tension $\Lambda$meta\textsuperscript{3} (fifteen orders above the CMB bound $\sigma$\textsuperscript{1/3} $\lesssim$ 1 MeV). The resolution has two layers.

\noindent \textbf{Layer 1 (infrared walls, [T2]). }The walls interpolate the condensate phase between adjacent vacua --- pseudo-Goldstone territory, stiffness f$\pi$\textsuperscript{2}, gap mZ6, both MeV scale (f$\pi$ = $\surd$(3/$\pi$)$\cdot$me = 0.4994 MeV, ratified: the §3 normalization is derived by the 2$\pi$ condition, f$\pi$\textsuperscript{2} = Nc$\cdot$me\textsuperscript{2}/(2$\pi$); mZ6 = 6f$\pi$ = 2.996 MeV, with the P18 route differing by 1.1\% --- audit Q2). The registered virial tension $\sigma$ = f$\pi$\textsuperscript{2}mZ6/$\pi$\textsuperscript{2} = 0.0757 MeV\textsuperscript{3} gives $\sigma$\textsuperscript{1/3} = 0.423 MeV: forty-six orders below the naive estimate and below the Zel'dovich bound. Even a fully persistent scaling network is invisible at current sensitivity: $\rho$wall/$\rho$crit $\approx$ 3 $\times$ 10\textsuperscript{-6} today, induced $\delta$T/T $\approx$ 3.5 $\times$ 10\textsuperscript{-7}, BBN-era impact $\approx$ 3 $\times$ 10\textsuperscript{-22}. The normalization is convention-sensitive (the sine-Gordon form gives 1.82 MeV, straddling the bound), so robustness rests on Layer 2; the residual mZ6 provenance audit (Q2) is thereby promoted to a cosmological-safety input.

\noindent \textbf{Layer 2 (chiral bias, [T2] shape / [T3] normalization). }The six vacua are not protected-degenerate: mJ(k) = k $-$ 5/2 is injective, so no nontrivial Z6 shift preserves the chiral structure. The splitting spectrum inherits the exact chiral autocorrelation $\Delta$C(s) = C(0) $-$ $\Sigma$k mJ(k)mJ(k+s) = \{0, 15, 24, 27, 24, 15\}: three wall classes in ratio 5 : 8 : 9, all positively biased, with the s = 3 walls (the pure L$\leftrightarrow$R exchange) maximally split and dying first. The normalization follows by spurion analysis: any vacuum-splitting operator crosses chirality twice, forcing $\kappa$CP = (me/$\Lambda$meta)\textsuperscript{2} --- the same CP-odd scale that powers baryogenesis. The estimated annihilation time, t\_ann \textasciitilde{} $\sigma$/$\varepsilon$ $\approx$ 4 yr, lands after BBN (where the walls are harmless) and well before recombination: the asymmetry that makes the matter kills the walls. Falsifiable branch: should the bias fail, the persistent network predicts $\delta$T/T \textasciitilde{} 4 $\times$ 10\textsuperscript{-7} ISW-class signatures, within reach of near-future CMB experiments.

\clearpage
\section*{L  Method Audit Record: Eliminated Routes and Negative Findings}

\noindent This appendix records the adversarial audit trail of the derivation programme: eliminated routes, self-caught circularity audits, and negative findings. A route killed by its own audit is a result, and is registered with the same care as a theorem; the complete session-level record, including all intermediate numerics, is preserved in the repository [2].

\subsection*{L.1  The $\varepsilon$$\psi$ Campaign: Audits and Kill Tables}

\noindent The $\varepsilon$$\psi$ = 25/169 closure campaign produced, beyond the final results carried in Appendix J, an audit trail whose density is itself information: every naive shortcut to T1 was tried and killed. Eight self-caught circularity audits are on record. C-I8.1: the `one hop, one quantum' route killed --- a single hop transports 2j = 5 flux quanta, not one. C-I8.2: flux-energetics linearity gap --- Schemes A (linear) and B (min-rule) registered as anchor-degenerate. C-I8.3: the holonomy definition g\textsuperscript{2} $\equiv$ 2$\pi$/N presupposes self-duality --- circular, killed before registration; the same audit corrected an overstatement of V.6 (equivalence, not support). C-I8.4: Scheme B's spectral identification was reverse-solved from the target --- killed as derivation, retained as exact equivalence. C-I8.5: architecture finding --- the §2 Lagrangian carries Z6 as a GLOBAL symmetry (domain walls, not gauge strings); the gauged-ladder line was reclassified as framework-ungrounded model study. C-I8.6: the framework J$\Sigma$ operator is not the Pauli J$\cdot$$\sigma$; the locked diagonal values are not eigenvalues of the full operator. C-I8.7: the combination 2f3 + f2 = j is target-forced --- registered as equivalence, killed as closure. C-I8.8: the self-duality route reduces to J = t, a restatement of Lemma I.1$'$.

\noindent Kill tables (exact). JR-anchor competitor table: $\langle$Q$\Sigma$$\rangle$/(2j)\textsuperscript{2} and $\langle$Q$\Sigma$$\rangle$/(2j) killed at j = 1/2 ($\varepsilon$$\psi$ = 16/25 $\ne$ 1); the linear form Nc$\langle$Q$\Sigma$$\rangle$/(2j) killed at j = 5/2 ($\varepsilon$$\psi$ > 1); the survivor is Nc$\langle$Q$\Sigma$$\rangle$/(2j)\textsuperscript{2} $\equiv$ Ind/(2j). Charge-channel survey for Eflip (flip (0,1,2) $\to$ (0,1,3)): pair channel $\Sigma$QaQb requires g = $-$27/8 (not derived, killed); self channel $\Sigma$Qa\textsuperscript{2} requires g = 27/10 (reverse-solved, killed as derivation); total-charge and charge--chirality channels require |g| = 9/2, opposite signs (not derived, killed). Verdict: the charge sector is forced as the splitter (Theorem L.1) but not computable from the derived structures.

\noindent Theorem L.1 (flip degeneracy, [T1]). The kink bond exchanges k = 2 $\leftrightarrow$ 3 (m = $-$1/2 $\to$ +1/2 on one constituent); since |m| and m\textsuperscript{2} are flip-invariant, both $\Sigma$|mi| and C2 = $\Sigma$mi\textsuperscript{2} are exactly degenerate between any Q = $-$1 composite and its flip partner --- e.g. (0,1,2) $\leftrightarrow$ (0,1,3), both at $\Sigma$|m| = 9/2, C2 = 35/4. The unique derived quantity split by the flip is the Z6 charge, $\Delta$Q = 5/3: the splitter is forced to be the charge sector.

\noindent Exact J$\Sigma$ spectrum (exact result). The framework operator Jz$\Sigma$z + $\lambda$(J+$\Sigma$$-$ + J$-$$\Sigma$+) with unit-element $\Sigma$$\pm$ has exact spectrum \{5/2 ($\times$2), $-$1/2 $\pm$ $\surd$3 ($\times$2 each), $-$1/2 $\pm$ $\surd$21/2 ($\times$2 each), 1, $-$2\}; the locked diagonal values 1/2 and 3/2 are NOT eigenvalues of the full operator --- the deep-gap projection is essential, and its physical origin is exactly the P26 locking mechanism. The naive continuum embedding Hint = M$\cdot$(J$\cdot$$\Sigma$) was killed by direct 24-component Dirac diagnostics in the Z6 kink (bulk thresholds at M$\cdot$\{1, 2, 2.5, 3.5\}, no zero mode).

\noindent S7 prototype and the three-body negative finding. The gauged k-ladder electrostatics scan (twelve configurations: string orientation $\times$ Z6 background) produced exactly ONE pass: Eflip = (1/2)(2$\cdot$d(4)\textsuperscript{2} $-$ 3$\cdot$d(1)\textsuperscript{2}) = 5/2 = j, with the Z6 conjugacy min-rule realized dynamically. Sensitivity ledger: $\kappa$ = 1, background c = 0, and string geometry were NOT derived --- three undetermined choices. The three-body extension ($\mu$ = (0,1,2) $\to$ (0,1,3), every simple prescription) FAILED to reproduce 5/2 (results ranged $-$23/2 to 19/2): the prototype hit is possibly a single-particle truncation artifact, and numerics were frozen pending the kink-as-gauge-object derivation. This stop condition --- no prescription-shopping --- is the campaign's discipline in action.

\noindent Operational definition of $\varepsilon$$\psi$(G) and correction constants. $\varepsilon$$\psi$(G) = sin\textsuperscript{2}(2$\theta$$\chi$(G)) with tan $\theta$$\chi$(G) = Nc$\langle$$\psi$0(G)|Q$\Sigma$z|$\psi$0(G)$\rangle$/(2j)\textsuperscript{2}, evaluated on the dynamical zero mode with unlocked penalty 2G. First-order symbolic PT reproduces the leading term $-$480/2197 fraction-for-fraction. Registered constants: $\kappa$6 = 0.228(1), with a reconstruction flag --- the re-derived value is 0.2208, a 3\% discrepancy attributed to subleading extraction, candidate corrected value $\kappa$6 $\approx$ 0.221; $\kappa$8 = 0.500(1). These numerics establish dynamical robustness of the formula's ingredients; they do not, and cannot, prove the Lemma.

\noindent Route eliminations. Ward identity: no canonical normalization produces $\kappa$ = 1 ($\chi$QQ $\to$ 0 in the deep gap; charge incompressible) --- $\kappa$ is a construction datum, not a symmetry consequence. Phase-stiffness discretization: $\kappa$/G \textasciitilde{} 10\textsuperscript{-5}, five orders short; the required enhancement carries the exact factor Nc\textsuperscript{6}/(4$\pi$\textsuperscript{2}) --- the first structural lead pointing to the vortex/winding mechanism, which then delivered $\kappa$/G = Nc\textsuperscript{2}/(4$\pi$\textsuperscript{2}) = sinc\textsuperscript{2}($\pi$/Nc) (§J.10). Kink energy (virial, Ekink = f$\pi$\textsuperscript{2}mZ6/$\pi$\textsuperscript{2}, a T2 by-product of the campaign): $\kappa$/G $\in$ \{0.203, 2\} depending on identification --- no route gives 1. P26 Round 2: the eight natural multiplicity assignments yield ce $\in$ \{11/6, 9/4, 31/12, 43/12, 2, 13/4, 3, 4\}; the single assignment hitting 2 contradicts the registered w3 = 2 --- the R1/R2 fork of §J.9 in its original data-first form.

\noindent Failed-route lesson. The leading-order field-theoretic estimate $\delta$$\chi$ \textasciitilde{} f$\pi$\textsuperscript{2}/($\Lambda$meta\textsuperscript{2} Ne\textsuperscript{2}) \textasciitilde{} 10\textsuperscript{-32} fails because it computes an energy ratio between microscopic scales, exponentially suppressed by S $\approx$ 213; the correct quantity is a geometric ratio of representation, dimensionless and $\Lambda$meta-independent --- the spectral width 2j = 5 provides exactly this. Wilson coincidence (registered, not promoted): for Z6 and only Z6, the unit-coupling Wilson plaquette cost equals the derived bond, 1 $-$ cos(2$\pi$/6) = 1/2 = t; Z6 is the largest ZN with rational Wilson cost (crystallographic N), consistent with Lemma V.3 singling out the physical group.

\subsection*{L.2  XL-1 Execution Record (the f$\pi$ fork, four routes)}

\noindent Four-route comparison (data-first): §3.2 explicit form $\surd$(3/$\pi$)$\cdot$me = 0.4994 MeV; XL-1/P24 exact form f$\pi$ = me = 0.5110 MeV; P18 table mZ6/6 = 0.5046 MeV (provenance unverified); CEC inversion 0.5069 MeV (provenance unverified). The two numerical routes sit between the two exact branches, \textasciitilde{}1\% from each --- inconclusive at current precision. CEC granularity: the inversion evidence equals 0.968$\cdot$me\textsuperscript{4} vs branch predictions 1.000$\cdot$me\textsuperscript{4} (A) and 0.912$\cdot$me\textsuperscript{4} (B) --- a slight lean toward branch A; the P18 route sits dead-centre. Decisive audits registered for the owner: Q1 --- provenance of the §3.2 normalisation (decides the branch); Q2 --- provenance of mZ6 in the P18 table; Q3 --- provenance of fplasma = 18 mHz and $\rho$cond. The branch bookkeeping and the $\eta$B quarantine are carried in §J.11.

\subsection*{L.3  Registered Observations}

\noindent 12$\alpha$ expansion. 1 $-$ sinc\textsuperscript{2}($\pi$/Nc) $\approx$ 0.08811 vs 12$\alpha$ $\approx$ 0.08757 (0.6\%); if exact, $\alpha$ = (4$\pi$\textsuperscript{2} $-$ 36)/(48$\pi$\textsuperscript{2}) $\approx$ 1/136.19 --- better than Eddington's 1/136, not a derivation (Wyler-class caution registered); using $\kappa$/G = 1 $-$ 12$\alpha$ in the pass criterion worsens $\eta$B from percent-level to 8.8\% --- the coincidence is an observation, not a route.

\noindent Chord/arc/sagitta geometry. $\surd$($\kappa$/G) = chord/arc of a Z6 sector (half-angle 30$^{\circ}$); the circle's own sin\textsuperscript{2} + cos\textsuperscript{2} = 1 supplies the Pythagorean relation (semi-chord, R $-$ sagitta, R); $\kappa$ lives in the flat lattice metric, G in the curved phase-circle metric. Nyquist reading: sinc is the universal sampling form factor --- the Z6 condensate literally samples the massless psion's phase at Nc = 6 points, and the P23 deficit is its sampling error; the geometry explains the gap completely and closes it only as Nc $\to$ $\infty$, which Nc = 6 [T1] forbids.

\subsection*{L.4  Charge-Conjugation Audit and the Lepton-Generation Mechanism}

\noindent \textbf{L.4.1 Composite identification (proton, neutron, exotic).} Proton is LRR $(0,4,5)$, $Q=1/9+4/9+4/9=+1$; neutron is LLR $(0,1,5)$, $Q=1/9-5/9+4/9=0$; RRR $(3,4,5)$ is a separate $Q=+2$ exotic BSM state. Both proton and neutron charges hold exactly under Theorem V.1 alone, no auxiliary scheme required. (The companion repository's \path{CHARGE_CONVENTIONS.md} is kept consistent with this identification; see its \S3a.)

\noindent \textbf{L.4.2 State-wise charge conjugation and Theorem V.1 (T1, label-independent).} Conjugation is the index map $k\to(6-k)\bmod 6$, fixing $k=0,3$. Requiring $Q_k=-Q_{\mathrm{conj}(k)}$ state-wise (not merely $\Sigma_k Q_k=0$) forces $Q_0=Q_3=0$ and $Q_1=-Q_5$, $Q_2=-Q_4$. Theorem V.1's values ($Q_0=1/9$, $Q_1=-5/9$, $Q_2=-5/9$, $Q_3=10/9$, $Q_4=4/9$, $Q_5=4/9$) satisfy none of these four constraints; the $(3,4,5)$ exotic state makes this concrete: V.1 gives $Q=+2$, while any conjugation-consistent assignment would give $Q=+1$ for the same triple --- a full unit apart. This is a property of the six $Q_k$ values themselves, independent of which composite is called proton or neutron.

\noindent \textbf{L.4.3 Adopted reading of C/CP.} Three repairs were evaluated: (i) retune the six $Q_k$ to satisfy state-wise conjugation exactly --- compatible with proton $(+1)$/neutron $(0)$ above, but changes the $(3,4,5)$ exotic charge from $+2$ to $+1$, a value not adopted in this revision; (ii) extend the conjugation map to include proto-colour --- any involution on 3 proto-colours has an odd fixed-point count, so this cannot remove both fixed points, and the one full involution tested trades the lepton-generation charge count for an unreachable neutron neutrality; (iii) retain conjugation only as an aggregate identity over the composite spectrum ($\Sigma_k Q_k$ tracked as an ensemble quantity), not as an exact per-state operator. Option (iii) is adopted: it requires no change to any $Q_k$ value, and \S13.1(2)'s $\varepsilon\psi$-based C/CP mechanism (footnote, \S13.1) is unaffected by it, since that mechanism never invokes state-wise $Q_k$-conjugation.

\noindent \textbf{L.4.4 Dual-sector labelling and the $\mu^-/\tau^-$ mechanism.} \S6 and Appendix~H \S H.3 both label LLR $(0,1,5)$, $(0,2,4)$, and $(1,2,3)$ as dual-sector (neutrino-like or baryon-like: n$/\nu_\tau$, $\nu_\mu/\Sigma^0$-like, $\nu_e/\Lambda^0$-like), with generation ordering for the third explicitly open. Theorem 6.2 ([T1]: exactly one $N=3$ configuration carries $Q=-1$) fixes the scope of this ambiguity precisely: it is neutrino-vs-baryon-like only, never charged-lepton-vs-baryon-like, since no $N=3$ triple other than $(0,1,2)$ can carry $Q=-1$.

\noindent The combinatorial identity of $\mu^-,\tau^-$ is therefore not an $N=3$ tuple at all, and is instead addressed by the Leptonic Transfer Matrix (P8b, Appendix A.2.5): $t_{12}=-8$ (exact, [T1]); $\langle 3|H_{dd}|1\rangle = t_{12}^2=8$ (three-body symmetry cancellation, exact, [T1]); diagonal set $C_2=\{10.75, 8.75, 6.75\}$ (exhaustive $Q=-1$-sector enumeration, [T1]) --- eigenvalues of a mixing operator on the lepton sector, mapped to $m_e:m_\mu:m_\tau$, rather than three separate psion composites. $N_{\mathrm{gen}}=3$ ([T1], \S9, kinematic bound $J_R=7/2>j_{\max}=5/2$) is consistent with this construction and bounds it from above; $m_\mu=105.66$~MeV additionally appears once (\S K.8, RCF formula) as an external CODATA reference, not a P8b-derived value.

\noindent \textbf{L.4.6 Dual sector and nucleon stability (v86).} If the same triple is a baryon or a neutrino according to its cascade level, and the level carries no internal quantum number, then $n\to\nu_\tau\gamma$ and its analogues are allowed by every symmetry of eq.~(1); stability requires an amplitude suppression $\lesssim5\times10^{-30}$. The tension recorded in L.4.4 is therefore a stability problem, resolved only by a conserved baryon (rung) label (§13.4, P14-B).

\noindent \textbf{L.4.7 Hadron size and geometry (v86, P14-C, [T4]).} An exploratory, pre-registered series (companion repository, WP-GEO to WP-GEO12 and WP-LJS) placed the six labels on geometric arrangements and computed charge-weighted mean-square radii with the charges $Q_k$ and equal constituent masses. The planar Z6 ring in natural label order (the hexagon, and equivalently any arrangement that keeps the cyclic order on one circle) is the only arrangement tested that reproduces, with one scale fixed by the proton, $R_n=\langle r^2\rangle_n/\langle r^2\rangle_p=-0.154$ (measured $-0.164$) and $R_{\Sigma^-}=0.846$ ($0.86\pm0.21$), and it gives $\langle r^2\rangle_{\Sigma^0}=0$ exactly; the scale is $a=1.212$ fm. Of the 720 orderings of the labels on the ring, only the natural cyclic order passes. The same arrangement fails for the triples containing $k=3$ ($\Lambda$, $\Sigma^+$, against chiral-perturbation-theory and lattice predictions), fails for magnetic moments, and cannot make leptons pointlike: leptonic sizes must come from the contact scale $1/\Lambda_{\rm meta}$, consistent with the exact $J\!\cdot\!\Sigma$-inertness of $(0,1,2)$. Three results prevent promotion. (a) $J\!\cdot\!\Sigma$ conserves $\sum k$ and therefore mixes $n$, $\Sigma^0$, $\Lambda$ (and $p$ with the two $\Sigma^+$ candidates); at the coupling of this paper the mixing destroys $R_n$ with adiabatic labels, and with ground-state labels gives $R_n\approx-0.08$ and $m_n<m_p$. (b) The scale $(\hslash c)/a=163$ MeV is not derived: in four dimensions the $J\!\cdot\!\Sigma$ contact term would need tuning to one part in $10^{24}$, and none of four pre-registered transmutation forms built from the constants of this paper lands within a factor 2 (nearest: rung $\ell=8$, factor 3.7; $e^{-S/g^2}$, factor 14). (c) The $k=3$ position that fits $\Lambda$ and $\Sigma^+$ (a $30^\circ$ shift on a dodecagon) is a discrete fit with no dynamical origin. This is recorded as open problem \textbf{P14-C}: a derived baryon binding scale and an internal structure consistent with the $J\!\cdot\!\Sigma$ dynamics. Nothing in the main text depends on it.

\noindent \textbf{L.4.5 Precision, Table 5 vs.\ \S A.3.1.} Table 5 (\S12) reports the leading-order transfer-matrix ratios, $m_\mu/m_e=218$ (5.4\% vs.\ 206.8) and $m_\tau/m_e=3795$ (9.1\% vs.\ 3477), explicitly captioned as targeting further refinement via P8b. \S A.3.1 reports the ratios after the P8b correction ($\Sigma_2=b_e\times 15/8$) at $+0.46\%$ and $+0.56\%$ respectively. The two are sequential stages of the same calculation --- leading order, then P8b-corrected --- not independent estimates; the P8b formal-integration step itself remains at [T2] (open-problems registry, App.~A), so the $0.46\%$/$0.56\%$ figures carry that status pending closure of Passo A.3, while the $5.4\%$/$9.1\%$ leading-order figures are unconditional.

\clearpage
\subsection*{L.5  v85 Revision Record (October 2026)}

\noindent This revision is an audit, not a new physical result. Every change below is traceable to a commit in the companion repository [2]; negative findings are recorded with the same weight as positive ones.

\noindent \textbf{L.5.1 Dark-matter census proof.} The module \path{proofs/P_dark_matter.py} tested stability by single-psion removal, a test that can never fire because no psion is neutral, and failed its own assertion. It now uses the fragmentation criterion (no proper sub-multiset is a neutral singlet) with two independent algorithms; Theorem DM.1 is reproduced unchanged (39 states, N $\in$ \{9, 12, 15, 18, 21, 27\}), and N = 6 is confirmed fully unstable.

\noindent \textbf{L.5.2 Dark-matter mass and cascade level.} Three incompatible dark-matter masses coexisted (Table 3, m(N) $\le$ 17.6 MeV; §14.2, 7.2 $\times$ 10\textsuperscript{10} GeV; a repository body-number law, 14.3 GeV). Because h\textsubscript{AB} and X17 share N = 2 but sit at cascade levels 8 and 9, the cascade level is not a function of body number; the Landau-pole mass of §14.2 is canonical, the 14.3 GeV result is retracted, and Table 3 is relabelled by cascade level (§8.1).

\noindent \textbf{L.5.3 Theorem KS.} Inequality (11) holds as arithmetic on the mass function (verified on all 1,295,920 partitions of every level up to 50), but physical fragmentation partitions body number, not level. Its physical application and Corollary 8.2 are suspended until a rule fixes the levels of fragments (§8.2). Dark-matter stability rests on DM.1 and psion-number conservation instead. Three k = 3 entries of Table 3 were wrong in v84 ($\ell$ = 2: no three-part partition exists; $\ell$ = 8 and 10 overstated); they are corrected.

\noindent \textbf{L.5.4 F5 (direct detection).} The bound $\sigma$ < 10\textsuperscript{-50} cm\textsuperscript{2} had no derivation and is retired. F5a (no dark-matter state outside the 39 singlets at m\textsubscript{DM}; [T2], E2) carries the falsification power; F5b (rate $\sim$ 10\textsuperscript{-54} cm\textsuperscript{2}, conditional on zero proto-colour coupling; [T4], E1) records the estimate. A kill criterion was registered ex ante (§9.1); the LZ event of September 2026 is below it.

\noindent \textbf{L.5.5 F3 and the PMNS programme.} F3 is a function of the PMNS matrix alone, which the PH does not derive, and the 7\% figure has no derivation in this release; F3 is carried at [T4], conditional on P-PMNS (§17.3). The companion repository opens a PMNS programme with pre-registered models: the flavour basis is fixed exactly (generation = proto-colour of the right-handed constituent; colour-diagonal charged current), and with charged leptons diagonal both terms of eq. (1) and the locking gap give sin\textsuperscript{2}$\theta_{12}$ near the measured value but sin\textsuperscript{2}$\theta_{23}$ $\approx$ 0.21. The charged-lepton side is blocked: the App. A.2.5 entries as printed do not reproduce the P8b mass ratios claimed in §A.3.1 (best m\textsubscript{$\mu$}/m\textsubscript{e} = 244 when m\textsubscript{$\tau$}/m\textsubscript{e} is matched), so the full P8b matrix is required.

\noindent \textbf{L.5.6 Unchanged.} No central numerical prediction of Classes A or B changes in this revision.

\subsection*{L.6  v86 Revision Record: the P14 Audit (October 2026)}

\noindent All work packages were pre-registered, with SHA-256 and a UTC timestamp, before execution. Failures and corrections are recorded with the same weight as results; the full log is in \path{p14d_psion_number_audit/README.md} of repository [2].

\noindent \textbf{L.6.1 Changes. }§13 rewritten (13.1--13.7); $\eta_B$ derivation withdrawn; F12 suspended; App.~J.1 and J.4 annotated; App.~A.4 antimatter statement corrected (antipsions); §8.2 status note: the missing level rule is a conserved baryon label (P14-B); §14.3 direct route marked unsupported; open-problem tables gain P14-B and P14-C; App.~L.4.7 added (hadron-size geometry, phenomenological); abstract, §4.1, §12 and §15--19 summaries updated accordingly. No other numerical prediction changes.

\noindent \textbf{L.6.2 Errors made during the audit and corrected. }(i) Sphalerons were first treated as external to the PH; they are psion composites and conserve $N_\psi$. (ii) Antimatter was first modelled as antipsions, then switched to index-inverted psion triples on a reading of A.5, and restored to antipsions by the reaction test of 13.1; the second step was the wrong one. (iii) In the domain test, unitarity was first imposed jointly across processes with different density scaling, which created asymmetry in equilibrium. A control caught this, the test was corrected, and the faulty version is retained. (iv) In the calibration test, the dark-matter block first patched a copy of the module globals and was redone in an addendum. (v) The anticipated count of neutral triples was wrong (it is five). (vi) In the bifurcation test, the registered endpoint bookkeeping had a sign error, which was corrected (two antineutrinos per baryon), and the anticipated efficiency was wrong by about three orders of magnitude. (vii) The bifurcation networks first used free psions below $\Lambda_{\rm meta}$, which App.~I forbids; the mechanism was reformulated as pair production of unconfined composites. (viii) Rung assignments in two tests contradicted the stability requirement for charged leptonic copies; corrected, they make the sign of $\eta_B$ law-dependent. (ix) A test of the vacuum phase counted the condensate-generated bond twice, as an independent real term; without the double counting the phase is removable by $U(1)_M$, and the claimed CP source does not exist.

\noindent \textbf{L.6.3 Registered control failures. }All are at or near the floating-point floor or due to integrator limits: assembly-network controls ($2\times10^{-13}$ against $10^{-15}$); temperature-programme step convergence (1--2\%); domain test C3 ($2.7\times10^{-6}$) and C4 ($1.3\times10^{-12}$); phase-in-$J\!\cdot\!\Sigma$ test C3 ($1.8\times10^{-6}$) and one sensitivity run not completed; bifurcation tests, $B+L=0$ at the numerical floor ($\le1.1\times10^{-8}$).

\section*{References}

\noindent [1] Planck Collaboration, Planck 2018 results. VI. Cosmological parameters, Astron. Astrophys. 641, A6 (2020). arXiv:1807.06209.

\noindent [2] K. M. Maia, psion-proofs: Verification repository for the Psion Hypothesis, GitHub, \url{https://github.com/Maia-KarlosMarden/psion-proofs} (2026).

\noindent [3] A. D. Sakharov, Vacuum quantum fluctuations in curved space, Sov. Phys. Dokl. 12, 1040 (1967).

\noindent [4] Y. Nambu and G. Jona-Lasinio, Dynamical model of elementary particles based on an analogy with superconductivity, Phys. Rev. 122, 345 (1961).

\noindent [5] S. Weinberg and E. Witten, Limits on massless particles, Phys. Lett. B 96, 59 (1980).

\noindent [6] I. M. Gel'fand and A. M. Yaglom, Integration in functional spaces and its applications in quantum physics, J. Math. Phys. 1, 48 (1960).

\noindent [7] R. Jackiw and C. Rebbi, Solitons with fermion number 1/2, Phys. Rev. D 13, 3398 (1976).

\noindent [8] A. J. Krasznahorkay et al., Observation of anomalous internal pair creation in \textsuperscript{8}Be, Phys. Rev. Lett. 116, 042501 (2016).

\noindent [9] MiniBooNE Collaboration, Significant excess of electronlike events in the MiniBooNE short-baseline neutrino experiment, Phys. Rev. Lett. 121, 221801 (2018).

\noindent [10] MicroBooNE Collaboration, Search for an excess of electron neutrino interactions in MicroBooNE using multiple final-state topologies, Phys. Rev. Lett. 128, 241801 (2022).

\noindent [11] E. Tiesinga, P. J. Mohr, D. B. Newell, and B. N. Taylor, CODATA recommended values of the fundamental physical constants: 2018, Rev. Mod. Phys. 93, 025010 (2021).

\noindent [12] K. M. Maia, Derivation of the RGD Kernel from the PH Lagrangian, v2 (2026). [companion paper]

\noindent [13] K. M. Maia, Dark Matter: Existence Without Necessity (2026). [companion paper]

\noindent [14] MicroBooNE Collaboration, Search for an anomalous excess of charged-current quasielastic $\nu$e interactions with the MicroBooNE experiment using Deep-Learning-based reconstruction, Phys. Rev. D 105, 112003 (2022).

\noindent [15] MicroBooNE Collaboration, Search for an anomalous excess of charged-current $\nu$e interactions without pions in the final state with the MicroBooNE experiment, Phys. Rev. D 105, 112004 (2022).


\noindent [16] A. D. Sakharov, Violation of CP invariance, C asymmetry, and baryon asymmetry of the universe, JETP Lett. 5, 24 (1967).

\noindent [17] Super-Kamiokande Collaboration, Search for nucleon decay via $n\to\bar\nu\pi^0$ and $p\to\bar\nu\pi^+$, arXiv:1305.4391.

\noindent [18] KamLAND Collaboration, Search for the invisible decay of neutrons with KamLAND, Phys. Rev. Lett. 96, 101802 (2006).

\noindent [19] A. G. Cohen, A. De R\'ujula, and S. L. Glashow, A matter--antimatter universe?, Astrophys. J. 495, 539 (1998).

\noindent [20] E. W. Kolb and S. Wolfram, Baryon number generation in the early universe, Nucl. Phys. B 172, 224 (1980).

\noindent [21] E. Witten, Current algebra, baryons, and quark confinement, Nucl. Phys. B 223, 433 (1983).

\end{document}
